
\newpage
{\fontsize{19}{22}\selectfont\bfseries\color{navy}Speed and emission-rate profiles}\par
{\color{muted}Detector 78 worked example. The three rows have different time meanings.}\par
\vspace{5pt}
\legendkey{methodO}{O: recorded}\quad\legendkey{methodS}{S: mean speed}\quad\legendkey{methodR}{R: reconstructed}\quad\legendkey{methodM}{M: two-speed}\quad\legendkey{methodA}{A: cubic}
\par\vspace{8pt}\begin{center}\resizebox{\linewidth}{!}{\begin{tikzpicture}[x=1pt,y=1pt]
\useasboundingbox (0,0) rectangle (509,540);
\node[anchor=west,font=\small\bfseries,text=navy] at (5,531) {(a) O and S: detector-clock speeds};
\draw[rulegray,line width=.25pt] (40,395.000) -- (239,395.000);
\node[anchor=east,font=\scriptsize] at (36,395.000) {0};
\draw[rulegray,line width=.25pt] (40,428.143) -- (239,428.143);
\node[anchor=east,font=\scriptsize] at (36,428.143) {20};
\draw[rulegray,line width=.25pt] (40,461.286) -- (239,461.286);
\node[anchor=east,font=\scriptsize] at (36,461.286) {40};
\draw[rulegray,line width=.25pt] (40,494.429) -- (239,494.429);
\node[anchor=east,font=\scriptsize] at (36,494.429) {60};
\draw[black,line width=.35pt] (50.474,395) -- (50.474,392);
\node[anchor=north,font=\scriptsize] at (50.474,391) {15:00};
\draw[black,line width=.35pt] (92.368,395) -- (92.368,392);
\node[anchor=north,font=\scriptsize] at (92.368,391) {16:00};
\draw[black,line width=.35pt] (134.263,395) -- (134.263,392);
\node[anchor=north,font=\scriptsize] at (134.263,391) {17:00};
\draw[black,line width=.35pt] (176.158,395) -- (176.158,392);
\node[anchor=north,font=\scriptsize] at (176.158,391) {18:00};
\draw[black,line width=.35pt] (218.053,395) -- (218.053,392);
\node[anchor=north,font=\scriptsize] at (218.053,391) {19:00};
\draw[black,line width=.45pt] (40,511) -- (40,395) -- (239,395);
\node[anchor=north,font=\scriptsize] at (139.5,376) {Clock time};
\node[rotate=90,font=\scriptsize] at (7,453.0) {Speed (mph)};
\begin{scope}
\clip (40,395) rectangle (239,511);
\draw[methodO,line width=1.05pt,] plot coordinates {(40.000,480.549) (43.491,480.549) (43.491,476.696) (46.982,476.696) (46.982,475.311) (50.474,475.311) (50.474,475.341) (53.965,475.341) (53.965,473.156) (57.456,473.156) (57.456,470.840) (60.947,470.840) (60.947,467.503) (64.439,467.503) (64.439,466.651) (67.930,466.651) (67.930,465.275) (71.421,465.275) (71.421,463.339) (74.912,463.339) (74.912,458.480) (78.404,458.480) (78.404,451.033) (81.895,451.033) (81.895,444.802) (85.386,444.802) (85.386,441.309) (88.877,441.309) (88.877,439.940) (92.368,439.940) (92.368,438.081) (95.860,438.081) (95.860,432.409) (99.351,432.409) (99.351,426.589) (102.842,426.589) (102.842,423.145) (106.333,423.145) (106.333,419.718) (109.825,419.718) (109.825,418.796) (113.316,418.796) (113.316,419.335) (116.807,419.335) (116.807,418.853) (120.298,418.853) (120.298,418.497) (123.789,418.497) (123.789,416.285) (127.281,416.285) (127.281,415.042) (130.772,415.042) (130.772,413.551) (134.263,413.551) (134.263,413.032) (137.754,413.032) (137.754,413.592) (141.246,413.592) (141.246,413.852) (144.737,413.852) (144.737,412.029) (148.228,412.029) (148.228,411.097) (151.719,411.097) (151.719,411.043) (155.211,411.043) (155.211,412.411) (158.702,412.411) (158.702,413.457) (162.193,413.457) (162.193,415.813) (165.684,415.813) (165.684,418.021) (169.175,418.021) (169.175,420.083) (172.667,420.083) (172.667,422.623) (176.158,422.623) (176.158,426.530) (179.649,426.530) (179.649,431.060) (183.140,431.060) (183.140,433.483) (186.632,433.483) (186.632,438.351) (190.123,438.351) (190.123,442.888) (193.614,442.888) (193.614,445.504) (197.105,445.504) (197.105,452.446) (200.596,452.446) (200.596,457.889) (204.088,457.889) (204.088,462.767) (207.579,462.767) (207.579,464.726) (211.070,464.726) (211.070,470.543) (214.561,470.543) (214.561,474.194) (218.053,474.194) (218.053,478.999) (221.544,478.999) (221.544,482.789) (225.035,482.789) (225.035,490.297) (228.526,490.297) (228.526,494.485) (232.018,494.485) (232.018,498.729) (235.509,498.729) (235.509,498.747) (239.000,498.747)};
\draw[methodS,line width=1.05pt,dashed] plot coordinates {(40.000,434.038) (239.000,434.038)};
\end{scope}
\node[anchor=west,font=\small\bfseries,text=navy] at (264,531) {(b) O and S: rates at these speeds};
\draw[rulegray,line width=.25pt] (299,395.000) -- (498,395.000);
\node[anchor=east,font=\scriptsize] at (295,395.000) {0};
\draw[rulegray,line width=.25pt] (299,421.364) -- (498,421.364);
\node[anchor=east,font=\scriptsize] at (295,421.364) {5};
\draw[rulegray,line width=.25pt] (299,447.727) -- (498,447.727);
\node[anchor=east,font=\scriptsize] at (295,447.727) {10};
\draw[rulegray,line width=.25pt] (299,474.091) -- (498,474.091);
\node[anchor=east,font=\scriptsize] at (295,474.091) {15};
\draw[rulegray,line width=.25pt] (299,500.455) -- (498,500.455);
\node[anchor=east,font=\scriptsize] at (295,500.455) {20};
\draw[black,line width=.35pt] (309.474,395) -- (309.474,392);
\node[anchor=north,font=\scriptsize] at (309.474,391) {15:00};
\draw[black,line width=.35pt] (351.368,395) -- (351.368,392);
\node[anchor=north,font=\scriptsize] at (351.368,391) {16:00};
\draw[black,line width=.35pt] (393.263,395) -- (393.263,392);
\node[anchor=north,font=\scriptsize] at (393.263,391) {17:00};
\draw[black,line width=.35pt] (435.158,395) -- (435.158,392);
\node[anchor=north,font=\scriptsize] at (435.158,391) {18:00};
\draw[black,line width=.35pt] (477.053,395) -- (477.053,392);
\node[anchor=north,font=\scriptsize] at (477.053,391) {19:00};
\draw[black,line width=.45pt] (299,511) -- (299,395) -- (498,395);
\node[anchor=north,font=\scriptsize] at (398.5,376) {Clock time};
\node[rotate=90,font=\scriptsize] at (266,453.0) {CO$_2$ rate (kg/veh h)};
\begin{scope}
\clip (299,395) rectangle (498,511);
\draw[methodO,line width=1.05pt,] plot coordinates {(299.000,478.549) (302.491,478.549) (302.491,461.106) (305.982,461.106) (305.982,461.106) (309.474,461.106) (309.474,461.106) (312.965,461.106) (312.965,446.280) (316.456,446.280) (316.456,446.280) (319.947,446.280) (319.947,446.280) (323.439,446.280) (323.439,446.280) (326.930,446.280) (326.930,446.280) (330.421,446.280) (330.421,446.280) (333.912,446.280) (333.912,446.280) (337.404,446.280) (337.404,446.280) (340.895,446.280) (340.895,436.414) (344.386,436.414) (344.386,436.414) (347.877,436.414) (347.877,436.414) (351.368,436.414) (351.368,436.414) (354.860,436.414) (354.860,431.550) (358.351,431.550) (358.351,431.550) (361.842,431.550) (361.842,431.550) (365.333,431.550) (365.333,431.550) (368.825,431.550) (368.825,431.550) (372.316,431.550) (372.316,431.550) (375.807,431.550) (375.807,431.550) (379.298,431.550) (379.298,431.550) (382.789,431.550) (382.789,431.550) (386.281,431.550) (386.281,431.550) (389.772,431.550) (389.772,431.550) (393.263,431.550) (393.263,431.550) (396.754,431.550) (396.754,431.550) (400.246,431.550) (400.246,431.550) (403.737,431.550) (403.737,431.550) (407.228,431.550) (407.228,431.550) (410.719,431.550) (410.719,431.550) (414.211,431.550) (414.211,431.550) (417.702,431.550) (417.702,431.550) (421.193,431.550) (421.193,431.550) (424.684,431.550) (424.684,431.550) (428.175,431.550) (428.175,431.550) (431.667,431.550) (431.667,431.550) (435.158,431.550) (435.158,431.550) (438.649,431.550) (438.649,431.550) (442.140,431.550) (442.140,431.550) (445.632,431.550) (445.632,436.414) (449.123,436.414) (449.123,436.414) (452.614,436.414) (452.614,436.414) (456.105,436.414) (456.105,446.280) (459.596,446.280) (459.596,446.280) (463.088,446.280) (463.088,446.280) (466.579,446.280) (466.579,446.280) (470.070,446.280) (470.070,446.280) (473.561,446.280) (473.561,461.106) (477.053,461.106) (477.053,478.549) (480.544,478.549) (480.544,478.549) (484.035,478.549) (484.035,478.549) (487.526,478.549) (487.526,478.549) (491.018,478.549) (491.018,478.549) (494.509,478.549) (494.509,478.549) (498.000,478.549)};
\draw[methodS,line width=1.05pt,dashed] plot coordinates {(299.000,431.550) (498.000,431.550)};
\end{scope}
\node[anchor=west,font=\small\bfseries,text=navy] at (5,348) {(c) R: trip-average speed by entry};
\draw[rulegray,line width=.25pt] (40,212.000) -- (239,212.000);
\node[anchor=east,font=\scriptsize] at (36,212.000) {0};
\draw[rulegray,line width=.25pt] (40,245.143) -- (239,245.143);
\node[anchor=east,font=\scriptsize] at (36,245.143) {20};
\draw[rulegray,line width=.25pt] (40,278.286) -- (239,278.286);
\node[anchor=east,font=\scriptsize] at (36,278.286) {40};
\draw[rulegray,line width=.25pt] (40,311.429) -- (239,311.429);
\node[anchor=east,font=\scriptsize] at (36,311.429) {60};
\draw[black,line width=.35pt] (40.000,212) -- (40.000,209);
\node[anchor=north,font=\scriptsize] at (40.000,208) {0};
\draw[black,line width=.35pt] (84.725,212) -- (84.725,209);
\node[anchor=north,font=\scriptsize] at (84.725,208) {1};
\draw[black,line width=.35pt] (129.449,212) -- (129.449,209);
\node[anchor=north,font=\scriptsize] at (129.449,208) {2};
\draw[black,line width=.35pt] (174.174,212) -- (174.174,209);
\node[anchor=north,font=\scriptsize] at (174.174,208) {3};
\draw[black,line width=.35pt] (218.898,212) -- (218.898,209);
\node[anchor=north,font=\scriptsize] at (218.898,208) {4};
\draw[black,line width=.45pt] (40,328) -- (40,212) -- (239,212);
\node[anchor=north,font=\scriptsize] at (139.5,193) {Elapsed entry time (h)};
\node[rotate=90,font=\scriptsize] at (7,270.0) {Speed (mph)};
\begin{scope}
\clip (40,212) rectangle (239,328);
\draw[methodR,line width=1.05pt,] plot coordinates {(40.000,318.325) (40.199,318.325) (40.398,318.325) (40.597,318.324) (40.796,318.324) (40.995,318.324) (41.194,318.323) (41.393,318.323) (41.592,318.321) (41.791,318.320) (41.990,318.318) (42.189,318.316) (42.388,318.313) (42.587,318.309) (42.786,318.304) (42.985,318.299) (43.184,318.293) (43.383,318.286) (43.582,318.278) (43.781,318.269) (43.980,318.259) (44.179,318.247) (44.378,318.234) (44.577,318.220) (44.776,318.204) (44.975,318.186) (45.174,318.167) (45.373,318.146) (45.572,318.123) (45.771,318.099) (45.970,318.072) (46.169,318.044) (46.368,318.013) (46.567,317.980) (46.766,317.944) (46.965,317.907) (47.164,317.867) (47.363,317.824) (47.562,317.779) (47.761,317.731) (47.960,317.680) (48.159,317.626) (48.358,317.570) (48.557,317.510) (48.756,317.448) (48.955,317.382) (49.154,317.314) (49.353,317.242) (49.552,317.166) (49.751,317.088) (49.950,317.005) (50.149,316.920) (50.348,316.831) (50.547,316.738) (50.746,316.642) (50.945,316.542) (51.144,316.438) (51.343,316.330) (51.542,316.219) (51.741,316.103) (51.940,315.984) (52.139,315.861) (52.338,315.733) (52.537,315.602) (52.736,315.466) (52.935,315.327) (53.134,315.183) (53.333,315.035) (53.532,314.883) (53.731,314.727) (53.930,314.566) (54.129,314.401) (54.328,314.232) (54.527,314.059) (54.726,313.881) (54.925,313.699) (55.124,313.513) (55.323,313.323) (55.522,313.128) (55.721,312.929) (55.920,312.725) (56.119,312.517) (56.318,312.305) (56.517,312.089) (56.716,311.868) (56.915,311.644) (57.114,311.415) (57.313,311.181) (57.512,310.944) (57.711,310.702) (57.910,310.457) (58.109,310.207) (58.308,309.953) (58.507,309.695) (58.706,309.433) (58.905,309.167) (59.104,308.898) (59.303,308.624) (59.502,308.347) (59.701,308.065) (59.900,307.781) (60.099,307.492) (60.298,307.200) (60.497,306.904) (60.696,306.605) (60.895,306.302) (61.094,305.996) (61.293,305.686) (61.492,305.374) (61.691,305.058) (61.890,304.739) (62.089,304.416) (62.288,304.091) (62.487,303.763) (62.686,303.432) (62.885,303.098) (63.084,302.761) (63.283,302.422) (63.482,302.080) (63.681,301.736) (63.880,301.389) (64.079,301.040) (64.278,300.688) (64.477,300.335) (64.676,299.979) (64.875,299.621) (65.074,299.261) (65.273,298.899) (65.472,298.535) (65.671,298.170) (65.870,297.803) (66.069,297.434) (66.268,297.063) (66.467,296.692) (66.666,296.318) (66.865,295.944) (67.064,295.568) (67.263,295.191) (67.462,294.813) (67.661,294.434) (67.860,294.054) (68.059,293.673) (68.258,293.292) (68.457,292.909) (68.656,292.526) (68.855,292.143) (69.054,291.759) (69.253,291.374) (69.452,290.989) (69.651,290.604) (69.850,290.218) (70.049,289.833) (70.248,289.447) (70.447,289.061) (70.646,288.676) (70.845,288.290) (71.044,287.904) (71.243,287.519) (71.442,287.134) (71.641,286.749) (71.840,286.365) (72.039,285.981) (72.238,285.598) (72.437,285.215) (72.636,284.833) (72.835,284.451) (73.034,284.070) (73.233,283.690) (73.432,283.311) (73.631,282.932) (73.830,282.555) (74.029,282.178) (74.228,281.803) (74.427,281.428) (74.626,281.054) (74.825,280.682) (75.024,280.311) (75.223,279.941) (75.422,279.572) (75.621,279.205) (75.820,278.838) (76.019,278.473) (76.218,278.110) (76.417,277.748) (76.616,277.387) (76.815,277.028) (77.014,276.670) (77.213,276.314) (77.412,275.960) (77.611,275.607) (77.810,275.255) (78.009,274.905) (78.208,274.557) (78.407,274.211) (78.606,273.866) (78.805,273.523) (79.004,273.182) (79.203,272.842) (79.402,272.505) (79.601,272.169) (79.800,271.834) (79.999,271.502) (80.198,271.172) (80.397,270.843) (80.596,270.516) (80.795,270.192) (80.994,269.869) (81.193,269.548) (81.392,269.228) (81.591,268.911) (81.790,268.596) (81.989,268.283) (82.188,267.971) (82.387,267.662) (82.586,267.354) (82.785,267.049) (82.984,266.745) (83.183,266.444) (83.382,266.144) (83.581,265.847) (83.780,265.551) (83.979,265.257) (84.178,264.966) (84.377,264.676) (84.576,264.388) (84.775,264.103) (84.974,263.819) (85.173,263.538) (85.372,263.258) (85.571,262.980) (85.770,262.705) (85.969,262.431) (86.168,262.159) (86.367,261.889) (86.566,261.622) (86.765,261.356) (86.964,261.092) (87.163,260.830) (87.362,260.570) (87.561,260.312) (87.760,260.056) (87.959,259.802) (88.158,259.550) (88.357,259.300) (88.556,259.052) (88.755,258.806) (88.954,258.561) (89.153,258.319) (89.352,258.078) (89.551,257.839) (89.750,257.603) (89.949,257.368) (90.148,257.135) (90.347,256.903) (90.546,256.674) (90.745,256.447) (90.944,256.221) (91.143,255.997) (91.342,255.775) (91.541,255.555) (91.740,255.336) (91.939,255.120) (92.138,254.905) (92.337,254.692) (92.536,254.480) (92.735,254.271) (92.934,254.063) (93.133,253.857) (93.332,253.652) (93.531,253.450) (93.730,253.249) (93.929,253.050) (94.128,252.852) (94.327,252.656) (94.526,252.462) (94.725,252.269) (94.924,252.078) (95.123,251.889) (95.322,251.701) (95.521,251.515) (95.720,251.331) (95.919,251.148) (96.118,250.966) (96.317,250.787) (96.516,250.608) (96.715,250.432) (96.914,250.257) (97.113,250.083) (97.312,249.911) (97.511,249.740) (97.710,249.571) (97.909,249.404) (98.108,249.238) (98.307,249.073) (98.506,248.910) (98.705,248.748) (98.904,248.588) (99.103,248.429) (99.302,248.272) (99.501,248.115) (99.700,247.961) (99.899,247.808) (100.098,247.656) (100.297,247.505) (100.496,247.356) (100.695,247.208) (100.894,247.062) (101.093,246.916) (101.292,246.773) (101.491,246.630) (101.690,246.489) (101.889,246.349) (102.088,246.210) (102.287,246.073) (102.486,245.937) (102.685,245.802) (102.884,245.668) (103.083,245.536) (103.282,245.404) (103.481,245.275) (103.680,245.146) (103.879,245.018) (104.078,244.892) (104.277,244.766) (104.476,244.642) (104.675,244.520) (104.874,244.398) (105.073,244.277) (105.272,244.158) (105.471,244.039) (105.670,243.922) (105.869,243.806) (106.068,243.691) (106.267,243.577) (106.466,243.464) (106.665,243.353) (106.864,243.242) (107.063,243.132) (107.262,243.024) (107.461,242.916) (107.660,242.810) (107.859,242.704) (108.058,242.600) (108.257,242.496) (108.456,242.394) (108.655,242.292) (108.854,242.192) (109.053,242.092) (109.252,241.994) (109.451,241.897) (109.650,241.800) (109.849,241.704) (110.048,241.610) (110.247,241.516) (110.446,241.423) (110.645,241.331) (110.844,241.241) (111.043,241.151) (111.242,241.062) (111.441,240.973) (111.640,240.886) (111.839,240.800) (112.038,240.714) (112.237,240.630) (112.436,240.546) (112.635,240.463) (112.834,240.381) (113.033,240.300) (113.232,240.219) (113.431,240.140) (113.630,240.061) (113.829,239.983) (114.028,239.906) (114.227,239.830) (114.426,239.755) (114.625,239.680) (114.824,239.606) (115.023,239.533) (115.222,239.461) (115.421,239.390) (115.620,239.319) (115.819,239.249) (116.018,239.180) (116.217,239.112) (116.416,239.044) (116.615,238.978) (116.814,238.912) (117.013,238.846) (117.212,238.782) (117.411,238.718) (117.610,238.655) (117.809,238.592) (118.008,238.531) (118.207,238.470) (118.406,238.410) (118.605,238.350) (118.804,238.291) (119.003,238.233) (119.202,238.176) (119.401,238.119) (119.600,238.063) (119.799,238.008) (119.998,237.953) (120.197,237.899) (120.396,237.846) (120.595,237.793) (120.794,237.741) (120.993,237.690) (121.192,237.639) (121.391,237.589) (121.590,237.540) (121.789,237.491) (121.988,237.443) (122.187,237.395) (122.386,237.349) (122.585,237.303) (122.784,237.257) (122.983,237.212) (123.182,237.168) (123.381,237.124) (123.580,237.081) (123.779,237.039) (123.978,236.997) (124.177,236.956) (124.376,236.915) (124.575,236.875) (124.774,236.835) (124.973,236.797) (125.172,236.758) (125.371,236.721) (125.570,236.684) (125.769,236.647) (125.968,236.611) (126.167,236.576) (126.366,236.541) (126.565,236.507) (126.764,236.473) (126.963,236.440) (127.162,236.408) (127.361,236.376) (127.560,236.345) (127.759,236.314) (127.958,236.284) (128.157,236.254) (128.356,236.225) (128.555,236.196) (128.754,236.168) (128.953,236.141) (129.152,236.114) (129.351,236.088) (129.550,236.062) (129.749,236.037) (129.948,236.012) (130.147,235.988) (130.346,235.964) (130.545,235.941) (130.744,235.918) (130.943,235.896) (131.142,235.875) (131.341,235.854) (131.540,235.833) (131.739,235.813) (131.938,235.794) (132.137,235.775) (132.336,235.756) (132.535,235.739) (132.734,235.721) (132.933,235.704) (133.132,235.688) (133.331,235.672) (133.530,235.657) (133.729,235.642) (133.928,235.628) (134.127,235.614) (134.326,235.601) (134.525,235.588) (134.724,235.576) (134.923,235.564) (135.122,235.553) (135.321,235.542) (135.520,235.532) (135.719,235.522) (135.918,235.513) (136.117,235.504) (136.316,235.496) (136.515,235.488) (136.714,235.481) (136.913,235.474) (137.112,235.468) (137.311,235.462) (137.510,235.457) (137.709,235.453) (137.908,235.448) (138.107,235.445) (138.306,235.441) (138.505,235.439) (138.704,235.436) (138.903,235.435) (139.102,235.433) (139.301,235.433) (139.500,235.432) (139.699,235.433) (139.898,235.433) (140.097,235.435) (140.296,235.436) (140.495,235.439) (140.694,235.441) (140.893,235.445) (141.092,235.448) (141.291,235.453) (141.490,235.457) (141.689,235.462) (141.888,235.468) (142.087,235.474) (142.286,235.481) (142.485,235.488) (142.684,235.496) (142.883,235.504) (143.082,235.513) (143.281,235.522) (143.480,235.532) (143.679,235.542) (143.878,235.553) (144.077,235.564) (144.276,235.576) (144.475,235.588) (144.674,235.601) (144.873,235.614) (145.072,235.628) (145.271,235.642) (145.470,235.657) (145.669,235.672) (145.868,235.688) (146.067,235.704) (146.266,235.721) (146.465,235.739) (146.664,235.756) (146.863,235.775) (147.062,235.794) (147.261,235.813) (147.460,235.833) (147.659,235.854) (147.858,235.875) (148.057,235.896) (148.256,235.918) (148.455,235.941) (148.654,235.964) (148.853,235.988) (149.052,236.012) (149.251,236.037) (149.450,236.062) (149.649,236.088) (149.848,236.114) (150.047,236.141) (150.246,236.168) (150.445,236.196) (150.644,236.225) (150.843,236.254) (151.042,236.284) (151.241,236.314) (151.440,236.345) (151.639,236.376) (151.838,236.408) (152.037,236.440) (152.236,236.473) (152.435,236.507) (152.634,236.541) (152.833,236.576) (153.032,236.611) (153.231,236.647) (153.430,236.684) (153.629,236.721) (153.828,236.758) (154.027,236.797) (154.226,236.835) (154.425,236.875) (154.624,236.915) (154.823,236.956) (155.022,236.997) (155.221,237.039) (155.420,237.081) (155.619,237.124) (155.818,237.168) (156.017,237.212) (156.216,237.257) (156.415,237.303) (156.614,237.349) (156.813,237.395) (157.012,237.443) (157.211,237.491) (157.410,237.540) (157.609,237.589) (157.808,237.639) (158.007,237.690) (158.206,237.741) (158.405,237.793) (158.604,237.846) (158.803,237.899) (159.002,237.953) (159.201,238.008) (159.400,238.063) (159.599,238.119) (159.798,238.176) (159.997,238.233) (160.196,238.291) (160.395,238.350) (160.594,238.410) (160.793,238.470) (160.992,238.531) (161.191,238.592) (161.390,238.655) (161.589,238.718) (161.788,238.782) (161.987,238.846) (162.186,238.912) (162.385,238.978) (162.584,239.044) (162.783,239.112) (162.982,239.180) (163.181,239.249) (163.380,239.319) (163.579,239.390) (163.778,239.461) (163.977,239.533) (164.176,239.606) (164.375,239.680) (164.574,239.755) (164.773,239.830) (164.972,239.906) (165.171,239.983) (165.370,240.061) (165.569,240.140) (165.768,240.219) (165.967,240.300) (166.166,240.381) (166.365,240.463) (166.564,240.546) (166.763,240.630) (166.962,240.714) (167.161,240.800) (167.360,240.886) (167.559,240.973) (167.758,241.062) (167.957,241.151) (168.156,241.241) (168.355,241.331) (168.554,241.423) (168.753,241.516) (168.952,241.610) (169.151,241.704) (169.350,241.800) (169.549,241.897) (169.748,241.994) (169.947,242.092) (170.146,242.192) (170.345,242.292) (170.544,242.394) (170.743,242.496) (170.942,242.600) (171.141,242.704) (171.340,242.810) (171.539,242.916) (171.738,243.024) (171.937,243.132) (172.136,243.242) (172.335,243.353) (172.534,243.464) (172.733,243.577) (172.932,243.691) (173.131,243.806) (173.330,243.922) (173.529,244.039) (173.728,244.158) (173.927,244.277) (174.126,244.398) (174.325,244.520) (174.524,244.642) (174.723,244.766) (174.922,244.892) (175.121,245.018) (175.320,245.146) (175.519,245.275) (175.718,245.404) (175.917,245.536) (176.116,245.668) (176.315,245.802) (176.514,245.937) (176.713,246.073) (176.912,246.210) (177.111,246.349) (177.310,246.489) (177.509,246.630) (177.708,246.773) (177.907,246.916) (178.106,247.062) (178.305,247.208) (178.504,247.356) (178.703,247.505) (178.902,247.656) (179.101,247.808) (179.300,247.961) (179.499,248.115) (179.698,248.272) (179.897,248.429) (180.096,248.588) (180.295,248.748) (180.494,248.910) (180.693,249.073) (180.892,249.238) (181.091,249.404) (181.290,249.571) (181.489,249.740) (181.688,249.911) (181.887,250.083) (182.086,250.257) (182.285,250.432) (182.484,250.608) (182.683,250.787) (182.882,250.966) (183.081,251.148) (183.280,251.331) (183.479,251.515) (183.678,251.701) (183.877,251.889) (184.076,252.078) (184.275,252.269) (184.474,252.462) (184.673,252.656) (184.872,252.852) (185.071,253.050) (185.270,253.249) (185.469,253.450) (185.668,253.652) (185.867,253.857) (186.066,254.063) (186.265,254.271) (186.464,254.480) (186.663,254.692) (186.862,254.905) (187.061,255.120) (187.260,255.336) (187.459,255.555) (187.658,255.775) (187.857,255.997) (188.056,256.221) (188.255,256.447) (188.454,256.674) (188.653,256.903) (188.852,257.135) (189.051,257.368) (189.250,257.603) (189.449,257.839) (189.648,258.078) (189.847,258.319) (190.046,258.561) (190.245,258.806) (190.444,259.052) (190.643,259.300) (190.842,259.550) (191.041,259.802) (191.240,260.056) (191.439,260.312) (191.638,260.570) (191.837,260.830) (192.036,261.092) (192.235,261.356) (192.434,261.622) (192.633,261.889) (192.832,262.159) (193.031,262.431) (193.230,262.705) (193.429,262.980) (193.628,263.258) (193.827,263.538) (194.026,263.819) (194.225,264.103) (194.424,264.388) (194.623,264.676) (194.822,264.966) (195.021,265.257) (195.220,265.551) (195.419,265.847) (195.618,266.144) (195.817,266.444) (196.016,266.745) (196.215,267.049) (196.414,267.354) (196.613,267.662) (196.812,267.971) (197.011,268.283) (197.210,268.596) (197.409,268.911) (197.608,269.228) (197.807,269.548) (198.006,269.869) (198.205,270.192) (198.404,270.516) (198.603,270.843) (198.802,271.172) (199.001,271.502) (199.200,271.834) (199.399,272.169) (199.598,272.505) (199.797,272.842) (199.996,273.182) (200.195,273.523) (200.394,273.866) (200.593,274.211) (200.792,274.557) (200.991,274.905) (201.190,275.255) (201.389,275.607) (201.588,275.960) (201.787,276.314) (201.986,276.670) (202.185,277.028) (202.384,277.387) (202.583,277.748) (202.782,278.110) (202.981,278.473) (203.180,278.838) (203.379,279.205) (203.578,279.572) (203.777,279.941) (203.976,280.311) (204.175,280.682) (204.374,281.054) (204.573,281.428) (204.772,281.803) (204.971,282.178) (205.170,282.555) (205.369,282.932) (205.568,283.311) (205.767,283.690) (205.966,284.070) (206.165,284.451) (206.364,284.833) (206.563,285.215) (206.762,285.598) (206.961,285.981) (207.160,286.365) (207.359,286.749) (207.558,287.134) (207.757,287.519) (207.956,287.904) (208.155,288.290) (208.354,288.676) (208.553,289.061) (208.752,289.447) (208.951,289.833) (209.150,290.218) (209.349,290.604) (209.548,290.989) (209.747,291.374) (209.946,291.759) (210.145,292.143) (210.344,292.526) (210.543,292.909) (210.742,293.292) (210.941,293.673) (211.140,294.054) (211.339,294.434) (211.538,294.813) (211.737,295.191) (211.936,295.568) (212.135,295.944) (212.334,296.318) (212.533,296.692) (212.732,297.063) (212.931,297.434) (213.130,297.803) (213.329,298.170) (213.528,298.535) (213.727,298.899) (213.926,299.261) (214.125,299.621) (214.324,299.979) (214.523,300.335) (214.722,300.688) (214.921,301.040) (215.120,301.389) (215.319,301.736) (215.518,302.080) (215.717,302.422) (215.916,302.761) (216.115,303.098) (216.314,303.432) (216.513,303.763) (216.712,304.091) (216.911,304.416) (217.110,304.739) (217.309,305.058) (217.508,305.374) (217.707,305.686) (217.906,305.996) (218.105,306.302) (218.304,306.605) (218.503,306.904) (218.702,307.200) (218.901,307.492) (219.100,307.781) (219.299,308.065) (219.498,308.347) (219.697,308.624) (219.896,308.898) (220.095,309.167) (220.294,309.433) (220.493,309.695) (220.692,309.953) (220.891,310.207) (221.090,310.457) (221.289,310.702) (221.488,310.944) (221.687,311.181) (221.886,311.415) (222.085,311.644) (222.284,311.868) (222.483,312.089) (222.682,312.305) (222.881,312.517) (223.080,312.725) (223.279,312.929) (223.478,313.128) (223.677,313.323) (223.876,313.513) (224.075,313.699) (224.274,313.881) (224.473,314.059) (224.672,314.232) (224.871,314.401) (225.070,314.566) (225.269,314.727) (225.468,314.883) (225.667,315.035) (225.866,315.183) (226.065,315.327) (226.264,315.466) (226.463,315.602) (226.662,315.733) (226.861,315.861) (227.060,315.984) (227.259,316.103) (227.458,316.219) (227.657,316.330) (227.856,316.438) (228.055,316.542) (228.254,316.642) (228.453,316.738) (228.652,316.831) (228.851,316.920) (229.050,317.005) (229.249,317.088) (229.448,317.166) (229.647,317.242) (229.846,317.314) (230.045,317.382) (230.244,317.448) (230.443,317.510) (230.642,317.570) (230.841,317.626) (231.040,317.680) (231.239,317.731) (231.438,317.779) (231.637,317.824) (231.836,317.867) (232.035,317.907) (232.234,317.944) (232.433,317.980) (232.632,318.013) (232.831,318.044) (233.030,318.072) (233.229,318.099) (233.428,318.123) (233.627,318.146) (233.826,318.167) (234.025,318.186) (234.224,318.204) (234.423,318.220) (234.622,318.234) (234.821,318.247) (235.020,318.259) (235.219,318.269) (235.418,318.278) (235.617,318.286) (235.816,318.293) (236.015,318.299) (236.214,318.304) (236.413,318.309) (236.612,318.313) (236.811,318.316) (237.010,318.318) (237.209,318.320) (237.408,318.321) (237.607,318.323) (237.806,318.323) (238.005,318.324) (238.204,318.324) (238.403,318.324) (238.602,318.325) (238.801,318.325) (239.000,318.325)};
\end{scope}
\node[anchor=west,font=\small\bfseries,text=navy] at (264,348) {(d) R: rate at each trip-average speed};
\draw[rulegray,line width=.25pt] (299,212.000) -- (498,212.000);
\node[anchor=east,font=\scriptsize] at (295,212.000) {0};
\draw[rulegray,line width=.25pt] (299,238.364) -- (498,238.364);
\node[anchor=east,font=\scriptsize] at (295,238.364) {5};
\draw[rulegray,line width=.25pt] (299,264.727) -- (498,264.727);
\node[anchor=east,font=\scriptsize] at (295,264.727) {10};
\draw[rulegray,line width=.25pt] (299,291.091) -- (498,291.091);
\node[anchor=east,font=\scriptsize] at (295,291.091) {15};
\draw[rulegray,line width=.25pt] (299,317.455) -- (498,317.455);
\node[anchor=east,font=\scriptsize] at (295,317.455) {20};
\draw[black,line width=.35pt] (299.000,212) -- (299.000,209);
\node[anchor=north,font=\scriptsize] at (299.000,208) {0};
\draw[black,line width=.35pt] (343.725,212) -- (343.725,209);
\node[anchor=north,font=\scriptsize] at (343.725,208) {1};
\draw[black,line width=.35pt] (388.449,212) -- (388.449,209);
\node[anchor=north,font=\scriptsize] at (388.449,208) {2};
\draw[black,line width=.35pt] (433.174,212) -- (433.174,209);
\node[anchor=north,font=\scriptsize] at (433.174,208) {3};
\draw[black,line width=.35pt] (477.898,212) -- (477.898,209);
\node[anchor=north,font=\scriptsize] at (477.898,208) {4};
\draw[black,line width=.45pt] (299,328) -- (299,212) -- (498,212);
\node[anchor=north,font=\scriptsize] at (398.5,193) {Elapsed entry time (h)};
\node[rotate=90,font=\scriptsize] at (266,270.0) {CO$_2$ rate (kg/veh h)};
\begin{scope}
\clip (299,212) rectangle (498,328);
\draw[methodR,line width=1.05pt,] plot coordinates {(299.000,320.701) (299.199,320.701) (299.398,320.701) (299.597,320.701) (299.796,320.701) (299.995,320.701) (300.194,320.701) (300.393,320.701) (300.592,320.701) (300.791,320.701) (300.990,320.701) (301.189,320.701) (301.388,320.701) (301.587,320.701) (301.786,320.701) (301.985,320.701) (302.184,320.701) (302.383,295.549) (302.582,295.549) (302.781,295.549) (302.980,295.549) (303.179,295.549) (303.378,295.549) (303.577,295.549) (303.776,295.549) (303.975,295.549) (304.174,295.549) (304.373,295.549) (304.572,295.549) (304.771,295.549) (304.970,295.549) (305.169,295.549) (305.368,295.549) (305.567,295.549) (305.766,295.549) (305.965,295.549) (306.164,295.549) (306.363,295.549) (306.562,295.549) (306.761,295.549) (306.960,295.549) (307.159,295.549) (307.358,295.549) (307.557,295.549) (307.756,295.549) (307.955,295.549) (308.154,295.549) (308.353,295.549) (308.552,295.549) (308.751,295.549) (308.950,295.549) (309.149,295.549) (309.348,295.549) (309.547,295.549) (309.746,295.549) (309.945,295.549) (310.144,295.549) (310.343,295.549) (310.542,295.549) (310.741,295.549) (310.940,295.549) (311.139,295.549) (311.338,295.549) (311.537,295.549) (311.736,295.549) (311.935,295.549) (312.134,295.549) (312.333,295.549) (312.532,295.549) (312.731,295.549) (312.930,295.549) (313.129,295.549) (313.328,295.549) (313.527,295.549) (313.726,295.549) (313.925,295.549) (314.124,295.549) (314.323,295.549) (314.522,295.549) (314.721,295.549) (314.920,295.549) (315.119,295.549) (315.318,295.549) (315.517,295.549) (315.716,295.549) (315.915,295.549) (316.114,295.549) (316.313,295.549) (316.512,295.549) (316.711,295.549) (316.910,295.549) (317.109,295.549) (317.308,295.549) (317.507,295.549) (317.706,295.549) (317.905,295.549) (318.104,295.549) (318.303,295.549) (318.502,295.549) (318.701,295.549) (318.900,295.549) (319.099,295.549) (319.298,295.549) (319.497,295.549) (319.696,295.549) (319.895,295.549) (320.094,295.549) (320.293,295.549) (320.492,295.549) (320.691,295.549) (320.890,295.549) (321.089,295.549) (321.288,295.549) (321.487,295.549) (321.686,295.549) (321.885,295.549) (322.084,295.549) (322.283,295.549) (322.482,295.549) (322.681,295.549) (322.880,295.549) (323.079,295.549) (323.278,295.549) (323.477,295.549) (323.676,295.549) (323.875,295.549) (324.074,295.549) (324.273,295.549) (324.472,295.549) (324.671,295.549) (324.870,295.549) (325.069,295.549) (325.268,295.549) (325.467,295.549) (325.666,295.549) (325.865,295.549) (326.064,295.549) (326.263,295.549) (326.462,278.106) (326.661,278.106) (326.860,278.106) (327.059,278.106) (327.258,278.106) (327.457,278.106) (327.656,278.106) (327.855,278.106) (328.054,278.106) (328.253,278.106) (328.452,278.106) (328.651,278.106) (328.850,263.280) (329.049,263.280) (329.248,263.280) (329.447,263.280) (329.646,263.280) (329.845,263.280) (330.044,263.280) (330.243,263.280) (330.442,263.280) (330.641,263.280) (330.840,263.280) (331.039,263.280) (331.238,263.280) (331.437,263.280) (331.636,263.280) (331.835,263.280) (332.034,263.280) (332.233,263.280) (332.432,263.280) (332.631,263.280) (332.830,263.280) (333.029,263.280) (333.228,263.280) (333.427,263.280) (333.626,263.280) (333.825,263.280) (334.024,263.280) (334.223,263.280) (334.422,263.280) (334.621,263.280) (334.820,263.280) (335.019,263.280) (335.218,263.280) (335.417,263.280) (335.616,263.280) (335.815,263.280) (336.014,263.280) (336.213,263.280) (336.412,263.280) (336.611,263.280) (336.810,263.280) (337.009,263.280) (337.208,263.280) (337.407,263.280) (337.606,263.280) (337.805,263.280) (338.004,263.280) (338.203,263.280) (338.402,263.280) (338.601,263.280) (338.800,263.280) (338.999,263.280) (339.198,263.280) (339.397,263.280) (339.596,263.280) (339.795,263.280) (339.994,263.280) (340.193,263.280) (340.392,263.280) (340.591,263.280) (340.790,263.280) (340.989,263.280) (341.188,263.280) (341.387,263.280) (341.586,253.414) (341.785,253.414) (341.984,253.414) (342.183,253.414) (342.382,253.414) (342.581,253.414) (342.780,253.414) (342.979,253.414) (343.178,253.414) (343.377,253.414) (343.576,253.414) (343.775,253.414) (343.974,253.414) (344.173,253.414) (344.372,253.414) (344.571,253.414) (344.770,253.414) (344.969,253.414) (345.168,253.414) (345.367,253.414) (345.566,253.414) (345.765,253.414) (345.964,253.414) (346.163,253.414) (346.362,253.414) (346.561,253.414) (346.760,253.414) (346.959,253.414) (347.158,253.414) (347.357,253.414) (347.556,253.414) (347.755,253.414) (347.954,253.414) (348.153,253.414) (348.352,253.414) (348.551,253.414) (348.750,253.414) (348.949,253.414) (349.148,253.414) (349.347,253.414) (349.546,253.414) (349.745,253.414) (349.944,253.414) (350.143,253.414) (350.342,253.414) (350.541,253.414) (350.740,253.414) (350.939,253.414) (351.138,253.414) (351.337,253.414) (351.536,253.414) (351.735,253.414) (351.934,253.414) (352.133,253.414) (352.332,253.414) (352.531,253.414) (352.730,248.550) (352.929,248.550) (353.128,248.550) (353.327,248.550) (353.526,248.550) (353.725,248.550) (353.924,248.550) (354.123,248.550) (354.322,248.550) (354.521,248.550) (354.720,248.550) (354.919,248.550) (355.118,248.550) (355.317,248.550) (355.516,248.550) (355.715,248.550) (355.914,248.550) (356.113,248.550) (356.312,248.550) (356.511,248.550) (356.710,248.550) (356.909,248.550) (357.108,248.550) (357.307,248.550) (357.506,248.550) (357.705,248.550) (357.904,248.550) (358.103,248.550) (358.302,248.550) (358.501,248.550) (358.700,248.550) (358.899,248.550) (359.098,248.550) (359.297,248.550) (359.496,248.550) (359.695,248.550) (359.894,248.550) (360.093,248.550) (360.292,248.550) (360.491,248.550) (360.690,248.550) (360.889,248.550) (361.088,248.550) (361.287,248.550) (361.486,248.550) (361.685,248.550) (361.884,248.550) (362.083,248.550) (362.282,248.550) (362.481,248.550) (362.680,248.550) (362.879,248.550) (363.078,248.550) (363.277,248.550) (363.476,248.550) (363.675,248.550) (363.874,248.550) (364.073,248.550) (364.272,248.550) (364.471,248.550) (364.670,248.550) (364.869,248.550) (365.068,248.550) (365.267,248.550) (365.466,248.550) (365.665,248.550) (365.864,248.550) (366.063,248.550) (366.262,248.550) (366.461,248.550) (366.660,248.550) (366.859,248.550) (367.058,248.550) (367.257,248.550) (367.456,248.550) (367.655,248.550) (367.854,248.550) (368.053,248.550) (368.252,248.550) (368.451,248.550) (368.650,248.550) (368.849,248.550) (369.048,248.550) (369.247,248.550) (369.446,248.550) (369.645,248.550) (369.844,248.550) (370.043,248.550) (370.242,248.550) (370.441,248.550) (370.640,248.550) (370.839,248.550) (371.038,248.550) (371.237,248.550) (371.436,248.550) (371.635,248.550) (371.834,248.550) (372.033,248.550) (372.232,248.550) (372.431,248.550) (372.630,248.550) (372.829,248.550) (373.028,248.550) (373.227,248.550) (373.426,248.550) (373.625,248.550) (373.824,248.550) (374.023,248.550) (374.222,248.550) (374.421,248.550) (374.620,248.550) (374.819,248.550) (375.018,248.550) (375.217,248.550) (375.416,248.550) (375.615,248.550) (375.814,248.550) (376.013,248.550) (376.212,248.550) (376.411,248.550) (376.610,248.550) (376.809,248.550) (377.008,248.550) (377.207,248.550) (377.406,248.550) (377.605,248.550) (377.804,248.550) (378.003,248.550) (378.202,248.550) (378.401,248.550) (378.600,248.550) (378.799,248.550) (378.998,248.550) (379.197,248.550) (379.396,248.550) (379.595,248.550) (379.794,248.550) (379.993,248.550) (380.192,248.550) (380.391,248.550) (380.590,248.550) (380.789,248.550) (380.988,248.550) (381.187,248.550) (381.386,248.550) (381.585,248.550) (381.784,248.550) (381.983,248.550) (382.182,248.550) (382.381,248.550) (382.580,248.550) (382.779,248.550) (382.978,248.550) (383.177,248.550) (383.376,248.550) (383.575,248.550) (383.774,248.550) (383.973,248.550) (384.172,248.550) (384.371,248.550) (384.570,248.550) (384.769,248.550) (384.968,248.550) (385.167,248.550) (385.366,248.550) (385.565,248.550) (385.764,248.550) (385.963,248.550) (386.162,248.550) (386.361,248.550) (386.560,248.550) (386.759,248.550) (386.958,248.550) (387.157,248.550) (387.356,248.550) (387.555,248.550) (387.754,248.550) (387.953,248.550) (388.152,248.550) (388.351,248.550) (388.550,248.550) (388.749,248.550) (388.948,248.550) (389.147,248.550) (389.346,248.550) (389.545,248.550) (389.744,248.550) (389.943,248.550) (390.142,248.550) (390.341,248.550) (390.540,248.550) (390.739,248.550) (390.938,248.550) (391.137,248.550) (391.336,248.550) (391.535,248.550) (391.734,248.550) (391.933,248.550) (392.132,248.550) (392.331,248.550) (392.530,248.550) (392.729,248.550) (392.928,248.550) (393.127,248.550) (393.326,248.550) (393.525,248.550) (393.724,248.550) (393.923,248.550) (394.122,248.550) (394.321,248.550) (394.520,248.550) (394.719,248.550) (394.918,248.550) (395.117,248.550) (395.316,248.550) (395.515,248.550) (395.714,248.550) (395.913,248.550) (396.112,248.550) (396.311,248.550) (396.510,248.550) (396.709,248.550) (396.908,248.550) (397.107,248.550) (397.306,248.550) (397.505,248.550) (397.704,248.550) (397.903,248.550) (398.102,248.550) (398.301,248.550) (398.500,248.550) (398.699,248.550) (398.898,248.550) (399.097,248.550) (399.296,248.550) (399.495,248.550) (399.694,248.550) (399.893,248.550) (400.092,248.550) (400.291,248.550) (400.490,248.550) (400.689,248.550) (400.888,248.550) (401.087,248.550) (401.286,248.550) (401.485,248.550) (401.684,248.550) (401.883,248.550) (402.082,248.550) (402.281,248.550) (402.480,248.550) (402.679,248.550) (402.878,248.550) (403.077,248.550) (403.276,248.550) (403.475,248.550) (403.674,248.550) (403.873,248.550) (404.072,248.550) (404.271,248.550) (404.470,248.550) (404.669,248.550) (404.868,248.550) (405.067,248.550) (405.266,248.550) (405.465,248.550) (405.664,248.550) (405.863,248.550) (406.062,248.550) (406.261,248.550) (406.460,248.550) (406.659,248.550) (406.858,248.550) (407.057,248.550) (407.256,248.550) (407.455,248.550) (407.654,248.550) (407.853,248.550) (408.052,248.550) (408.251,248.550) (408.450,248.550) (408.649,248.550) (408.848,248.550) (409.047,248.550) (409.246,248.550) (409.445,248.550) (409.644,248.550) (409.843,248.550) (410.042,248.550) (410.241,248.550) (410.440,248.550) (410.639,248.550) (410.838,248.550) (411.037,248.550) (411.236,248.550) (411.435,248.550) (411.634,248.550) (411.833,248.550) (412.032,248.550) (412.231,248.550) (412.430,248.550) (412.629,248.550) (412.828,248.550) (413.027,248.550) (413.226,248.550) (413.425,248.550) (413.624,248.550) (413.823,248.550) (414.022,248.550) (414.221,248.550) (414.420,248.550) (414.619,248.550) (414.818,248.550) (415.017,248.550) (415.216,248.550) (415.415,248.550) (415.614,248.550) (415.813,248.550) (416.012,248.550) (416.211,248.550) (416.410,248.550) (416.609,248.550) (416.808,248.550) (417.007,248.550) (417.206,248.550) (417.405,248.550) (417.604,248.550) (417.803,248.550) (418.002,248.550) (418.201,248.550) (418.400,248.550) (418.599,248.550) (418.798,248.550) (418.997,248.550) (419.196,248.550) (419.395,248.550) (419.594,248.550) (419.793,248.550) (419.992,248.550) (420.191,248.550) (420.390,248.550) (420.589,248.550) (420.788,248.550) (420.987,248.550) (421.186,248.550) (421.385,248.550) (421.584,248.550) (421.783,248.550) (421.982,248.550) (422.181,248.550) (422.380,248.550) (422.579,248.550) (422.778,248.550) (422.977,248.550) (423.176,248.550) (423.375,248.550) (423.574,248.550) (423.773,248.550) (423.972,248.550) (424.171,248.550) (424.370,248.550) (424.569,248.550) (424.768,248.550) (424.967,248.550) (425.166,248.550) (425.365,248.550) (425.564,248.550) (425.763,248.550) (425.962,248.550) (426.161,248.550) (426.360,248.550) (426.559,248.550) (426.758,248.550) (426.957,248.550) (427.156,248.550) (427.355,248.550) (427.554,248.550) (427.753,248.550) (427.952,248.550) (428.151,248.550) (428.350,248.550) (428.549,248.550) (428.748,248.550) (428.947,248.550) (429.146,248.550) (429.345,248.550) (429.544,248.550) (429.743,248.550) (429.942,248.550) (430.141,248.550) (430.340,248.550) (430.539,248.550) (430.738,248.550) (430.937,248.550) (431.136,248.550) (431.335,248.550) (431.534,248.550) (431.733,248.550) (431.932,248.550) (432.131,248.550) (432.330,248.550) (432.529,248.550) (432.728,248.550) (432.927,248.550) (433.126,248.550) (433.325,248.550) (433.524,248.550) (433.723,248.550) (433.922,248.550) (434.121,248.550) (434.320,248.550) (434.519,248.550) (434.718,248.550) (434.917,248.550) (435.116,248.550) (435.315,248.550) (435.514,248.550) (435.713,248.550) (435.912,248.550) (436.111,248.550) (436.310,248.550) (436.509,248.550) (436.708,248.550) (436.907,248.550) (437.106,248.550) (437.305,248.550) (437.504,248.550) (437.703,248.550) (437.902,248.550) (438.101,248.550) (438.300,248.550) (438.499,248.550) (438.698,248.550) (438.897,248.550) (439.096,248.550) (439.295,248.550) (439.494,248.550) (439.693,248.550) (439.892,248.550) (440.091,248.550) (440.290,248.550) (440.489,248.550) (440.688,248.550) (440.887,248.550) (441.086,248.550) (441.285,248.550) (441.484,248.550) (441.683,248.550) (441.882,248.550) (442.081,248.550) (442.280,248.550) (442.479,248.550) (442.678,248.550) (442.877,248.550) (443.076,248.550) (443.275,248.550) (443.474,248.550) (443.673,248.550) (443.872,248.550) (444.071,248.550) (444.270,248.550) (444.469,253.414) (444.668,253.414) (444.867,253.414) (445.066,253.414) (445.265,253.414) (445.464,253.414) (445.663,253.414) (445.862,253.414) (446.061,253.414) (446.260,253.414) (446.459,253.414) (446.658,253.414) (446.857,253.414) (447.056,253.414) (447.255,253.414) (447.454,253.414) (447.653,253.414) (447.852,253.414) (448.051,253.414) (448.250,253.414) (448.449,253.414) (448.648,253.414) (448.847,253.414) (449.046,253.414) (449.245,253.414) (449.444,253.414) (449.643,253.414) (449.842,253.414) (450.041,253.414) (450.240,253.414) (450.439,253.414) (450.638,253.414) (450.837,253.414) (451.036,253.414) (451.235,253.414) (451.434,253.414) (451.633,253.414) (451.832,253.414) (452.031,253.414) (452.230,253.414) (452.429,253.414) (452.628,253.414) (452.827,253.414) (453.026,253.414) (453.225,253.414) (453.424,253.414) (453.623,253.414) (453.822,253.414) (454.021,253.414) (454.220,253.414) (454.419,253.414) (454.618,253.414) (454.817,253.414) (455.016,253.414) (455.215,253.414) (455.414,253.414) (455.613,263.280) (455.812,263.280) (456.011,263.280) (456.210,263.280) (456.409,263.280) (456.608,263.280) (456.807,263.280) (457.006,263.280) (457.205,263.280) (457.404,263.280) (457.603,263.280) (457.802,263.280) (458.001,263.280) (458.200,263.280) (458.399,263.280) (458.598,263.280) (458.797,263.280) (458.996,263.280) (459.195,263.280) (459.394,263.280) (459.593,263.280) (459.792,263.280) (459.991,263.280) (460.190,263.280) (460.389,263.280) (460.588,263.280) (460.787,263.280) (460.986,263.280) (461.185,263.280) (461.384,263.280) (461.583,263.280) (461.782,263.280) (461.981,263.280) (462.180,263.280) (462.379,263.280) (462.578,263.280) (462.777,263.280) (462.976,263.280) (463.175,263.280) (463.374,263.280) (463.573,263.280) (463.772,263.280) (463.971,263.280) (464.170,263.280) (464.369,263.280) (464.568,263.280) (464.767,263.280) (464.966,263.280) (465.165,263.280) (465.364,263.280) (465.563,263.280) (465.762,263.280) (465.961,263.280) (466.160,263.280) (466.359,263.280) (466.558,263.280) (466.757,263.280) (466.956,263.280) (467.155,263.280) (467.354,263.280) (467.553,263.280) (467.752,263.280) (467.951,263.280) (468.150,263.280) (468.349,278.106) (468.548,278.106) (468.747,278.106) (468.946,278.106) (469.145,278.106) (469.344,278.106) (469.543,278.106) (469.742,278.106) (469.941,278.106) (470.140,278.106) (470.339,278.106) (470.538,278.106) (470.737,295.549) (470.936,295.549) (471.135,295.549) (471.334,295.549) (471.533,295.549) (471.732,295.549) (471.931,295.549) (472.130,295.549) (472.329,295.549) (472.528,295.549) (472.727,295.549) (472.926,295.549) (473.125,295.549) (473.324,295.549) (473.523,295.549) (473.722,295.549) (473.921,295.549) (474.120,295.549) (474.319,295.549) (474.518,295.549) (474.717,295.549) (474.916,295.549) (475.115,295.549) (475.314,295.549) (475.513,295.549) (475.712,295.549) (475.911,295.549) (476.110,295.549) (476.309,295.549) (476.508,295.549) (476.707,295.549) (476.906,295.549) (477.105,295.549) (477.304,295.549) (477.503,295.549) (477.702,295.549) (477.901,295.549) (478.100,295.549) (478.299,295.549) (478.498,295.549) (478.697,295.549) (478.896,295.549) (479.095,295.549) (479.294,295.549) (479.493,295.549) (479.692,295.549) (479.891,295.549) (480.090,295.549) (480.289,295.549) (480.488,295.549) (480.687,295.549) (480.886,295.549) (481.085,295.549) (481.284,295.549) (481.483,295.549) (481.682,295.549) (481.881,295.549) (482.080,295.549) (482.279,295.549) (482.478,295.549) (482.677,295.549) (482.876,295.549) (483.075,295.549) (483.274,295.549) (483.473,295.549) (483.672,295.549) (483.871,295.549) (484.070,295.549) (484.269,295.549) (484.468,295.549) (484.667,295.549) (484.866,295.549) (485.065,295.549) (485.264,295.549) (485.463,295.549) (485.662,295.549) (485.861,295.549) (486.060,295.549) (486.259,295.549) (486.458,295.549) (486.657,295.549) (486.856,295.549) (487.055,295.549) (487.254,295.549) (487.453,295.549) (487.652,295.549) (487.851,295.549) (488.050,295.549) (488.249,295.549) (488.448,295.549) (488.647,295.549) (488.846,295.549) (489.045,295.549) (489.244,295.549) (489.443,295.549) (489.642,295.549) (489.841,295.549) (490.040,295.549) (490.239,295.549) (490.438,295.549) (490.637,295.549) (490.836,295.549) (491.035,295.549) (491.234,295.549) (491.433,295.549) (491.632,295.549) (491.831,295.549) (492.030,295.549) (492.229,295.549) (492.428,295.549) (492.627,295.549) (492.826,295.549) (493.025,295.549) (493.224,295.549) (493.423,295.549) (493.622,295.549) (493.821,295.549) (494.020,295.549) (494.219,295.549) (494.418,295.549) (494.617,295.549) (494.816,320.701) (495.015,320.701) (495.214,320.701) (495.413,320.701) (495.612,320.701) (495.811,320.701) (496.010,320.701) (496.209,320.701) (496.408,320.701) (496.607,320.701) (496.806,320.701) (497.005,320.701) (497.204,320.701) (497.403,320.701) (497.602,320.701) (497.801,320.701) (498.000,320.701)};
\end{scope}
\node[anchor=west,font=\small\bfseries,text=navy] at (5,165) {(e) M and A: same two operating speeds};
\draw[rulegray,line width=.25pt] (40,29.000) -- (239,29.000);
\node[anchor=east,font=\scriptsize] at (36,29.000) {0};
\draw[rulegray,line width=.25pt] (40,62.143) -- (239,62.143);
\node[anchor=east,font=\scriptsize] at (36,62.143) {20};
\draw[rulegray,line width=.25pt] (40,95.286) -- (239,95.286);
\node[anchor=east,font=\scriptsize] at (36,95.286) {40};
\draw[rulegray,line width=.25pt] (40,128.429) -- (239,128.429);
\node[anchor=east,font=\scriptsize] at (36,128.429) {60};
\draw[black,line width=.35pt] (40.000,29) -- (40.000,26);
\node[anchor=north,font=\scriptsize] at (40.000,25) {0};
\draw[black,line width=.35pt] (79.020,29) -- (79.020,26);
\node[anchor=north,font=\scriptsize] at (79.020,25) {0.5};
\draw[black,line width=.35pt] (118.039,29) -- (118.039,26);
\node[anchor=north,font=\scriptsize] at (118.039,25) {1};
\draw[black,line width=.35pt] (157.059,29) -- (157.059,26);
\node[anchor=north,font=\scriptsize] at (157.059,25) {1.5};
\draw[black,line width=.35pt] (196.078,29) -- (196.078,26);
\node[anchor=north,font=\scriptsize] at (196.078,25) {2};
\draw[black,line width=.35pt] (235.098,29) -- (235.098,26);
\node[anchor=north,font=\scriptsize] at (235.098,25) {2.5};
\draw[black,line width=.45pt] (40,145) -- (40,29) -- (239,29);
\node[anchor=north,font=\scriptsize] at (139.5,10) {Elapsed time within one trip (min)};
\node[rotate=90,font=\scriptsize] at (7,87.0) {Speed (mph)};
\begin{scope}
\clip (40,29) rectangle (239,145);
\draw[methodM,line width=1.05pt,] plot coordinates {(40.000,135.325) (62.231,135.325) (62.231,51.556) (210.411,51.556) (210.411,135.325) (232.642,135.325)};
\draw[methodA,line width=1.05pt,dashed] plot coordinates {(40.000,135.325) (62.231,135.325) (62.231,51.556) (210.411,51.556) (210.411,135.325) (232.642,135.325)};
\end{scope}
\node[anchor=west,font=\small\bfseries,text=navy] at (264,165) {(f) M and A: lookup versus cubic rates};
\draw[rulegray,line width=.25pt] (299,29.000) -- (498,29.000);
\node[anchor=east,font=\scriptsize] at (295,29.000) {0};
\draw[rulegray,line width=.25pt] (299,55.364) -- (498,55.364);
\node[anchor=east,font=\scriptsize] at (295,55.364) {5};
\draw[rulegray,line width=.25pt] (299,81.727) -- (498,81.727);
\node[anchor=east,font=\scriptsize] at (295,81.727) {10};
\draw[rulegray,line width=.25pt] (299,108.091) -- (498,108.091);
\node[anchor=east,font=\scriptsize] at (295,108.091) {15};
\draw[rulegray,line width=.25pt] (299,134.455) -- (498,134.455);
\node[anchor=east,font=\scriptsize] at (295,134.455) {20};
\draw[black,line width=.35pt] (299.000,29) -- (299.000,26);
\node[anchor=north,font=\scriptsize] at (299.000,25) {0};
\draw[black,line width=.35pt] (338.020,29) -- (338.020,26);
\node[anchor=north,font=\scriptsize] at (338.020,25) {0.5};
\draw[black,line width=.35pt] (377.039,29) -- (377.039,26);
\node[anchor=north,font=\scriptsize] at (377.039,25) {1};
\draw[black,line width=.35pt] (416.059,29) -- (416.059,26);
\node[anchor=north,font=\scriptsize] at (416.059,25) {1.5};
\draw[black,line width=.35pt] (455.078,29) -- (455.078,26);
\node[anchor=north,font=\scriptsize] at (455.078,25) {2};
\draw[black,line width=.35pt] (494.098,29) -- (494.098,26);
\node[anchor=north,font=\scriptsize] at (494.098,25) {2.5};
\draw[black,line width=.45pt] (299,145) -- (299,29) -- (498,29);
\node[anchor=north,font=\scriptsize] at (398.5,10) {Elapsed time within one trip (min)};
\node[rotate=90,font=\scriptsize] at (266,87.0) {CO$_2$ rate (kg/veh h)};
\begin{scope}
\clip (299,29) rectangle (498,145);
\draw[methodM,line width=1.05pt,] plot coordinates {(299.000,137.701) (321.231,137.701) (321.231,65.550) (469.411,65.550) (469.411,137.701) (491.642,137.701)};
\draw[methodA,line width=1.05pt,dashed] plot coordinates {(299.000,124.818) (321.231,124.818) (321.231,64.224) (469.411,64.224) (469.411,124.818) (491.642,124.818)};
\end{scope}
\end{tikzpicture}
}\end{center}
\textbf{How to read the panels.}
(a,b) O uses the recorded five-minute link-speed proxies; S holds their harmonic mean constant.
(c,d) R assigns a whole-link average speed to each entering group over its fitted 4.449 h duration.
This is not a vehicle's within-trip speed history. (e,f) M and A share a 2.469 min mean trip and identical speeds;
only the rate function changes. Free-flow time is split equally before and after the queue for illustration.
This ordering is not identified by the method and does not change its total; instantaneous transitions add no acceleration emissions.

\notebox{\textbf{A rate curve is not an episode emission total.}
O's rate must be weighted by $d_iL/v_i$; R's by each group's count and trip time;
M and A by free/queue time and total demand. The area under (b) or (d) over the displayed horizontal axis
alone is not the episode total. The cumulative contributions on the next page include the required weights.}

\newpage
{\fontsize{19}{22}\selectfont\bfseries\color{navy}From speed choice to emission totals}\par
{\color{muted}Detector 78: lookup sensitivity and correctly weighted cumulative contributions.}\par
\vspace{5pt}
\legendkey{methodO}{O: recorded}\quad\legendkey{methodS}{S: mean speed}\quad\legendkey{methodR}{R: reconstructed}\quad\legendkey{methodM}{M: two-speed}\quad\legendkey{methodA}{A: cubic}
\par\vspace{8pt}\begin{center}\resizebox{\linewidth}{!}{\begin{tikzpicture}[x=1pt,y=1pt]
\useasboundingbox (0,0) rectangle (509,485);
\node[anchor=west,font=\small\bfseries,text=navy] at (5,476) {(a) CO$_2$: speed-to-rate mapping};
\draw[rulegray,line width=.25pt] (40,304.000) -- (239,304.000);
\node[anchor=east,font=\scriptsize] at (36,304.000) {0};
\draw[rulegray,line width=.25pt] (40,354.667) -- (239,354.667);
\node[anchor=east,font=\scriptsize] at (36,354.667) {10};
\draw[rulegray,line width=.25pt] (40,405.333) -- (239,405.333);
\node[anchor=east,font=\scriptsize] at (36,405.333) {20};
\draw[rulegray,line width=.25pt] (40,456.000) -- (239,456.000);
\node[anchor=east,font=\scriptsize] at (36,456.000) {30};
\draw[black,line width=.35pt] (54.214,304) -- (54.214,301);
\node[anchor=north,font=\scriptsize] at (54.214,300) {10};
\draw[black,line width=.35pt] (96.857,304) -- (96.857,301);
\node[anchor=north,font=\scriptsize] at (96.857,300) {25};
\draw[black,line width=.35pt] (139.500,304) -- (139.500,301);
\node[anchor=north,font=\scriptsize] at (139.500,300) {40};
\draw[black,line width=.35pt] (182.143,304) -- (182.143,301);
\node[anchor=north,font=\scriptsize] at (182.143,300) {55};
\draw[black,line width=.35pt] (224.786,304) -- (224.786,301);
\node[anchor=north,font=\scriptsize] at (224.786,300) {70};
\draw[black,line width=.45pt] (40,456) -- (40,304) -- (239,304);
\node[anchor=north,font=\scriptsize] at (139.5,285) {Speed (mph)};
\node[rotate=90,font=\scriptsize] at (7,380.0) {Rate (kg/veh h)};
\begin{scope}
\clip (40,304) rectangle (239,456);
\draw[methodO,line width=1.05pt,] plot coordinates {(40.000,339.121) (40.142,339.121) (40.284,339.121) (40.426,339.121) (40.569,339.121) (40.711,339.121) (40.853,339.121) (40.995,339.121) (41.137,339.121) (41.279,339.121) (41.421,339.121) (41.564,339.121) (41.706,339.121) (41.848,339.121) (41.990,339.121) (42.132,339.121) (42.274,339.121) (42.416,339.121) (42.559,339.121) (42.701,339.121) (42.843,339.121) (42.985,339.121) (43.127,339.121) (43.269,339.121) (43.411,339.121) (43.554,339.121) (43.696,339.121) (43.838,339.121) (43.980,339.121) (44.122,339.121) (44.264,339.121) (44.406,339.121) (44.549,339.121) (44.691,339.121) (44.833,339.121) (44.975,339.121) (45.117,339.121) (45.259,339.121) (45.401,339.121) (45.544,339.121) (45.686,339.121) (45.828,339.121) (45.970,339.121) (46.112,339.121) (46.254,339.121) (46.396,339.121) (46.539,339.121) (46.681,339.121) (46.823,339.121) (46.965,339.121) (47.107,339.121) (47.249,339.121) (47.391,339.121) (47.534,339.121) (47.676,339.121) (47.818,339.121) (47.960,339.121) (48.102,339.121) (48.244,339.121) (48.386,339.121) (48.529,339.121) (48.671,339.121) (48.813,339.121) (48.955,339.121) (49.097,339.121) (49.239,339.121) (49.381,339.121) (49.524,339.121) (49.666,339.121) (49.808,339.121) (49.950,339.121) (50.092,339.121) (50.234,339.121) (50.376,339.121) (50.519,339.121) (50.661,339.121) (50.803,339.121) (50.945,339.121) (51.087,339.121) (51.229,339.121) (51.371,339.121) (51.514,339.121) (51.656,339.121) (51.798,339.121) (51.940,339.121) (52.082,339.121) (52.224,339.121) (52.366,339.121) (52.509,339.121) (52.651,339.121) (52.793,339.121) (52.935,339.121) (53.077,339.121) (53.219,339.121) (53.361,339.121) (53.504,339.121) (53.646,339.121) (53.788,339.121) (53.930,339.121) (54.072,339.121) (54.214,339.121) (54.356,339.121) (54.499,339.121) (54.641,339.121) (54.783,339.121) (54.925,339.121) (55.067,339.121) (55.209,339.121) (55.351,339.121) (55.494,339.121) (55.636,339.121) (55.778,339.121) (55.920,339.121) (56.062,339.121) (56.204,339.121) (56.346,339.121) (56.489,339.121) (56.631,339.121) (56.773,339.121) (56.915,339.121) (57.057,339.121) (57.199,339.121) (57.341,339.121) (57.484,339.121) (57.626,339.121) (57.768,339.121) (57.910,339.121) (58.052,339.121) (58.194,339.121) (58.336,339.121) (58.479,339.121) (58.621,339.121) (58.763,339.121) (58.905,339.121) (59.047,339.121) (59.189,339.121) (59.331,339.121) (59.474,339.121) (59.616,339.121) (59.758,339.121) (59.900,339.121) (60.042,339.121) (60.184,339.121) (60.326,339.121) (60.469,339.121) (60.611,339.121) (60.753,339.121) (60.895,339.121) (61.037,339.121) (61.179,339.121) (61.321,339.121) (61.464,339.121) (61.606,339.121) (61.748,339.121) (61.890,339.121) (62.032,339.121) (62.174,339.121) (62.316,339.121) (62.459,339.121) (62.601,339.121) (62.743,339.121) (62.885,339.121) (63.027,339.121) (63.169,339.121) (63.311,339.121) (63.454,339.121) (63.596,339.121) (63.738,339.121) (63.880,339.121) (64.022,339.121) (64.164,339.121) (64.306,339.121) (64.449,339.121) (64.591,339.121) (64.733,339.121) (64.875,339.121) (65.017,339.121) (65.159,339.121) (65.301,339.121) (65.444,339.121) (65.586,339.121) (65.728,339.121) (65.870,339.121) (66.012,339.121) (66.154,339.121) (66.296,339.121) (66.439,339.121) (66.581,339.121) (66.723,339.121) (66.865,339.121) (67.007,339.121) (67.149,339.121) (67.291,339.121) (67.434,339.121) (67.576,339.121) (67.718,339.121) (67.860,339.121) (68.002,339.121) (68.144,339.121) (68.286,339.121) (68.429,339.121) (68.571,339.121) (68.713,339.121) (68.855,339.121) (68.997,339.121) (69.139,339.121) (69.281,339.121) (69.424,339.121) (69.566,339.121) (69.708,339.121) (69.850,339.121) (69.992,339.121) (70.134,339.121) (70.276,339.121) (70.419,339.121) (70.561,339.121) (70.703,339.121) (70.845,339.121) (70.987,339.121) (71.129,339.121) (71.271,339.121) (71.414,339.121) (71.556,339.121) (71.698,339.121) (71.840,339.121) (71.982,339.121) (72.124,339.121) (72.266,339.121) (72.409,339.121) (72.551,339.121) (72.693,339.121) (72.835,339.121) (72.977,339.121) (73.119,339.121) (73.261,339.121) (73.404,339.121) (73.546,339.121) (73.688,339.121) (73.830,339.121) (73.972,339.121) (74.114,339.121) (74.256,339.121) (74.399,339.121) (74.541,339.121) (74.683,339.121) (74.825,339.121) (74.967,339.121) (75.109,339.121) (75.251,339.121) (75.394,339.121) (75.536,339.121) (75.678,339.121) (75.820,339.121) (75.962,339.121) (76.104,339.121) (76.246,339.121) (76.389,339.121) (76.531,339.121) (76.673,339.121) (76.815,339.121) (76.957,339.121) (77.099,339.121) (77.241,339.121) (77.384,339.121) (77.526,339.121) (77.668,339.121) (77.810,339.121) (77.952,339.121) (78.094,339.121) (78.236,339.121) (78.379,339.121) (78.521,339.121) (78.663,339.121) (78.805,339.121) (78.947,339.121) (79.089,339.121) (79.231,339.121) (79.374,339.121) (79.516,339.121) (79.658,339.121) (79.800,339.121) (79.942,339.121) (80.084,339.121) (80.226,339.121) (80.369,339.121) (80.511,339.121) (80.653,339.121) (80.795,339.121) (80.937,339.121) (81.079,339.121) (81.221,339.121) (81.364,339.121) (81.506,339.121) (81.648,339.121) (81.790,339.121) (81.932,339.121) (82.074,339.121) (82.216,339.121) (82.359,339.121) (82.501,339.121) (82.643,339.121) (82.785,339.121) (82.927,339.121) (83.069,339.121) (83.211,339.121) (83.354,339.121) (83.496,339.121) (83.638,339.121) (83.780,339.121) (83.922,339.121) (84.064,339.121) (84.206,339.121) (84.349,339.121) (84.491,339.121) (84.633,339.121) (84.775,339.121) (84.917,339.121) (85.059,339.121) (85.201,339.121) (85.344,339.121) (85.486,339.121) (85.628,339.121) (85.770,339.121) (85.912,339.121) (86.054,339.121) (86.196,339.121) (86.339,339.121) (86.481,339.121) (86.623,339.121) (86.765,339.121) (86.907,339.121) (87.049,339.121) (87.191,339.121) (87.334,339.121) (87.476,339.121) (87.618,339.121) (87.760,339.121) (87.902,339.121) (88.044,339.121) (88.186,339.121) (88.329,339.121) (88.471,339.121) (88.613,339.121) (88.755,339.121) (88.897,339.121) (89.039,339.121) (89.181,339.121) (89.324,339.121) (89.466,339.121) (89.608,339.121) (89.750,339.121) (89.892,339.121) (90.034,339.121) (90.176,339.121) (90.319,339.121) (90.461,339.121) (90.603,339.121) (90.745,339.121) (90.887,339.121) (91.029,339.121) (91.171,339.121) (91.314,339.121) (91.456,339.121) (91.598,339.121) (91.740,339.121) (91.882,339.121) (92.024,339.121) (92.166,339.121) (92.309,339.121) (92.451,339.121) (92.593,339.121) (92.735,339.121) (92.877,339.121) (93.019,339.121) (93.161,339.121) (93.304,339.121) (93.446,339.121) (93.588,339.121) (93.730,339.121) (93.872,339.121) (94.014,339.121) (94.156,339.121) (94.299,339.121) (94.441,339.121) (94.583,339.121) (94.725,339.121) (94.867,339.121) (95.009,339.121) (95.151,339.121) (95.294,339.121) (95.436,339.121) (95.578,339.121) (95.720,339.121) (95.862,339.121) (96.004,339.121) (96.146,339.121) (96.289,339.121) (96.431,339.121) (96.573,339.121) (96.715,339.121) (96.857,343.796) (96.999,343.796) (97.141,343.796) (97.284,343.796) (97.426,343.796) (97.568,343.796) (97.710,343.796) (97.852,343.796) (97.994,343.796) (98.136,343.796) (98.279,343.796) (98.421,343.796) (98.563,343.796) (98.705,343.796) (98.847,343.796) (98.989,343.796) (99.131,343.796) (99.274,343.796) (99.416,343.796) (99.558,343.796) (99.700,343.796) (99.842,343.796) (99.984,343.796) (100.126,343.796) (100.269,343.796) (100.411,343.796) (100.553,343.796) (100.695,343.796) (100.837,343.796) (100.979,343.796) (101.121,343.796) (101.264,343.796) (101.406,343.796) (101.548,343.796) (101.690,343.796) (101.832,343.796) (101.974,343.796) (102.116,343.796) (102.259,343.796) (102.401,343.796) (102.543,343.796) (102.685,343.796) (102.827,343.796) (102.969,343.796) (103.111,343.796) (103.254,343.796) (103.396,343.796) (103.538,343.796) (103.680,343.796) (103.822,343.796) (103.964,343.796) (104.106,343.796) (104.249,343.796) (104.391,343.796) (104.533,343.796) (104.675,343.796) (104.817,343.796) (104.959,343.796) (105.101,343.796) (105.244,343.796) (105.386,343.796) (105.528,343.796) (105.670,343.796) (105.812,343.796) (105.954,343.796) (106.096,343.796) (106.239,343.796) (106.381,343.796) (106.523,343.796) (106.665,343.796) (106.807,343.796) (106.949,343.796) (107.091,343.796) (107.234,343.796) (107.376,343.796) (107.518,343.796) (107.660,343.796) (107.802,343.796) (107.944,343.796) (108.086,343.796) (108.229,343.796) (108.371,343.796) (108.513,343.796) (108.655,343.796) (108.797,343.796) (108.939,343.796) (109.081,343.796) (109.224,343.796) (109.366,343.796) (109.508,343.796) (109.650,343.796) (109.792,343.796) (109.934,343.796) (110.076,343.796) (110.219,343.796) (110.361,343.796) (110.503,343.796) (110.645,343.796) (110.787,343.796) (110.929,343.796) (111.071,343.796) (111.214,343.796) (111.356,343.796) (111.498,343.796) (111.640,343.796) (111.782,343.796) (111.924,343.796) (112.066,343.796) (112.209,343.796) (112.351,343.796) (112.493,343.796) (112.635,343.796) (112.777,343.796) (112.919,343.796) (113.061,343.796) (113.204,343.796) (113.346,343.796) (113.488,343.796) (113.630,343.796) (113.772,343.796) (113.914,343.796) (114.056,343.796) (114.199,343.796) (114.341,343.796) (114.483,343.796) (114.625,343.796) (114.767,343.796) (114.909,343.796) (115.051,343.796) (115.194,343.796) (115.336,343.796) (115.478,343.796) (115.620,343.796) (115.762,343.796) (115.904,343.796) (116.046,343.796) (116.189,343.796) (116.331,343.796) (116.473,343.796) (116.615,343.796) (116.757,343.796) (116.899,343.796) (117.041,343.796) (117.184,343.796) (117.326,343.796) (117.468,343.796) (117.610,343.796) (117.752,343.796) (117.894,343.796) (118.036,343.796) (118.179,343.796) (118.321,343.796) (118.463,343.796) (118.605,343.796) (118.747,343.796) (118.889,343.796) (119.031,343.796) (119.174,343.796) (119.316,343.796) (119.458,343.796) (119.600,343.796) (119.742,343.796) (119.884,343.796) (120.026,343.796) (120.169,343.796) (120.311,343.796) (120.453,343.796) (120.595,343.796) (120.737,343.796) (120.879,343.796) (121.021,353.276) (121.164,353.276) (121.306,353.276) (121.448,353.276) (121.590,353.276) (121.732,353.276) (121.874,353.276) (122.016,353.276) (122.159,353.276) (122.301,353.276) (122.443,353.276) (122.585,353.276) (122.727,353.276) (122.869,353.276) (123.011,353.276) (123.154,353.276) (123.296,353.276) (123.438,353.276) (123.580,353.276) (123.722,353.276) (123.864,353.276) (124.006,353.276) (124.149,353.276) (124.291,353.276) (124.433,353.276) (124.575,353.276) (124.717,353.276) (124.859,353.276) (125.001,353.276) (125.144,353.276) (125.286,353.276) (125.428,353.276) (125.570,353.276) (125.712,353.276) (125.854,353.276) (125.996,353.276) (126.139,353.276) (126.281,353.276) (126.423,353.276) (126.565,353.276) (126.707,353.276) (126.849,353.276) (126.991,353.276) (127.134,353.276) (127.276,353.276) (127.418,353.276) (127.560,353.276) (127.702,353.276) (127.844,353.276) (127.986,353.276) (128.129,353.276) (128.271,353.276) (128.413,353.276) (128.555,353.276) (128.697,353.276) (128.839,353.276) (128.981,353.276) (129.124,353.276) (129.266,353.276) (129.408,353.276) (129.550,353.276) (129.692,353.276) (129.834,353.276) (129.976,353.276) (130.119,353.276) (130.261,353.276) (130.403,353.276) (130.545,353.276) (130.687,353.276) (130.829,353.276) (130.971,353.276) (131.114,353.276) (131.256,353.276) (131.398,353.276) (131.540,353.276) (131.682,353.276) (131.824,353.276) (131.966,353.276) (132.109,353.276) (132.251,353.276) (132.393,353.276) (132.535,353.276) (132.677,353.276) (132.819,353.276) (132.961,353.276) (133.104,353.276) (133.246,353.276) (133.388,353.276) (133.530,353.276) (133.672,353.276) (133.814,353.276) (133.956,353.276) (134.099,353.276) (134.241,353.276) (134.383,353.276) (134.525,353.276) (134.667,353.276) (134.809,353.276) (134.951,353.276) (135.094,353.276) (135.236,353.276) (135.378,353.276) (135.520,353.276) (135.662,353.276) (135.804,353.276) (135.946,353.276) (136.089,353.276) (136.231,353.276) (136.373,353.276) (136.515,353.276) (136.657,353.276) (136.799,353.276) (136.941,353.276) (137.084,353.276) (137.226,353.276) (137.368,353.276) (137.510,353.276) (137.652,353.276) (137.794,353.276) (137.936,353.276) (138.079,353.276) (138.221,353.276) (138.363,353.276) (138.505,353.276) (138.647,353.276) (138.789,353.276) (138.931,353.276) (139.074,353.276) (139.216,353.276) (139.358,353.276) (139.500,353.276) (139.642,353.276) (139.784,353.276) (139.926,353.276) (140.069,353.276) (140.211,353.276) (140.353,353.276) (140.495,353.276) (140.637,353.276) (140.779,353.276) (140.921,353.276) (141.064,353.276) (141.206,353.276) (141.348,353.276) (141.490,353.276) (141.632,353.276) (141.774,353.276) (141.916,353.276) (142.059,353.276) (142.201,353.276) (142.343,353.276) (142.485,353.276) (142.627,353.276) (142.769,353.276) (142.911,353.276) (143.054,353.276) (143.196,353.276) (143.338,353.276) (143.480,353.276) (143.622,353.276) (143.764,353.276) (143.906,353.276) (144.049,353.276) (144.191,353.276) (144.333,353.276) (144.475,353.276) (144.617,353.276) (144.759,353.276) (144.901,353.276) (145.044,353.276) (145.186,353.276) (145.328,353.276) (145.470,353.276) (145.612,353.276) (145.754,353.276) (145.896,353.276) (146.039,353.276) (146.181,353.276) (146.323,353.276) (146.465,353.276) (146.607,353.276) (146.749,353.276) (146.891,353.276) (147.034,353.276) (147.176,353.276) (147.318,353.276) (147.460,353.276) (147.602,353.276) (147.744,353.276) (147.886,353.276) (148.029,353.276) (148.171,353.276) (148.313,353.276) (148.455,353.276) (148.597,353.276) (148.739,353.276) (148.881,353.276) (149.024,353.276) (149.166,353.276) (149.308,353.276) (149.450,353.276) (149.592,353.276) (149.734,353.276) (149.876,353.276) (150.019,353.276) (150.161,353.276) (150.303,353.276) (150.445,353.276) (150.587,353.276) (150.729,353.276) (150.871,353.276) (151.014,353.276) (151.156,353.276) (151.298,353.276) (151.440,353.276) (151.582,353.276) (151.724,353.276) (151.866,353.276) (152.009,353.276) (152.151,353.276) (152.293,353.276) (152.435,353.276) (152.577,353.276) (152.719,353.276) (152.861,353.276) (153.004,353.276) (153.146,353.276) (153.288,353.276) (153.430,353.276) (153.572,353.276) (153.714,353.276) (153.856,353.276) (153.999,353.276) (154.141,353.276) (154.283,353.276) (154.425,353.276) (154.567,353.276) (154.709,353.276) (154.851,353.276) (154.994,353.276) (155.136,353.276) (155.278,353.276) (155.420,353.276) (155.562,353.276) (155.704,353.276) (155.846,353.276) (155.989,353.276) (156.131,353.276) (156.273,353.276) (156.415,353.276) (156.557,353.276) (156.699,353.276) (156.841,353.276) (156.984,353.276) (157.126,353.276) (157.268,353.276) (157.410,353.276) (157.552,353.276) (157.694,353.276) (157.836,353.276) (157.979,353.276) (158.121,353.276) (158.263,353.276) (158.405,353.276) (158.547,353.276) (158.689,353.276) (158.831,353.276) (158.974,353.276) (159.116,353.276) (159.258,353.276) (159.400,353.276) (159.542,353.276) (159.684,353.276) (159.826,353.276) (159.969,353.276) (160.111,353.276) (160.253,353.276) (160.395,367.522) (160.537,367.522) (160.679,367.522) (160.821,367.522) (160.964,367.522) (161.106,367.522) (161.248,367.522) (161.390,367.522) (161.532,367.522) (161.674,367.522) (161.816,367.522) (161.959,367.522) (162.101,367.522) (162.243,367.522) (162.385,367.522) (162.527,367.522) (162.669,367.522) (162.811,367.522) (162.954,367.522) (163.096,367.522) (163.238,367.522) (163.380,367.522) (163.522,367.522) (163.664,367.522) (163.806,367.522) (163.949,367.522) (164.091,367.522) (164.233,367.522) (164.375,367.522) (164.517,367.522) (164.659,367.522) (164.801,367.522) (164.944,367.522) (165.086,367.522) (165.228,367.522) (165.370,367.522) (165.512,367.522) (165.654,367.522) (165.796,367.522) (165.939,367.522) (166.081,367.522) (166.223,367.522) (166.365,367.522) (166.507,367.522) (166.649,367.522) (166.791,367.522) (166.934,367.522) (167.076,367.522) (167.218,367.522) (167.360,367.522) (167.502,367.522) (167.644,367.522) (167.786,367.522) (167.929,384.284) (168.071,384.284) (168.213,384.284) (168.355,384.284) (168.497,384.284) (168.639,384.284) (168.781,384.284) (168.924,384.284) (169.066,384.284) (169.208,384.284) (169.350,384.284) (169.492,384.284) (169.634,384.284) (169.776,384.284) (169.919,384.284) (170.061,384.284) (170.203,384.284) (170.345,384.284) (170.487,384.284) (170.629,384.284) (170.771,384.284) (170.914,384.284) (171.056,384.284) (171.198,384.284) (171.340,384.284) (171.482,384.284) (171.624,384.284) (171.766,384.284) (171.909,384.284) (172.051,384.284) (172.193,384.284) (172.335,384.284) (172.477,384.284) (172.619,384.284) (172.761,384.284) (172.904,384.284) (173.046,384.284) (173.188,384.284) (173.330,384.284) (173.472,384.284) (173.614,384.284) (173.756,384.284) (173.899,384.284) (174.041,384.284) (174.183,384.284) (174.325,384.284) (174.467,384.284) (174.609,384.284) (174.751,384.284) (174.894,384.284) (175.036,384.284) (175.178,384.284) (175.320,384.284) (175.462,384.284) (175.604,384.284) (175.746,384.284) (175.889,384.284) (176.031,384.284) (176.173,384.284) (176.315,384.284) (176.457,384.284) (176.599,384.284) (176.741,384.284) (176.884,384.284) (177.026,384.284) (177.168,384.284) (177.310,384.284) (177.452,384.284) (177.594,384.284) (177.736,384.284) (177.879,384.284) (178.021,384.284) (178.163,384.284) (178.305,384.284) (178.447,384.284) (178.589,384.284) (178.731,384.284) (178.874,384.284) (179.016,384.284) (179.158,384.284) (179.300,384.284) (179.442,384.284) (179.584,384.284) (179.726,384.284) (179.869,384.284) (180.011,384.284) (180.153,384.284) (180.295,384.284) (180.437,384.284) (180.579,384.284) (180.721,384.284) (180.864,384.284) (181.006,384.284) (181.148,384.284) (181.290,384.284) (181.432,384.284) (181.574,384.284) (181.716,384.284) (181.859,384.284) (182.001,384.284) (182.143,384.284) (182.285,384.284) (182.427,384.284) (182.569,384.284) (182.711,384.284) (182.854,384.284) (182.996,384.284) (183.138,384.284) (183.280,384.284) (183.422,384.284) (183.564,384.284) (183.706,384.284) (183.849,384.284) (183.991,384.284) (184.133,384.284) (184.275,384.284) (184.417,384.284) (184.559,384.284) (184.701,384.284) (184.844,384.284) (184.986,384.284) (185.128,384.284) (185.270,384.284) (185.412,384.284) (185.554,384.284) (185.696,384.284) (185.839,384.284) (185.981,384.284) (186.123,384.284) (186.265,384.284) (186.407,384.284) (186.549,384.284) (186.691,384.284) (186.834,384.284) (186.976,384.284) (187.118,384.284) (187.260,384.284) (187.402,384.284) (187.544,384.284) (187.686,384.284) (187.829,384.284) (187.971,384.284) (188.113,384.284) (188.255,384.284) (188.397,384.284) (188.539,384.284) (188.681,384.284) (188.824,384.284) (188.966,384.284) (189.108,384.284) (189.250,384.284) (189.392,384.284) (189.534,384.284) (189.676,384.284) (189.819,384.284) (189.961,384.284) (190.103,384.284) (190.245,384.284) (190.387,384.284) (190.529,384.284) (190.671,384.284) (190.814,384.284) (190.956,384.284) (191.098,384.284) (191.240,384.284) (191.382,384.284) (191.524,384.284) (191.666,384.284) (191.809,384.284) (191.951,384.284) (192.093,384.284) (192.235,384.284) (192.377,384.284) (192.519,384.284) (192.661,384.284) (192.804,384.284) (192.946,384.284) (193.088,384.284) (193.230,384.284) (193.372,384.284) (193.514,384.284) (193.656,384.284) (193.799,384.284) (193.941,384.284) (194.083,384.284) (194.225,384.284) (194.367,384.284) (194.509,384.284) (194.651,384.284) (194.794,384.284) (194.936,384.284) (195.078,384.284) (195.220,384.284) (195.362,384.284) (195.504,384.284) (195.646,384.284) (195.789,384.284) (195.931,384.284) (196.073,384.284) (196.215,384.284) (196.357,384.284) (196.499,384.284) (196.641,384.284) (196.784,384.284) (196.926,384.284) (197.068,384.284) (197.210,384.284) (197.352,384.284) (197.494,384.284) (197.636,384.284) (197.779,384.284) (197.921,384.284) (198.063,384.284) (198.205,384.284) (198.347,384.284) (198.489,384.284) (198.631,384.284) (198.774,384.284) (198.916,384.284) (199.058,384.284) (199.200,384.284) (199.342,384.284) (199.484,384.284) (199.626,384.284) (199.769,384.284) (199.911,384.284) (200.053,384.284) (200.195,384.284) (200.337,384.284) (200.479,384.284) (200.621,384.284) (200.764,384.284) (200.906,384.284) (201.048,384.284) (201.190,384.284) (201.332,384.284) (201.474,384.284) (201.616,384.284) (201.759,384.284) (201.901,384.284) (202.043,384.284) (202.185,384.284) (202.327,384.284) (202.469,384.284) (202.611,384.284) (202.754,384.284) (202.896,384.284) (203.038,384.284) (203.180,384.284) (203.322,384.284) (203.464,384.284) (203.606,384.284) (203.749,384.284) (203.891,384.284) (204.033,384.284) (204.175,384.284) (204.317,384.284) (204.459,384.284) (204.601,384.284) (204.744,384.284) (204.886,384.284) (205.028,384.284) (205.170,384.284) (205.312,384.284) (205.454,384.284) (205.596,384.284) (205.739,384.284) (205.881,384.284) (206.023,384.284) (206.165,384.284) (206.307,384.284) (206.449,384.284) (206.591,384.284) (206.734,384.284) (206.876,384.284) (207.018,384.284) (207.160,384.284) (207.302,384.284) (207.444,384.284) (207.586,384.284) (207.729,384.284) (207.871,384.284) (208.013,384.284) (208.155,408.453) (208.297,408.453) (208.439,408.453) (208.581,408.453) (208.724,408.453) (208.866,408.453) (209.008,408.453) (209.150,408.453) (209.292,408.453) (209.434,408.453) (209.576,408.453) (209.719,408.453) (209.861,408.453) (210.003,408.453) (210.145,408.453) (210.287,408.453) (210.429,408.453) (210.571,408.453) (210.714,408.453) (210.856,408.453) (210.998,408.453) (211.140,408.453) (211.282,408.453) (211.424,408.453) (211.566,408.453) (211.709,408.453) (211.851,408.453) (211.993,408.453) (212.135,408.453) (212.277,408.453) (212.419,408.453) (212.561,408.453) (212.704,408.453) (212.846,408.453) (212.988,408.453) (213.130,408.453) (213.272,408.453) (213.414,408.453) (213.556,408.453) (213.699,408.453) (213.841,408.453) (213.983,408.453) (214.125,408.453) (214.267,408.453) (214.409,408.453) (214.551,408.453) (214.694,408.453) (214.836,408.453) (214.978,408.453) (215.120,408.453) (215.262,408.453) (215.404,408.453) (215.546,408.453) (215.689,408.453) (215.831,408.453) (215.973,408.453) (216.115,408.453) (216.257,408.453) (216.399,408.453) (216.541,408.453) (216.684,408.453) (216.826,408.453) (216.968,408.453) (217.110,408.453) (217.252,408.453) (217.394,408.453) (217.536,408.453) (217.679,408.453) (217.821,408.453) (217.963,408.453) (218.105,408.453) (218.247,408.453) (218.389,408.453) (218.531,408.453) (218.674,408.453) (218.816,408.453) (218.958,408.453) (219.100,408.453) (219.242,408.453) (219.384,408.453) (219.526,408.453) (219.669,408.453) (219.811,408.453) (219.953,408.453) (220.095,408.453) (220.237,408.453) (220.379,408.453) (220.521,408.453) (220.664,408.453) (220.806,408.453) (220.948,408.453) (221.090,408.453) (221.232,408.453) (221.374,408.453) (221.516,408.453) (221.659,408.453) (221.801,408.453) (221.943,408.453) (222.085,408.453) (222.227,408.453) (222.369,408.453) (222.511,408.453) (222.654,408.453) (222.796,408.453) (222.938,408.453) (223.080,408.453) (223.222,408.453) (223.364,408.453) (223.506,408.453) (223.649,408.453) (223.791,408.453) (223.933,408.453) (224.075,408.453) (224.217,408.453) (224.359,408.453) (224.501,408.453) (224.644,408.453) (224.786,408.453) (224.928,408.453) (225.070,408.453) (225.212,408.453) (225.354,408.453) (225.496,408.453) (225.639,408.453) (225.781,408.453) (225.923,408.453) (226.065,408.453) (226.207,408.453) (226.349,408.453) (226.491,408.453) (226.634,408.453) (226.776,408.453) (226.918,408.453) (227.060,408.453) (227.202,408.453) (227.344,408.453) (227.486,408.453) (227.629,408.453) (227.771,408.453) (227.913,408.453) (228.055,408.453) (228.197,408.453) (228.339,408.453) (228.481,408.453) (228.624,408.453) (228.766,408.453) (228.908,408.453) (229.050,408.453) (229.192,408.453) (229.334,408.453) (229.476,408.453) (229.619,408.453) (229.761,408.453) (229.903,408.453) (230.045,408.453) (230.187,408.453) (230.329,408.453) (230.471,408.453) (230.614,408.453) (230.756,408.453) (230.898,408.453) (231.040,408.453) (231.182,408.453) (231.324,408.453) (231.466,408.453) (231.609,408.453) (231.751,408.453) (231.893,408.453) (232.035,408.453) (232.177,408.453) (232.319,408.453) (232.461,408.453) (232.604,408.453) (232.746,408.453) (232.888,408.453) (233.030,408.453) (233.172,408.453) (233.314,408.453) (233.456,408.453) (233.599,408.453) (233.741,408.453) (233.883,408.453) (234.025,408.453) (234.167,408.453) (234.309,408.453) (234.451,408.453) (234.594,408.453) (234.736,408.453) (234.878,408.453) (235.020,408.453) (235.162,408.453) (235.304,408.453) (235.446,408.453) (235.589,408.453) (235.731,408.453) (235.873,408.453) (236.015,408.453) (236.157,408.453) (236.299,408.453) (236.441,408.453) (236.584,408.453) (236.726,408.453) (236.868,408.453) (237.010,408.453) (237.152,408.453) (237.294,408.453) (237.436,408.453) (237.579,408.453) (237.721,408.453) (237.863,408.453) (238.005,408.453) (238.147,408.453) (238.289,408.453) (238.431,408.453) (238.574,408.453) (238.716,408.453) (238.858,408.453) (239.000,408.453)};
\draw[methodA,line width=1.05pt,dashed] plot coordinates {(40.000,342.682) (40.142,342.635) (40.284,342.588) (40.426,342.541) (40.569,342.495) (40.711,342.449) (40.853,342.403) (40.995,342.357) (41.137,342.312) (41.279,342.267) (41.421,342.222) (41.564,342.177) (41.706,342.133) (41.848,342.088) (41.990,342.044) (42.132,342.001) (42.274,341.957) (42.416,341.914) (42.559,341.871) (42.701,341.828) (42.843,341.785) (42.985,341.743) (43.127,341.700) (43.269,341.658) (43.411,341.617) (43.554,341.575) (43.696,341.534) (43.838,341.493) (43.980,341.452) (44.122,341.411) (44.264,341.371) (44.406,341.331) (44.549,341.291) (44.691,341.251) (44.833,341.212) (44.975,341.173) (45.117,341.134) (45.259,341.095) (45.401,341.056) (45.544,341.018) (45.686,340.980) (45.828,340.942) (45.970,340.904) (46.112,340.867) (46.254,340.830) (46.396,340.793) (46.539,340.756) (46.681,340.719) (46.823,340.683) (46.965,340.647) (47.107,340.611) (47.249,340.575) (47.391,340.540) (47.534,340.505) (47.676,340.470) (47.818,340.435) (47.960,340.400) (48.102,340.366) (48.244,340.332) (48.386,340.298) (48.529,340.265) (48.671,340.231) (48.813,340.198) (48.955,340.165) (49.097,340.132) (49.239,340.100) (49.381,340.067) (49.524,340.035) (49.666,340.003) (49.808,339.972) (49.950,339.940) (50.092,339.909) (50.234,339.878) (50.376,339.847) (50.519,339.817) (50.661,339.786) (50.803,339.756) (50.945,339.726) (51.087,339.696) (51.229,339.667) (51.371,339.638) (51.514,339.609) (51.656,339.580) (51.798,339.551) (51.940,339.523) (52.082,339.495) (52.224,339.467) (52.366,339.439) (52.509,339.411) (52.651,339.384) (52.793,339.357) (52.935,339.330) (53.077,339.304) (53.219,339.277) (53.361,339.251) (53.504,339.225) (53.646,339.199) (53.788,339.173) (53.930,339.148) (54.072,339.123) (54.214,339.098) (54.356,339.073) (54.499,339.049) (54.641,339.024) (54.783,339.000) (54.925,338.976) (55.067,338.953) (55.209,338.929) (55.351,338.906) (55.494,338.883) (55.636,338.860) (55.778,338.837) (55.920,338.815) (56.062,338.793) (56.204,338.771) (56.346,338.749) (56.489,338.728) (56.631,338.706) (56.773,338.685) (56.915,338.664) (57.057,338.643) (57.199,338.623) (57.341,338.603) (57.484,338.583) (57.626,338.563) (57.768,338.543) (57.910,338.524) (58.052,338.504) (58.194,338.485) (58.336,338.466) (58.479,338.448) (58.621,338.429) (58.763,338.411) (58.905,338.393) (59.047,338.375) (59.189,338.358) (59.331,338.340) (59.474,338.323) (59.616,338.306) (59.758,338.289) (59.900,338.273) (60.042,338.256) (60.184,338.240) (60.326,338.224) (60.469,338.209) (60.611,338.193) (60.753,338.178) (60.895,338.163) (61.037,338.148) (61.179,338.133) (61.321,338.118) (61.464,338.104) (61.606,338.090) (61.748,338.076) (61.890,338.062) (62.032,338.049) (62.174,338.036) (62.316,338.022) (62.459,338.010) (62.601,337.997) (62.743,337.984) (62.885,337.972) (63.027,337.960) (63.169,337.948) (63.311,337.936) (63.454,337.925) (63.596,337.914) (63.738,337.903) (63.880,337.892) (64.022,337.881) (64.164,337.870) (64.306,337.860) (64.449,337.850) (64.591,337.840) (64.733,337.831) (64.875,337.821) (65.017,337.812) (65.159,337.803) (65.301,337.794) (65.444,337.785) (65.586,337.777) (65.728,337.768) (65.870,337.760) (66.012,337.752) (66.154,337.745) (66.296,337.737) (66.439,337.730) (66.581,337.723) (66.723,337.716) (66.865,337.709) (67.007,337.702) (67.149,337.696) (67.291,337.690) (67.434,337.684) (67.576,337.678) (67.718,337.673) (67.860,337.667) (68.002,337.662) (68.144,337.657) (68.286,337.653) (68.429,337.648) (68.571,337.644) (68.713,337.639) (68.855,337.635) (68.997,337.632) (69.139,337.628) (69.281,337.625) (69.424,337.621) (69.566,337.618) (69.708,337.615) (69.850,337.613) (69.992,337.610) (70.134,337.608) (70.276,337.606) (70.419,337.604) (70.561,337.602) (70.703,337.601) (70.845,337.600) (70.987,337.599) (71.129,337.598) (71.271,337.597) (71.414,337.596) (71.556,337.596) (71.698,337.596) (71.840,337.596) (71.982,337.596) (72.124,337.596) (72.266,337.597) (72.409,337.598) (72.551,337.599) (72.693,337.600) (72.835,337.601) (72.977,337.603) (73.119,337.605) (73.261,337.606) (73.404,337.609) (73.546,337.611) (73.688,337.613) (73.830,337.616) (73.972,337.619) (74.114,337.622) (74.256,337.625) (74.399,337.628) (74.541,337.632) (74.683,337.636) (74.825,337.640) (74.967,337.644) (75.109,337.648) (75.251,337.653) (75.394,337.657) (75.536,337.662) (75.678,337.667) (75.820,337.673) (75.962,337.678) (76.104,337.684) (76.246,337.689) (76.389,337.695) (76.531,337.702) (76.673,337.708) (76.815,337.715) (76.957,337.721) (77.099,337.728) (77.241,337.735) (77.384,337.742) (77.526,337.750) (77.668,337.758) (77.810,337.765) (77.952,337.773) (78.094,337.781) (78.236,337.790) (78.379,337.798) (78.521,337.807) (78.663,337.816) (78.805,337.825) (78.947,337.834) (79.089,337.844) (79.231,337.853) (79.374,337.863) (79.516,337.873) (79.658,337.883) (79.800,337.893) (79.942,337.904) (80.084,337.915) (80.226,337.925) (80.369,337.937) (80.511,337.948) (80.653,337.959) (80.795,337.971) (80.937,337.982) (81.079,337.994) (81.221,338.006) (81.364,338.019) (81.506,338.031) (81.648,338.044) (81.790,338.057) (81.932,338.070) (82.074,338.083) (82.216,338.096) (82.359,338.110) (82.501,338.123) (82.643,338.137) (82.785,338.151) (82.927,338.165) (83.069,338.180) (83.211,338.194) (83.354,338.209) (83.496,338.224) (83.638,338.239) (83.780,338.254) (83.922,338.269) (84.064,338.285) (84.206,338.301) (84.349,338.317) (84.491,338.333) (84.633,338.349) (84.775,338.366) (84.917,338.382) (85.059,338.399) (85.201,338.416) (85.344,338.433) (85.486,338.450) (85.628,338.468) (85.770,338.485) (85.912,338.503) (86.054,338.521) (86.196,338.539) (86.339,338.558) (86.481,338.576) (86.623,338.595) (86.765,338.614) (86.907,338.633) (87.049,338.652) (87.191,338.671) (87.334,338.691) (87.476,338.710) (87.618,338.730) (87.760,338.750) (87.902,338.770) (88.044,338.791) (88.186,338.811) (88.329,338.832) (88.471,338.853) (88.613,338.874) (88.755,338.895) (88.897,338.916) (89.039,338.938) (89.181,338.959) (89.324,338.981) (89.466,339.003) (89.608,339.025) (89.750,339.048) (89.892,339.070) (90.034,339.093) (90.176,339.116) (90.319,339.139) (90.461,339.162) (90.603,339.185) (90.745,339.209) (90.887,339.232) (91.029,339.256) (91.171,339.280) (91.314,339.304) (91.456,339.328) (91.598,339.353) (91.740,339.378) (91.882,339.402) (92.024,339.427) (92.166,339.452) (92.309,339.478) (92.451,339.503) (92.593,339.529) (92.735,339.554) (92.877,339.580) (93.019,339.606) (93.161,339.633) (93.304,339.659) (93.446,339.686) (93.588,339.712) (93.730,339.739) (93.872,339.766) (94.014,339.793) (94.156,339.821) (94.299,339.848) (94.441,339.876) (94.583,339.904) (94.725,339.932) (94.867,339.960) (95.009,339.988) (95.151,340.016) (95.294,340.045) (95.436,340.074) (95.578,340.103) (95.720,340.132) (95.862,340.161) (96.004,340.190) (96.146,340.220) (96.289,340.250) (96.431,340.280) (96.573,340.310) (96.715,340.340) (96.857,340.370) (96.999,340.401) (97.141,340.431) (97.284,340.462) (97.426,340.493) (97.568,340.524) (97.710,340.555) (97.852,340.587) (97.994,340.618) (98.136,340.650) (98.279,340.682) (98.421,340.714) (98.563,340.746) (98.705,340.778) (98.847,340.811) (98.989,340.843) (99.131,340.876) (99.274,340.909) (99.416,340.942) (99.558,340.975) (99.700,341.009) (99.842,341.042) (99.984,341.076) (100.126,341.110) (100.269,341.144) (100.411,341.178) (100.553,341.212) (100.695,341.247) (100.837,341.281) (100.979,341.316) (101.121,341.351) (101.264,341.386) (101.406,341.421) (101.548,341.456) (101.690,341.492) (101.832,341.528) (101.974,341.563) (102.116,341.599) (102.259,341.635) (102.401,341.672) (102.543,341.708) (102.685,341.744) (102.827,341.781) (102.969,341.818) (103.111,341.855) (103.254,341.892) (103.396,341.929) (103.538,341.967) (103.680,342.004) (103.822,342.042) (103.964,342.080) (104.106,342.118) (104.249,342.156) (104.391,342.194) (104.533,342.232) (104.675,342.271) (104.817,342.310) (104.959,342.349) (105.101,342.388) (105.244,342.427) (105.386,342.466) (105.528,342.505) (105.670,342.545) (105.812,342.585) (105.954,342.624) (106.096,342.664) (106.239,342.705) (106.381,342.745) (106.523,342.785) (106.665,342.826) (106.807,342.866) (106.949,342.907) (107.091,342.948) (107.234,342.989) (107.376,343.031) (107.518,343.072) (107.660,343.114) (107.802,343.155) (107.944,343.197) (108.086,343.239) (108.229,343.281) (108.371,343.323) (108.513,343.366) (108.655,343.408) (108.797,343.451) (108.939,343.494) (109.081,343.537) (109.224,343.580) (109.366,343.623) (109.508,343.666) (109.650,343.710) (109.792,343.753) (109.934,343.797) (110.076,343.841) (110.219,343.885) (110.361,343.929) (110.503,343.974) (110.645,344.018) (110.787,344.063) (110.929,344.107) (111.071,344.152) (111.214,344.197) (111.356,344.242) (111.498,344.287) (111.640,344.333) (111.782,344.378) (111.924,344.424) (112.066,344.470) (112.209,344.516) (112.351,344.562) (112.493,344.608) (112.635,344.654) (112.777,344.701) (112.919,344.747) (113.061,344.794) (113.204,344.841) (113.346,344.888) (113.488,344.935) (113.630,344.982) (113.772,345.030) (113.914,345.077) (114.056,345.125) (114.199,345.173) (114.341,345.221) (114.483,345.269) (114.625,345.317) (114.767,345.365) (114.909,345.414) (115.051,345.462) (115.194,345.511) (115.336,345.560) (115.478,345.609) (115.620,345.658) (115.762,345.707) (115.904,345.756) (116.046,345.806) (116.189,345.855) (116.331,345.905) (116.473,345.955) (116.615,346.005) (116.757,346.055) (116.899,346.105) (117.041,346.156) (117.184,346.206) (117.326,346.257) (117.468,346.307) (117.610,346.358) (117.752,346.409) (117.894,346.460) (118.036,346.512) (118.179,346.563) (118.321,346.615) (118.463,346.666) (118.605,346.718) (118.747,346.770) (118.889,346.822) (119.031,346.874) (119.174,346.926) (119.316,346.979) (119.458,347.031) (119.600,347.084) (119.742,347.136) (119.884,347.189) (120.026,347.242) (120.169,347.295) (120.311,347.349) (120.453,347.402) (120.595,347.456) (120.737,347.509) (120.879,347.563) (121.021,347.617) (121.164,347.671) (121.306,347.725) (121.448,347.779) (121.590,347.833) (121.732,347.888) (121.874,347.942) (122.016,347.997) (122.159,348.052) (122.301,348.107) (122.443,348.162) (122.585,348.217) (122.727,348.272) (122.869,348.328) (123.011,348.383) (123.154,348.439) (123.296,348.495) (123.438,348.551) (123.580,348.607) (123.722,348.663) (123.864,348.719) (124.006,348.775) (124.149,348.832) (124.291,348.889) (124.433,348.945) (124.575,349.002) (124.717,349.059) (124.859,349.116) (125.001,349.173) (125.144,349.231) (125.286,349.288) (125.428,349.346) (125.570,349.403) (125.712,349.461) (125.854,349.519) (125.996,349.577) (126.139,349.635) (126.281,349.693) (126.423,349.752) (126.565,349.810) (126.707,349.869) (126.849,349.927) (126.991,349.986) (127.134,350.045) (127.276,350.104) (127.418,350.163) (127.560,350.222) (127.702,350.282) (127.844,350.341) (127.986,350.401) (128.129,350.461) (128.271,350.520) (128.413,350.580) (128.555,350.640) (128.697,350.701) (128.839,350.761) (128.981,350.821) (129.124,350.882) (129.266,350.942) (129.408,351.003) (129.550,351.064) (129.692,351.125) (129.834,351.186) (129.976,351.247) (130.119,351.308) (130.261,351.369) (130.403,351.431) (130.545,351.492) (130.687,351.554) (130.829,351.616) (130.971,351.678) (131.114,351.740) (131.256,351.802) (131.398,351.864) (131.540,351.927) (131.682,351.989) (131.824,352.052) (131.966,352.114) (132.109,352.177) (132.251,352.240) (132.393,352.303) (132.535,352.366) (132.677,352.429) (132.819,352.492) (132.961,352.556) (133.104,352.619) (133.246,352.683) (133.388,352.747) (133.530,352.810) (133.672,352.874) (133.814,352.938) (133.956,353.003) (134.099,353.067) (134.241,353.131) (134.383,353.196) (134.525,353.260) (134.667,353.325) (134.809,353.390) (134.951,353.454) (135.094,353.519) (135.236,353.584) (135.378,353.650) (135.520,353.715) (135.662,353.780) (135.804,353.846) (135.946,353.911) (136.089,353.977) (136.231,354.043) (136.373,354.109) (136.515,354.175) (136.657,354.241) (136.799,354.307) (136.941,354.373) (137.084,354.439) (137.226,354.506) (137.368,354.573) (137.510,354.639) (137.652,354.706) (137.794,354.773) (137.936,354.840) (138.079,354.907) (138.221,354.974) (138.363,355.041) (138.505,355.109) (138.647,355.176) (138.789,355.244) (138.931,355.311) (139.074,355.379) (139.216,355.447) (139.358,355.515) (139.500,355.583) (139.642,355.651) (139.784,355.719) (139.926,355.788) (140.069,355.856) (140.211,355.925) (140.353,355.993) (140.495,356.062) (140.637,356.131) (140.779,356.200) (140.921,356.269) (141.064,356.338) (141.206,356.407) (141.348,356.476) (141.490,356.546) (141.632,356.615) (141.774,356.685) (141.916,356.754) (142.059,356.824) (142.201,356.894) (142.343,356.964) (142.485,357.034) (142.627,357.104) (142.769,357.174) (142.911,357.245) (143.054,357.315) (143.196,357.385) (143.338,357.456) (143.480,357.527) (143.622,357.597) (143.764,357.668) (143.906,357.739) (144.049,357.810) (144.191,357.881) (144.333,357.953) (144.475,358.024) (144.617,358.095) (144.759,358.167) (144.901,358.238) (145.044,358.310) (145.186,358.382) (145.328,358.454) (145.470,358.525) (145.612,358.597) (145.754,358.670) (145.896,358.742) (146.039,358.814) (146.181,358.886) (146.323,358.959) (146.465,359.031) (146.607,359.104) (146.749,359.177) (146.891,359.250) (147.034,359.322) (147.176,359.395) (147.318,359.468) (147.460,359.542) (147.602,359.615) (147.744,359.688) (147.886,359.762) (148.029,359.835) (148.171,359.909) (148.313,359.982) (148.455,360.056) (148.597,360.130) (148.739,360.204) (148.881,360.278) (149.024,360.352) (149.166,360.426) (149.308,360.500) (149.450,360.574) (149.592,360.649) (149.734,360.723) (149.876,360.798) (150.019,360.873) (150.161,360.947) (150.303,361.022) (150.445,361.097) (150.587,361.172) (150.729,361.247) (150.871,361.322) (151.014,361.397) (151.156,361.473) (151.298,361.548) (151.440,361.624) (151.582,361.699) (151.724,361.775) (151.866,361.850) (152.009,361.926) (152.151,362.002) (152.293,362.078) (152.435,362.154) (152.577,362.230) (152.719,362.306) (152.861,362.383) (153.004,362.459) (153.146,362.535) (153.288,362.612) (153.430,362.688) (153.572,362.765) (153.714,362.842) (153.856,362.919) (153.999,362.995) (154.141,363.072) (154.283,363.149) (154.425,363.226) (154.567,363.304) (154.709,363.381) (154.851,363.458) (154.994,363.536) (155.136,363.613) (155.278,363.691) (155.420,363.768) (155.562,363.846) (155.704,363.924) (155.846,364.002) (155.989,364.080) (156.131,364.158) (156.273,364.236) (156.415,364.314) (156.557,364.392) (156.699,364.470) (156.841,364.549) (156.984,364.627) (157.126,364.706) (157.268,364.784) (157.410,364.863) (157.552,364.942) (157.694,365.020) (157.836,365.099) (157.979,365.178) (158.121,365.257) (158.263,365.336) (158.405,365.415) (158.547,365.495) (158.689,365.574) (158.831,365.653) (158.974,365.733) (159.116,365.812) (159.258,365.892) (159.400,365.971) (159.542,366.051) (159.684,366.131) (159.826,366.211) (159.969,366.291) (160.111,366.371) (160.253,366.451) (160.395,366.531) (160.537,366.611) (160.679,366.691) (160.821,366.772) (160.964,366.852) (161.106,366.932) (161.248,367.013) (161.390,367.094) (161.532,367.174) (161.674,367.255) (161.816,367.336) (161.959,367.417) (162.101,367.498) (162.243,367.579) (162.385,367.660) (162.527,367.741) (162.669,367.822) (162.811,367.903) (162.954,367.984) (163.096,368.066) (163.238,368.147) (163.380,368.229) (163.522,368.310) (163.664,368.392) (163.806,368.474) (163.949,368.555) (164.091,368.637) (164.233,368.719) (164.375,368.801) (164.517,368.883) (164.659,368.965) (164.801,369.047) (164.944,369.129) (165.086,369.212) (165.228,369.294) (165.370,369.376) (165.512,369.459) (165.654,369.541) (165.796,369.624) (165.939,369.707) (166.081,369.789) (166.223,369.872) (166.365,369.955) (166.507,370.038) (166.649,370.120) (166.791,370.203) (166.934,370.287) (167.076,370.370) (167.218,370.453) (167.360,370.536) (167.502,370.619) (167.644,370.703) (167.786,370.786) (167.929,370.869) (168.071,370.953) (168.213,371.036) (168.355,371.120) (168.497,371.204) (168.639,371.287) (168.781,371.371) (168.924,371.455) (169.066,371.539) (169.208,371.623) (169.350,371.707) (169.492,371.791) (169.634,371.875) (169.776,371.959) (169.919,372.044) (170.061,372.128) (170.203,372.212) (170.345,372.297) (170.487,372.381) (170.629,372.465) (170.771,372.550) (170.914,372.635) (171.056,372.719) (171.198,372.804) (171.340,372.889) (171.482,372.974) (171.624,373.058) (171.766,373.143) (171.909,373.228) (172.051,373.313) (172.193,373.398) (172.335,373.484) (172.477,373.569) (172.619,373.654) (172.761,373.739) (172.904,373.825) (173.046,373.910) (173.188,373.995) (173.330,374.081) (173.472,374.166) (173.614,374.252) (173.756,374.338) (173.899,374.423) (174.041,374.509) (174.183,374.595) (174.325,374.681) (174.467,374.767) (174.609,374.853) (174.751,374.939) (174.894,375.025) (175.036,375.111) (175.178,375.197) (175.320,375.283) (175.462,375.369) (175.604,375.455) (175.746,375.542) (175.889,375.628) (176.031,375.714) (176.173,375.801) (176.315,375.887) (176.457,375.974) (176.599,376.060) (176.741,376.147) (176.884,376.234) (177.026,376.321) (177.168,376.407) (177.310,376.494) (177.452,376.581) (177.594,376.668) (177.736,376.755) (177.879,376.842) (178.021,376.929) (178.163,377.016) (178.305,377.103) (178.447,377.190) (178.589,377.277) (178.731,377.365) (178.874,377.452) (179.016,377.539) (179.158,377.627) (179.300,377.714) (179.442,377.802) (179.584,377.889) (179.726,377.977) (179.869,378.064) (180.011,378.152) (180.153,378.239) (180.295,378.327) (180.437,378.415) (180.579,378.503) (180.721,378.590) (180.864,378.678) (181.006,378.766) (181.148,378.854) (181.290,378.942) (181.432,379.030) (181.574,379.118) (181.716,379.206) (181.859,379.295) (182.001,379.383) (182.143,379.471) (182.285,379.559) (182.427,379.647) (182.569,379.736) (182.711,379.824) (182.854,379.913) (182.996,380.001) (183.138,380.089) (183.280,380.178) (183.422,380.267) (183.564,380.355) (183.706,380.444) (183.849,380.532) (183.991,380.621) (184.133,380.710) (184.275,380.799) (184.417,380.887) (184.559,380.976) (184.701,381.065) (184.844,381.154) (184.986,381.243) (185.128,381.332) (185.270,381.421) (185.412,381.510) (185.554,381.599) (185.696,381.688) (185.839,381.777) (185.981,381.866) (186.123,381.956) (186.265,382.045) (186.407,382.134) (186.549,382.223) (186.691,382.313) (186.834,382.402) (186.976,382.491) (187.118,382.581) (187.260,382.670) (187.402,382.760) (187.544,382.849) (187.686,382.939) (187.829,383.028) (187.971,383.118) (188.113,383.208) (188.255,383.297) (188.397,383.387) (188.539,383.477) (188.681,383.567) (188.824,383.656) (188.966,383.746) (189.108,383.836) (189.250,383.926) (189.392,384.016) (189.534,384.106) (189.676,384.196) (189.819,384.286) (189.961,384.376) (190.103,384.466) (190.245,384.556) (190.387,384.646) (190.529,384.736) (190.671,384.826) (190.814,384.916) (190.956,385.006) (191.098,385.097) (191.240,385.187) (191.382,385.277) (191.524,385.368) (191.666,385.458) (191.809,385.548) (191.951,385.639) (192.093,385.729) (192.235,385.819) (192.377,385.910) (192.519,386.000) (192.661,386.091) (192.804,386.181) (192.946,386.272) (193.088,386.362) (193.230,386.453) (193.372,386.544) (193.514,386.634) (193.656,386.725) (193.799,386.816) (193.941,386.906) (194.083,386.997) (194.225,387.088) (194.367,387.178) (194.509,387.269) (194.651,387.360) (194.794,387.451) (194.936,387.542) (195.078,387.633) (195.220,387.723) (195.362,387.814) (195.504,387.905) (195.646,387.996) (195.789,388.087) (195.931,388.178) (196.073,388.269) (196.215,388.360) (196.357,388.451) (196.499,388.542) (196.641,388.633) (196.784,388.724) (196.926,388.816) (197.068,388.907) (197.210,388.998) (197.352,389.089) (197.494,389.180) (197.636,389.271) (197.779,389.363) (197.921,389.454) (198.063,389.545) (198.205,389.636) (198.347,389.728) (198.489,389.819) (198.631,389.910) (198.774,390.002) (198.916,390.093) (199.058,390.184) (199.200,390.276) (199.342,390.367) (199.484,390.458) (199.626,390.550) (199.769,390.641) (199.911,390.733) (200.053,390.824) (200.195,390.916) (200.337,391.007) (200.479,391.099) (200.621,391.190) (200.764,391.282) (200.906,391.373) (201.048,391.465) (201.190,391.556) (201.332,391.648) (201.474,391.739) (201.616,391.831) (201.759,391.923) (201.901,392.014) (202.043,392.106) (202.185,392.197) (202.327,392.289) (202.469,392.381) (202.611,392.472) (202.754,392.564) (202.896,392.656) (203.038,392.747) (203.180,392.839) (203.322,392.931) (203.464,393.023) (203.606,393.114) (203.749,393.206) (203.891,393.298) (204.033,393.389) (204.175,393.481) (204.317,393.573) (204.459,393.665) (204.601,393.757) (204.744,393.848) (204.886,393.940) (205.028,394.032) (205.170,394.124) (205.312,394.216) (205.454,394.307) (205.596,394.399) (205.739,394.491) (205.881,394.583) (206.023,394.675) (206.165,394.767) (206.307,394.858) (206.449,394.950) (206.591,395.042) (206.734,395.134) (206.876,395.226) (207.018,395.318) (207.160,395.410) (207.302,395.502) (207.444,395.593) (207.586,395.685) (207.729,395.777) (207.871,395.869) (208.013,395.961) (208.155,396.053) (208.297,396.145) (208.439,396.237) (208.581,396.329) (208.724,396.421) (208.866,396.512) (209.008,396.604) (209.150,396.696) (209.292,396.788) (209.434,396.880) (209.576,396.972) (209.719,397.064) (209.861,397.156) (210.003,397.248) (210.145,397.340) (210.287,397.432) (210.429,397.523) (210.571,397.615) (210.714,397.707) (210.856,397.799) (210.998,397.891) (211.140,397.983) (211.282,398.075) (211.424,398.167) (211.566,398.259) (211.709,398.351) (211.851,398.442) (211.993,398.534) (212.135,398.626) (212.277,398.718) (212.419,398.810) (212.561,398.902) (212.704,398.994) (212.846,399.086) (212.988,399.178) (213.130,399.269) (213.272,399.361) (213.414,399.453) (213.556,399.545) (213.699,399.637) (213.841,399.729) (213.983,399.820) (214.125,399.912) (214.267,400.004) (214.409,400.096) (214.551,400.188) (214.694,400.280) (214.836,400.371) (214.978,400.463) (215.120,400.555) (215.262,400.647) (215.404,400.738) (215.546,400.830) (215.689,400.922) (215.831,401.014) (215.973,401.105) (216.115,401.197) (216.257,401.289) (216.399,401.381) (216.541,401.472) (216.684,401.564) (216.826,401.656) (216.968,401.747) (217.110,401.839) (217.252,401.931) (217.394,402.022) (217.536,402.114) (217.679,402.205) (217.821,402.297) (217.963,402.389) (218.105,402.480) (218.247,402.572) (218.389,402.663) (218.531,402.755) (218.674,402.846) (218.816,402.938) (218.958,403.029) (219.100,403.121) (219.242,403.212) (219.384,403.304) (219.526,403.395) (219.669,403.487) (219.811,403.578) (219.953,403.670) (220.095,403.761) (220.237,403.852) (220.379,403.944) (220.521,404.035) (220.664,404.127) (220.806,404.218) (220.948,404.309) (221.090,404.400) (221.232,404.492) (221.374,404.583) (221.516,404.674) (221.659,404.766) (221.801,404.857) (221.943,404.948) (222.085,405.039) (222.227,405.130) (222.369,405.221) (222.511,405.313) (222.654,405.404) (222.796,405.495) (222.938,405.586) (223.080,405.677) (223.222,405.768) (223.364,405.859) (223.506,405.950) (223.649,406.041) (223.791,406.132) (223.933,406.223) (224.075,406.314) (224.217,406.405) (224.359,406.496) (224.501,406.587) (224.644,406.677) (224.786,406.768) (224.928,406.859) (225.070,406.950) (225.212,407.041) (225.354,407.131) (225.496,407.222) (225.639,407.313) (225.781,407.404) (225.923,407.494) (226.065,407.585) (226.207,407.675) (226.349,407.766) (226.491,407.857) (226.634,407.947) (226.776,408.038) (226.918,408.128) (227.060,408.219) (227.202,408.309) (227.344,408.400) (227.486,408.490) (227.629,408.580) (227.771,408.671) (227.913,408.761) (228.055,408.851) (228.197,408.942) (228.339,409.032) (228.481,409.122) (228.624,409.212) (228.766,409.303) (228.908,409.393) (229.050,409.483) (229.192,409.573) (229.334,409.663) (229.476,409.753) (229.619,409.843) (229.761,409.933) (229.903,410.023) (230.045,410.113) (230.187,410.203) (230.329,410.293) (230.471,410.383) (230.614,410.472) (230.756,410.562) (230.898,410.652) (231.040,410.742) (231.182,410.831) (231.324,410.921) (231.466,411.011) (231.609,411.100) (231.751,411.190) (231.893,411.280) (232.035,411.369) (232.177,411.459) (232.319,411.548) (232.461,411.638) (232.604,411.727) (232.746,411.816) (232.888,411.906) (233.030,411.995) (233.172,412.084) (233.314,412.174) (233.456,412.263) (233.599,412.352) (233.741,412.441) (233.883,412.530) (234.025,412.619) (234.167,412.708) (234.309,412.797) (234.451,412.886) (234.594,412.975) (234.736,413.064) (234.878,413.153) (235.020,413.242) (235.162,413.331) (235.304,413.420) (235.446,413.508) (235.589,413.597) (235.731,413.686) (235.873,413.774) (236.015,413.863) (236.157,413.952) (236.299,414.040) (236.441,414.129) (236.584,414.217) (236.726,414.305) (236.868,414.394) (237.010,414.482) (237.152,414.570) (237.294,414.659) (237.436,414.747) (237.579,414.835) (237.721,414.923) (237.863,415.011) (238.005,415.100) (238.147,415.188) (238.289,415.276) (238.431,415.364) (238.574,415.451) (238.716,415.539) (238.858,415.627) (239.000,415.715)};
\fill[methodM] (64.481,339.121) circle (2pt);
\node[anchor=south east,font=\scriptsize,text=methodM,fill=white,inner sep=1pt] at (64.481,344.121) {$v_q$};
\fill[methodS] (92.756,339.121) circle (2pt);
\node[anchor=south west,font=\scriptsize,text=methodS,fill=white,inner sep=1pt] at (92.756,344.121) {$\bar v$};
\fill[methodM] (208.187,408.453) circle (2pt);
\node[anchor=south,font=\scriptsize,text=methodM,fill=white,inner sep=1pt] at (208.187,413.453) {$v_f$};
\end{scope}
\node[anchor=west,font=\small\bfseries,text=navy] at (5,226) {(c) CO$_2$: cumulative contribution};
\draw[rulegray,line width=.25pt] (40,29.000) -- (239,29.000);
\node[anchor=east,font=\scriptsize] at (36,29.000) {0};
\draw[rulegray,line width=.25pt] (40,69.227) -- (239,69.227);
\node[anchor=east,font=\scriptsize] at (36,69.227) {500};
\draw[rulegray,line width=.25pt] (40,109.455) -- (239,109.455);
\node[anchor=east,font=\scriptsize] at (36,109.455) {1,000};
\draw[rulegray,line width=.25pt] (40,149.682) -- (239,149.682);
\node[anchor=east,font=\scriptsize] at (36,149.682) {1,500};
\draw[rulegray,line width=.25pt] (40,189.909) -- (239,189.909);
\node[anchor=east,font=\scriptsize] at (36,189.909) {2,000};
\draw[black,line width=.35pt] (40.000,29) -- (40.000,26);
\node[anchor=north,font=\scriptsize] at (40.000,25) {0};
\draw[black,line width=.35pt] (89.750,29) -- (89.750,26);
\node[anchor=north,font=\scriptsize] at (89.750,25) {25};
\draw[black,line width=.35pt] (139.500,29) -- (139.500,26);
\node[anchor=north,font=\scriptsize] at (139.500,25) {50};
\draw[black,line width=.35pt] (189.250,29) -- (189.250,26);
\node[anchor=north,font=\scriptsize] at (189.250,25) {75};
\draw[black,line width=.35pt] (239.000,29) -- (239.000,26);
\node[anchor=north,font=\scriptsize] at (239.000,25) {100};
\draw[black,line width=.45pt] (40,206) -- (40,29) -- (239,29);
\node[anchor=north,font=\scriptsize] at (139.5,10) {Cumulative counted-vehicle share (\%)};
\node[rotate=90,font=\scriptsize] at (7,117.5) {Cumulative emission (kg/lane)};
\begin{scope}
\clip (40,29) rectangle (239,206);
\draw[methodO,line width=1.05pt,] plot coordinates {(40.000,29.000) (45.216,32.224) (50.212,34.783) (55.041,37.298) (59.607,39.676) (64.049,41.520) (68.392,43.379) (72.641,45.281) (76.916,47.217) (81.045,49.124) (85.172,51.084) (89.099,53.092) (92.904,55.295) (96.631,57.257) (100.375,59.376) (104.109,61.554) (107.659,63.714) (110.974,65.764) (114.063,68.025) (117.019,70.455) (119.943,73.191) (122.942,76.106) (125.913,78.930) (128.698,81.632) (131.366,84.258) (134.071,87.198) (136.980,90.555) (139.871,94.160) (142.622,97.689) (145.131,100.810) (147.571,103.805) (149.993,107.095) (152.539,110.753) (155.077,114.412) (157.665,117.850) (160.250,121.090) (162.934,124.073) (165.817,126.970) (168.815,129.734) (171.924,132.338) (175.062,134.641) (178.312,136.725) (181.680,138.750) (185.165,140.857) (188.696,142.789) (192.306,144.663) (195.944,146.718) (199.659,148.636) (203.401,150.428) (207.207,152.199) (211.191,153.911) (215.249,156.055) (219.341,158.630) (223.365,161.054) (227.456,163.324) (231.466,165.456) (235.359,167.440) (239.000,169.296)};
\draw[methodR,line width=1.05pt,] plot coordinates {(40.000,29.000) (40.005,29.003) (40.017,29.011) (40.035,29.023) (40.060,29.039) (40.091,29.059) (40.130,29.084) (40.174,29.113) (40.226,29.146) (40.284,29.184) (40.349,29.226) (40.420,29.272) (40.498,29.322) (40.582,29.377) (40.673,29.436) (40.771,29.499) (40.876,29.566) (40.986,29.638) (41.104,29.714) (41.228,29.795) (41.359,29.879) (41.496,29.968) (41.640,30.061) (41.790,30.158) (41.947,30.259) (42.110,30.365) (42.280,30.475) (42.456,30.589) (42.639,30.707) (42.828,30.830) (43.024,30.957) (43.226,31.087) (43.434,31.191) (43.649,31.298) (43.871,31.408) (44.099,31.522) (44.333,31.638) (44.573,31.758) (44.820,31.881) (45.074,32.007) (45.333,32.136) (45.599,32.269) (45.871,32.404) (46.150,32.543) (46.435,32.685) (46.726,32.830) (47.023,32.979) (47.326,33.130) (47.636,33.285) (47.952,33.443) (48.274,33.604) (48.602,33.769) (48.937,33.936) (49.277,34.107) (49.624,34.282) (49.976,34.459) (50.335,34.640) (50.700,34.824) (51.071,35.011) (51.448,35.202) (51.830,35.397) (52.219,35.595) (52.614,35.796) (53.015,36.001) (53.422,36.209) (53.834,36.422) (54.253,36.637) (54.677,36.857) (55.108,37.081) (55.544,37.308) (55.986,37.540) (56.434,37.776) (56.887,38.016) (57.347,38.260) (57.812,38.508) (58.283,38.761) (58.759,39.019) (59.241,39.282) (59.729,39.549) (60.223,39.822) (60.722,40.099) (61.227,40.383) (61.738,40.671) (62.254,40.965) (62.776,41.265) (63.303,41.572) (63.835,41.884) (64.374,42.203) (64.917,42.528) (65.466,42.861) (66.021,43.200) (66.581,43.547) (67.146,43.902) (67.716,44.264) (68.292,44.557) (68.873,44.857) (69.460,45.164) (70.051,45.478) (70.648,45.727) (71.250,45.982) (71.857,46.243) (72.470,46.510) (73.087,46.784) (73.709,47.064) (74.337,47.351) (74.969,47.646) (75.606,47.947) (76.248,48.256) (76.895,48.573) (77.547,48.898) (78.203,49.231) (78.864,49.572) (79.530,49.922) (80.200,50.281) (80.875,50.649) (81.555,51.026) (82.239,51.412) (82.927,51.809) (83.620,52.137) (84.317,52.474) (85.019,52.818) (85.724,53.172) (86.434,53.534) (87.148,53.905) (87.866,54.286) (88.588,54.675) (89.313,55.074) (90.043,55.483) (90.776,55.901) (91.514,56.329) (92.254,56.768) (92.999,57.216) (93.747,57.675) (94.498,58.144) (95.253,58.567) (96.011,59.000) (96.772,59.442) (97.537,59.893) (98.304,60.355) (99.075,60.825) (99.848,61.306) (100.624,61.796) (101.404,62.296) (102.185,62.806) (102.970,63.326) (103.757,63.855) (104.546,64.395) (105.338,64.944) (106.132,65.503) (106.928,66.071) (107.726,66.649) (108.527,67.237) (109.329,67.834) (110.133,68.441) (110.939,69.057) (111.747,69.682) (112.556,70.317) (113.367,70.960) (114.179,71.612) (114.993,72.273) (115.808,72.943) (116.624,73.621) (117.441,74.307) (118.259,75.000) (119.078,75.702) (119.898,76.411) (120.719,77.128) (121.540,77.851) (122.362,78.582) (123.184,79.319) (124.007,80.062) (124.831,80.811) (125.654,81.566) (126.478,82.327) (127.301,83.092) (128.125,83.863) (128.949,84.638) (129.773,85.417) (130.596,86.201) (131.419,86.987) (132.242,87.778) (133.064,88.571) (133.886,89.367) (134.708,90.165) (135.528,90.964) (136.348,91.766) (137.168,92.569) (137.986,93.373) (138.804,94.177) (139.621,94.981) (140.436,95.786) (141.251,96.589) (142.065,97.393) (142.877,98.194) (143.688,98.995) (144.498,99.794) (145.307,100.590) (146.114,101.384) (146.920,102.175) (147.724,102.963) (148.527,103.748) (149.328,104.529) (150.128,105.306) (150.926,106.079) (151.722,106.847) (152.517,107.610) (153.310,108.368) (154.101,109.120) (154.890,109.867) (155.678,110.608) (156.463,111.343) (157.247,112.071) (158.028,112.793) (158.808,113.508) (159.586,114.216) (160.362,114.917) (161.135,115.610) (161.907,116.295) (162.676,116.973) (163.444,117.643) (164.209,118.305) (164.972,118.959) (165.733,119.604) (166.491,120.241) (167.248,120.869) (168.002,121.488) (168.753,122.099) (169.503,122.701) (170.250,123.294) (170.995,123.878) (171.738,124.453) (172.478,125.018) (173.216,125.575) (173.951,126.123) (174.684,126.661) (175.414,127.190) (176.143,127.710) (176.868,128.220) (177.591,128.722) (178.312,129.214) (179.030,129.697) (179.746,130.171) (180.459,130.636) (181.169,131.092) (181.877,131.539) (182.582,131.977) (183.284,132.407) (183.984,132.827) (184.681,133.239) (185.376,133.642) (186.068,134.037) (186.757,134.424) (187.443,134.852) (188.126,135.271) (188.807,135.681) (189.484,136.082) (190.159,136.474) (190.831,136.857) (191.500,137.232) (192.166,137.598) (192.829,137.956) (193.489,138.305) (194.146,138.647) (194.800,138.981) (195.451,139.307) (196.098,139.625) (196.742,139.936) (197.383,140.240) (198.021,140.607) (198.656,140.966) (199.287,141.316) (199.915,141.658) (200.539,141.992) (201.160,142.319) (201.778,142.637) (202.392,142.949) (203.002,143.253) (203.609,143.550) (204.212,143.840) (204.811,144.124) (205.406,144.401) (205.998,144.672) (206.586,144.937) (207.170,145.196) (207.750,145.449) (208.326,145.697) (208.899,145.939) (209.467,146.176) (210.031,146.476) (210.590,146.768) (211.146,147.055) (211.697,147.336) (212.244,147.683) (212.787,148.023) (213.325,148.357) (213.859,148.684) (214.388,149.004) (214.913,149.318) (215.433,149.627) (215.949,149.929) (216.460,150.226) (216.966,150.517) (217.467,150.803) (217.964,151.084) (218.456,151.359) (218.942,151.630) (219.424,151.896) (219.901,152.157) (220.373,152.414) (220.840,152.667) (221.301,152.915) (221.757,153.159) (222.209,153.398) (222.655,153.634) (223.095,153.866) (223.530,154.094) (223.960,154.319) (224.385,154.539) (224.804,154.756) (225.217,154.969) (225.625,155.179) (226.028,155.385) (226.425,155.588) (226.816,155.788) (227.201,155.984) (227.581,156.177) (227.955,156.366) (228.323,156.552) (228.685,156.735) (229.042,156.915) (229.393,157.091) (229.737,157.265) (230.076,157.435) (230.409,157.602) (230.736,157.765) (231.057,157.926) (231.372,158.083) (231.680,158.238) (231.983,158.389) (232.279,158.537) (232.570,158.682) (232.854,158.823) (233.132,158.962) (233.404,159.097) (233.669,159.230) (233.928,159.359) (234.181,159.485) (234.428,159.607) (234.668,159.727) (234.902,159.844) (235.130,159.957) (235.351,160.067) (235.566,160.174) (235.775,160.278) (235.977,160.408) (236.172,160.535) (236.361,160.657) (236.544,160.776) (236.720,160.890) (236.890,161.000) (237.054,161.105) (237.210,161.207) (237.361,161.304) (237.504,161.397) (237.641,161.486) (237.772,161.570) (237.896,161.650) (238.014,161.726) (238.124,161.798) (238.229,161.866) (238.327,161.929) (238.418,161.988) (238.502,162.043) (238.580,162.093) (238.651,162.139) (238.716,162.181) (238.774,162.219) (238.826,162.252) (238.870,162.281) (238.909,162.306) (238.940,162.326) (238.965,162.342) (238.983,162.354) (238.995,162.361) (239.000,162.365)};
\draw[methodS,line width=1.05pt,dashed] plot coordinates {(40.000,29.000) (239.000,146.916)};
\draw[methodM,line width=1.05pt,] plot coordinates {(40.000,29.000) (40.005,29.003) (40.017,29.011) (40.035,29.023) (40.060,29.039) (40.091,29.059) (40.130,29.084) (40.174,29.113) (40.226,29.146) (40.284,29.184) (40.349,29.226) (40.420,29.272) (40.498,29.322) (40.582,29.377) (40.673,29.436) (40.771,29.499) (40.876,29.566) (40.986,29.638) (41.104,29.714) (41.228,29.795) (41.359,29.879) (41.496,29.968) (41.640,30.061) (41.790,30.158) (41.947,30.259) (42.110,30.365) (42.280,30.475) (42.456,30.589) (42.639,30.707) (42.828,30.830) (43.024,30.956) (43.226,31.087) (43.434,31.222) (43.649,31.361) (43.871,31.505) (44.099,31.652) (44.333,31.804) (44.573,31.959) (44.820,32.119) (45.074,32.283) (45.333,32.451) (45.599,32.623) (45.871,32.799) (46.150,32.980) (46.435,33.164) (46.726,33.352) (47.023,33.545) (47.326,33.741) (47.636,33.942) (47.952,34.146) (48.274,34.355) (48.602,34.567) (48.937,34.784) (49.277,35.005) (49.624,35.229) (49.976,35.458) (50.335,35.691) (50.700,35.927) (51.071,36.168) (51.448,36.412) (51.830,36.661) (52.219,36.913) (52.614,37.170) (53.015,37.430) (53.422,37.695) (53.834,37.963) (54.253,38.236) (54.677,38.512) (55.108,38.792) (55.544,39.077) (55.986,39.365) (56.434,39.658) (56.887,39.954) (57.347,40.254) (57.812,40.559) (58.283,40.867) (58.759,41.180) (59.241,41.497) (59.729,41.818) (60.223,42.143) (60.722,42.472) (61.227,42.805) (61.738,43.142) (62.254,43.484) (62.776,43.830) (63.303,44.180) (63.835,44.534) (64.374,44.893) (64.917,45.256) (65.466,45.623) (66.021,45.995) (66.581,46.372) (67.146,46.752) (67.716,47.138) (68.292,47.527) (68.873,47.922) (69.460,48.321) (70.051,48.725) (70.648,49.133) (71.250,49.547) (71.857,49.965) (72.470,50.388) (73.087,50.816) (73.709,51.248) (74.337,51.686) (74.969,52.129) (75.606,52.577) (76.248,53.030) (76.895,53.489) (77.547,53.952) (78.203,54.421) (78.864,54.896) (79.530,55.375) (80.200,55.861) (80.875,56.352) (81.555,56.848) (82.239,57.350) (82.927,57.857) (83.620,58.371) (84.317,58.890) (85.019,59.415) (85.724,59.946) (86.434,60.482) (87.148,61.025) (87.866,61.573) (88.588,62.128) (89.313,62.689) (90.043,63.255) (90.776,63.828) (91.514,64.407) (92.254,64.992) (92.999,65.583) (93.747,66.180) (94.498,66.784) (95.253,67.394) (96.011,68.009) (96.772,68.632) (97.537,69.260) (98.304,69.894) (99.075,70.535) (99.848,71.182) (100.624,71.834) (101.404,72.493) (102.185,73.158) (102.970,73.829) (103.757,74.506) (104.546,75.189) (105.338,75.878) (106.132,76.573) (106.928,77.273) (107.726,77.979) (108.527,78.691) (109.329,79.408) (110.133,80.131) (110.939,80.859) (111.747,81.592) (112.556,82.330) (113.367,83.074) (114.179,83.822) (114.993,84.575) (115.808,85.333) (116.624,86.095) (117.441,86.862) (118.259,87.633) (119.078,88.408) (119.898,89.187) (120.719,89.970) (121.540,90.757) (122.362,91.547) (123.184,92.341) (124.007,93.137) (124.831,93.937) (125.654,94.740) (126.478,95.545) (127.301,96.353) (128.125,97.163) (128.949,97.975) (129.773,98.789) (130.596,99.605) (131.419,100.422) (132.242,101.241) (133.064,102.060) (133.886,102.881) (134.708,103.703) (135.528,104.525) (136.348,105.348) (137.168,106.170) (137.986,106.993) (138.804,107.815) (139.621,108.637) (140.436,109.459) (141.251,110.280) (142.065,111.099) (142.877,111.918) (143.688,112.735) (144.498,113.550) (145.307,114.364) (146.114,115.176) (146.920,115.986) (147.724,116.794) (148.527,117.599) (149.328,118.402) (150.128,119.202) (150.926,119.999) (151.722,120.792) (152.517,121.583) (153.310,122.370) (154.101,123.154) (154.890,123.934) (155.678,124.711) (156.463,125.483) (157.247,126.252) (158.028,127.016) (158.808,127.776) (159.586,128.531) (160.362,129.282) (161.135,130.029) (161.907,130.771) (162.676,131.508) (163.444,132.240) (164.209,132.967) (164.972,133.689) (165.733,134.406) (166.491,135.118) (167.248,135.824) (168.002,136.526) (168.753,137.222) (169.503,137.912) (170.250,138.597) (170.995,139.277) (171.738,139.951) (172.478,140.619) (173.216,141.282) (173.951,141.940) (174.684,142.591) (175.414,143.237) (176.143,143.878) (176.868,144.513) (177.591,145.142) (178.312,145.765) (179.030,146.383) (179.746,146.995) (180.459,147.601) (181.169,148.202) (181.877,148.798) (182.582,149.387) (183.284,149.971) (183.984,150.550) (184.681,151.123) (185.376,151.690) (186.068,152.252) (186.757,152.809) (187.443,153.360) (188.126,153.906) (188.807,154.446) (189.484,154.981) (190.159,155.511) (190.831,156.036) (191.500,156.556) (192.166,157.070) (192.829,157.579) (193.489,158.084) (194.146,158.583) (194.800,159.077) (195.451,159.567) (196.098,160.051) (196.742,160.531) (197.383,161.006) (198.021,161.476) (198.656,161.942) (199.287,162.403) (199.915,162.860) (200.539,163.312) (201.160,163.759) (201.778,164.202) (202.392,164.641) (203.002,165.075) (203.609,165.505) (204.212,165.930) (204.811,166.352) (205.406,166.769) (205.998,167.182) (206.586,167.591) (207.170,167.996) (207.750,168.397) (208.326,168.793) (208.899,169.186) (209.467,169.575) (210.031,169.960) (210.590,170.341) (211.146,170.718) (211.697,171.091) (212.244,171.461) (212.787,171.826) (213.325,172.188) (213.859,172.546) (214.388,172.900) (214.913,173.250) (215.433,173.597) (215.949,173.940) (216.460,174.280) (216.966,174.615) (217.467,174.947) (217.964,175.275) (218.456,175.600) (218.942,175.921) (219.424,176.238) (219.901,176.551) (220.373,176.861) (220.840,177.167) (221.301,177.470) (221.757,177.768) (222.209,178.064) (222.655,178.355) (223.095,178.643) (223.530,178.927) (223.960,179.207) (224.385,179.484) (224.804,179.756) (225.217,180.026) (225.625,180.291) (226.028,180.553) (226.425,180.810) (226.816,181.064) (227.201,181.315) (227.581,181.561) (227.955,181.804) (228.323,182.043) (228.685,182.278) (229.042,182.509) (229.393,182.736) (229.737,182.960) (230.076,183.179) (230.409,183.395) (230.736,183.607) (231.057,183.815) (231.372,184.018) (231.680,184.218) (231.983,184.414) (232.279,184.606) (232.570,184.794) (232.854,184.978) (233.132,185.158) (233.404,185.334) (233.669,185.506) (233.928,185.673) (234.181,185.837) (234.428,185.997) (234.668,186.152) (234.902,186.304) (235.130,186.451) (235.351,186.594) (235.566,186.733) (235.775,186.868) (235.977,186.999) (236.172,187.126) (236.361,187.248) (236.544,187.366) (236.720,187.480) (236.890,187.590) (237.054,187.696) (237.210,187.797) (237.361,187.894) (237.504,187.987) (237.641,188.076) (237.772,188.161) (237.896,188.241) (238.014,188.317) (238.124,188.389) (238.229,188.456) (238.327,188.519) (238.418,188.578) (238.502,188.633) (238.580,188.683) (238.651,188.730) (238.716,188.771) (238.774,188.809) (238.826,188.842) (238.870,188.871) (238.909,188.896) (238.940,188.916) (238.965,188.933) (238.983,188.944) (238.995,188.952) (239.000,188.955)};
\draw[methodA,line width=1.05pt,dashed] plot coordinates {(40.000,29.000) (40.005,29.003) (40.017,29.009) (40.035,29.020) (40.060,29.034) (40.091,29.052) (40.130,29.074) (40.174,29.099) (40.226,29.129) (40.284,29.162) (40.349,29.199) (40.420,29.239) (40.498,29.284) (40.582,29.332) (40.673,29.384) (40.771,29.440) (40.876,29.499) (40.986,29.563) (41.104,29.630) (41.228,29.700) (41.359,29.775) (41.496,29.853) (41.640,29.935) (41.790,30.021) (41.947,30.110) (42.110,30.203) (42.280,30.300) (42.456,30.401) (42.639,30.505) (42.828,30.613) (43.024,30.725) (43.226,30.840) (43.434,30.959) (43.649,31.081) (43.871,31.208) (44.099,31.338) (44.333,31.471) (44.573,31.609) (44.820,31.749) (45.074,31.894) (45.333,32.042) (45.599,32.194) (45.871,32.349) (46.150,32.508) (46.435,32.671) (46.726,32.837) (47.023,33.006) (47.326,33.179) (47.636,33.356) (47.952,33.537) (48.274,33.721) (48.602,33.908) (48.937,34.099) (49.277,34.294) (49.624,34.492) (49.976,34.693) (50.335,34.898) (50.700,35.107) (51.071,35.319) (51.448,35.535) (51.830,35.754) (52.219,35.977) (52.614,36.203) (53.015,36.433) (53.422,36.667) (53.834,36.904) (54.253,37.144) (54.677,37.388) (55.108,37.636) (55.544,37.887) (55.986,38.142) (56.434,38.400) (56.887,38.662) (57.347,38.928) (57.812,39.197) (58.283,39.470) (58.759,39.746) (59.241,40.026) (59.729,40.310) (60.223,40.598) (60.722,40.889) (61.227,41.185) (61.738,41.484) (62.254,41.786) (62.776,42.093) (63.303,42.404) (63.835,42.718) (64.374,43.037) (64.917,43.359) (65.466,43.686) (66.021,44.017) (66.581,44.351) (67.146,44.690) (67.716,45.034) (68.292,45.381) (68.873,45.733) (69.460,46.089) (70.051,46.450) (70.648,46.815) (71.250,47.184) (71.857,47.559) (72.470,47.937) (73.087,48.321) (73.709,48.709) (74.337,49.102) (74.969,49.500) (75.606,49.903) (76.248,50.311) (76.895,50.724) (77.547,51.142) (78.203,51.565) (78.864,51.994) (79.530,52.428) (80.200,52.867) (80.875,53.312) (81.555,53.762) (82.239,54.217) (82.927,54.679) (83.620,55.146) (84.317,55.619) (85.019,56.097) (85.724,56.581) (86.434,57.072) (87.148,57.568) (87.866,58.070) (88.588,58.579) (89.313,59.093) (90.043,59.614) (90.776,60.140) (91.514,60.673) (92.254,61.213) (92.999,61.758) (93.747,62.310) (94.498,62.868) (95.253,63.433) (96.011,64.003) (96.772,64.581) (97.537,65.164) (98.304,65.754) (99.075,66.351) (99.848,66.953) (100.624,67.562) (101.404,68.178) (102.185,68.799) (102.970,69.427) (103.757,70.062) (104.546,70.702) (105.338,71.348) (106.132,72.001) (106.928,72.660) (107.726,73.324) (108.527,73.995) (109.329,74.671) (110.133,75.353) (110.939,76.041) (111.747,76.734) (112.556,77.433) (113.367,78.137) (114.179,78.846) (114.993,79.560) (115.808,80.280) (116.624,81.004) (117.441,81.733) (118.259,82.466) (119.078,83.204) (119.898,83.947) (120.719,84.693) (121.540,85.443) (122.362,86.197) (123.184,86.955) (124.007,87.716) (124.831,88.481) (125.654,89.248) (126.478,90.019) (127.301,90.792) (128.125,91.568) (128.949,92.346) (129.773,93.126) (130.596,93.908) (131.419,94.692) (132.242,95.478) (133.064,96.265) (133.886,97.053) (134.708,97.841) (135.528,98.631) (136.348,99.421) (137.168,100.212) (137.986,101.002) (138.804,101.792) (139.621,102.583) (140.436,103.372) (141.251,104.161) (142.065,104.949) (142.877,105.736) (143.688,106.521) (144.498,107.305) (145.307,108.087) (146.114,108.867) (146.920,109.646) (147.724,110.422) (148.527,111.195) (149.328,111.966) (150.128,112.734) (150.926,113.499) (151.722,114.261) (152.517,115.019) (153.310,115.775) (154.101,116.526) (154.890,117.274) (155.678,118.018) (156.463,118.757) (157.247,119.493) (158.028,120.224) (158.808,120.951) (159.586,121.673) (160.362,122.391) (161.135,123.104) (161.907,123.812) (162.676,124.515) (163.444,125.212) (164.209,125.905) (164.972,126.593) (165.733,127.275) (166.491,127.951) (167.248,128.623) (168.002,129.288) (168.753,129.948) (169.503,130.603) (170.250,131.252) (170.995,131.895) (171.738,132.532) (172.478,133.164) (173.216,133.790) (173.951,134.409) (174.684,135.023) (175.414,135.632) (176.143,136.234) (176.868,136.830) (177.591,137.421) (178.312,138.005) (179.030,138.584) (179.746,139.157) (180.459,139.724) (181.169,140.285) (181.877,140.840) (182.582,141.390) (183.284,141.933) (183.984,142.471) (184.681,143.003) (185.376,143.530) (186.068,144.051) (186.757,144.566) (187.443,145.076) (188.126,145.580) (188.807,146.079) (189.484,146.572) (190.159,147.060) (190.831,147.543) (191.500,148.020) (192.166,148.492) (192.829,148.959) (193.489,149.421) (194.146,149.877) (194.800,150.329) (195.451,150.776) (196.098,151.218) (196.742,151.654) (197.383,152.087) (198.021,152.514) (198.656,152.937) (199.287,153.355) (199.915,153.769) (200.539,154.178) (201.160,154.582) (201.778,154.982) (202.392,155.378) (203.002,155.770) (203.609,156.157) (204.212,156.540) (204.811,156.919) (205.406,157.294) (205.998,157.665) (206.586,158.032) (207.170,158.395) (207.750,158.753) (208.326,159.109) (208.899,159.460) (209.467,159.807) (210.031,160.151) (210.590,160.491) (211.146,160.827) (211.697,161.160) (212.244,161.489) (212.787,161.814) (213.325,162.136) (213.859,162.455) (214.388,162.770) (214.913,163.081) (215.433,163.389) (215.949,163.693) (216.460,163.994) (216.966,164.292) (217.467,164.586) (217.964,164.877) (218.456,165.165) (218.942,165.449) (219.424,165.729) (219.901,166.007) (220.373,166.281) (220.840,166.552) (221.301,166.819) (221.757,167.083) (222.209,167.344) (222.655,167.602) (223.095,167.856) (223.530,168.107) (223.960,168.354) (224.385,168.599) (224.804,168.840) (225.217,169.077) (225.625,169.311) (226.028,169.542) (226.425,169.770) (226.816,169.994) (227.201,170.215) (227.581,170.432) (227.955,170.646) (228.323,170.857) (228.685,171.064) (229.042,171.268) (229.393,171.469) (229.737,171.666) (230.076,171.859) (230.409,172.050) (230.736,172.236) (231.057,172.419) (231.372,172.599) (231.680,172.775) (231.983,172.948) (232.279,173.117) (232.570,173.283) (232.854,173.445) (233.132,173.604) (233.404,173.759) (233.669,173.910) (233.928,174.058) (234.181,174.203) (234.428,174.343) (234.668,174.480) (234.902,174.614) (235.130,174.744) (235.351,174.870) (235.566,174.993) (235.775,175.111) (235.977,175.227) (236.172,175.338) (236.361,175.446) (236.544,175.550) (236.720,175.651) (236.890,175.748) (237.054,175.841) (237.210,175.930) (237.361,176.016) (237.504,176.098) (237.641,176.176) (237.772,176.251) (237.896,176.321) (238.014,176.389) (238.124,176.452) (238.229,176.511) (238.327,176.567) (238.418,176.619) (238.502,176.667) (238.580,176.712) (238.651,176.752) (238.716,176.789) (238.774,176.822) (238.826,176.852) (238.870,176.877) (238.909,176.899) (238.940,176.917) (238.965,176.931) (238.983,176.942) (238.995,176.948) (239.000,176.951)};
\end{scope}
\node[anchor=west,font=\small\bfseries,text=navy] at (264,476) {(b) NO$_x$: speed-to-rate mapping};
\draw[rulegray,line width=.25pt] (299,304.000) -- (498,304.000);
\node[anchor=east,font=\scriptsize] at (295,304.000) {0};
\draw[rulegray,line width=.25pt] (299,334.400) -- (498,334.400);
\node[anchor=east,font=\scriptsize] at (295,334.400) {10};
\draw[rulegray,line width=.25pt] (299,364.800) -- (498,364.800);
\node[anchor=east,font=\scriptsize] at (295,364.800) {20};
\draw[rulegray,line width=.25pt] (299,395.200) -- (498,395.200);
\node[anchor=east,font=\scriptsize] at (295,395.200) {30};
\draw[rulegray,line width=.25pt] (299,425.600) -- (498,425.600);
\node[anchor=east,font=\scriptsize] at (295,425.600) {40};
\draw[rulegray,line width=.25pt] (299,456.000) -- (498,456.000);
\node[anchor=east,font=\scriptsize] at (295,456.000) {50};
\draw[black,line width=.35pt] (313.214,304) -- (313.214,301);
\node[anchor=north,font=\scriptsize] at (313.214,300) {10};
\draw[black,line width=.35pt] (355.857,304) -- (355.857,301);
\node[anchor=north,font=\scriptsize] at (355.857,300) {25};
\draw[black,line width=.35pt] (398.500,304) -- (398.500,301);
\node[anchor=north,font=\scriptsize] at (398.500,300) {40};
\draw[black,line width=.35pt] (441.143,304) -- (441.143,301);
\node[anchor=north,font=\scriptsize] at (441.143,300) {55};
\draw[black,line width=.35pt] (483.786,304) -- (483.786,301);
\node[anchor=north,font=\scriptsize] at (483.786,300) {70};
\draw[black,line width=.45pt] (299,456) -- (299,304) -- (498,304);
\node[anchor=north,font=\scriptsize] at (398.5,285) {Speed (mph)};
\node[rotate=90,font=\scriptsize] at (266,380.0) {Rate (g/veh h)};
\begin{scope}
\clip (299,304) rectangle (498,456);
\draw[methodO,line width=1.05pt,] plot coordinates {(299.000,316.391) (299.142,316.391) (299.284,316.391) (299.426,316.391) (299.569,316.391) (299.711,316.391) (299.853,316.391) (299.995,316.391) (300.137,316.391) (300.279,316.391) (300.421,316.391) (300.564,316.391) (300.706,316.391) (300.848,316.391) (300.990,316.391) (301.132,316.391) (301.274,316.391) (301.416,316.391) (301.559,316.391) (301.701,316.391) (301.843,316.391) (301.985,316.391) (302.127,316.391) (302.269,316.391) (302.411,316.391) (302.554,316.391) (302.696,316.391) (302.838,316.391) (302.980,316.391) (303.122,316.391) (303.264,316.391) (303.406,316.391) (303.549,316.391) (303.691,316.391) (303.833,316.391) (303.975,316.391) (304.117,316.391) (304.259,316.391) (304.401,316.391) (304.544,316.391) (304.686,316.391) (304.828,316.391) (304.970,316.391) (305.112,316.391) (305.254,316.391) (305.396,316.391) (305.539,316.391) (305.681,316.391) (305.823,316.391) (305.965,316.391) (306.107,316.391) (306.249,316.391) (306.391,316.391) (306.534,316.391) (306.676,316.391) (306.818,316.391) (306.960,316.391) (307.102,316.391) (307.244,316.391) (307.386,316.391) (307.529,316.391) (307.671,316.391) (307.813,316.391) (307.955,316.391) (308.097,316.391) (308.239,316.391) (308.381,316.391) (308.524,316.391) (308.666,316.391) (308.808,316.391) (308.950,316.391) (309.092,316.391) (309.234,316.391) (309.376,316.391) (309.519,316.391) (309.661,316.391) (309.803,316.391) (309.945,316.391) (310.087,316.391) (310.229,316.391) (310.371,316.391) (310.514,316.391) (310.656,316.391) (310.798,316.391) (310.940,316.391) (311.082,316.391) (311.224,316.391) (311.366,316.391) (311.509,316.391) (311.651,316.391) (311.793,316.391) (311.935,316.391) (312.077,316.391) (312.219,316.391) (312.361,316.391) (312.504,316.391) (312.646,316.391) (312.788,316.391) (312.930,316.391) (313.072,316.391) (313.214,316.391) (313.356,316.391) (313.499,316.391) (313.641,316.391) (313.783,316.391) (313.925,316.391) (314.067,316.391) (314.209,316.391) (314.351,316.391) (314.494,316.391) (314.636,316.391) (314.778,316.391) (314.920,316.391) (315.062,316.391) (315.204,316.391) (315.346,316.391) (315.489,316.391) (315.631,316.391) (315.773,316.391) (315.915,316.391) (316.057,316.391) (316.199,316.391) (316.341,316.391) (316.484,316.391) (316.626,316.391) (316.768,316.391) (316.910,316.391) (317.052,316.391) (317.194,316.391) (317.336,316.391) (317.479,316.391) (317.621,316.391) (317.763,316.391) (317.905,316.391) (318.047,316.391) (318.189,316.391) (318.331,316.391) (318.474,316.391) (318.616,316.391) (318.758,316.391) (318.900,316.391) (319.042,316.391) (319.184,316.391) (319.326,316.391) (319.469,316.391) (319.611,316.391) (319.753,316.391) (319.895,316.391) (320.037,316.391) (320.179,316.391) (320.321,316.391) (320.464,316.391) (320.606,316.391) (320.748,316.391) (320.890,316.391) (321.032,316.391) (321.174,316.391) (321.316,316.391) (321.459,316.391) (321.601,316.391) (321.743,316.391) (321.885,316.391) (322.027,316.391) (322.169,316.391) (322.311,316.391) (322.454,316.391) (322.596,316.391) (322.738,316.391) (322.880,316.391) (323.022,316.391) (323.164,316.391) (323.306,316.391) (323.449,316.391) (323.591,316.391) (323.733,316.391) (323.875,316.391) (324.017,316.391) (324.159,316.391) (324.301,316.391) (324.444,316.391) (324.586,316.391) (324.728,316.391) (324.870,316.391) (325.012,316.391) (325.154,316.391) (325.296,316.391) (325.439,316.391) (325.581,316.391) (325.723,316.391) (325.865,316.391) (326.007,316.391) (326.149,316.391) (326.291,316.391) (326.434,316.391) (326.576,316.391) (326.718,316.391) (326.860,316.391) (327.002,316.391) (327.144,316.391) (327.286,316.391) (327.429,316.391) (327.571,316.391) (327.713,316.391) (327.855,316.391) (327.997,316.391) (328.139,316.391) (328.281,316.391) (328.424,316.391) (328.566,316.391) (328.708,316.391) (328.850,316.391) (328.992,316.391) (329.134,316.391) (329.276,316.391) (329.419,316.391) (329.561,316.391) (329.703,316.391) (329.845,316.391) (329.987,316.391) (330.129,316.391) (330.271,316.391) (330.414,316.391) (330.556,316.391) (330.698,316.391) (330.840,316.391) (330.982,316.391) (331.124,316.391) (331.266,316.391) (331.409,316.391) (331.551,316.391) (331.693,316.391) (331.835,316.391) (331.977,316.391) (332.119,316.391) (332.261,316.391) (332.404,316.391) (332.546,316.391) (332.688,316.391) (332.830,316.391) (332.972,316.391) (333.114,316.391) (333.256,316.391) (333.399,316.391) (333.541,316.391) (333.683,316.391) (333.825,316.391) (333.967,316.391) (334.109,316.391) (334.251,316.391) (334.394,316.391) (334.536,316.391) (334.678,316.391) (334.820,316.391) (334.962,316.391) (335.104,316.391) (335.246,316.391) (335.389,316.391) (335.531,316.391) (335.673,316.391) (335.815,316.391) (335.957,316.391) (336.099,316.391) (336.241,316.391) (336.384,316.391) (336.526,316.391) (336.668,316.391) (336.810,316.391) (336.952,316.391) (337.094,316.391) (337.236,316.391) (337.379,316.391) (337.521,316.391) (337.663,316.391) (337.805,316.391) (337.947,316.391) (338.089,316.391) (338.231,316.391) (338.374,316.391) (338.516,316.391) (338.658,316.391) (338.800,316.391) (338.942,316.391) (339.084,316.391) (339.226,316.391) (339.369,316.391) (339.511,316.391) (339.653,316.391) (339.795,316.391) (339.937,316.391) (340.079,316.391) (340.221,316.391) (340.364,316.391) (340.506,316.391) (340.648,316.391) (340.790,316.391) (340.932,316.391) (341.074,316.391) (341.216,316.391) (341.359,316.391) (341.501,316.391) (341.643,316.391) (341.785,316.391) (341.927,316.391) (342.069,316.391) (342.211,316.391) (342.354,316.391) (342.496,316.391) (342.638,316.391) (342.780,316.391) (342.922,316.391) (343.064,316.391) (343.206,316.391) (343.349,316.391) (343.491,316.391) (343.633,316.391) (343.775,316.391) (343.917,316.391) (344.059,316.391) (344.201,316.391) (344.344,316.391) (344.486,316.391) (344.628,316.391) (344.770,316.391) (344.912,316.391) (345.054,316.391) (345.196,316.391) (345.339,316.391) (345.481,316.391) (345.623,316.391) (345.765,316.391) (345.907,316.391) (346.049,316.391) (346.191,316.391) (346.334,316.391) (346.476,316.391) (346.618,316.391) (346.760,316.391) (346.902,316.391) (347.044,316.391) (347.186,316.391) (347.329,316.391) (347.471,316.391) (347.613,316.391) (347.755,316.391) (347.897,316.391) (348.039,316.391) (348.181,316.391) (348.324,316.391) (348.466,316.391) (348.608,316.391) (348.750,316.391) (348.892,316.391) (349.034,316.391) (349.176,316.391) (349.319,316.391) (349.461,316.391) (349.603,316.391) (349.745,316.391) (349.887,316.391) (350.029,316.391) (350.171,316.391) (350.314,316.391) (350.456,316.391) (350.598,316.391) (350.740,316.391) (350.882,316.391) (351.024,316.391) (351.166,316.391) (351.309,316.391) (351.451,316.391) (351.593,316.391) (351.735,316.391) (351.877,316.391) (352.019,316.391) (352.161,316.391) (352.304,316.391) (352.446,316.391) (352.588,316.391) (352.730,316.391) (352.872,316.391) (353.014,316.391) (353.156,316.391) (353.299,316.391) (353.441,316.391) (353.583,316.391) (353.725,316.391) (353.867,316.391) (354.009,316.391) (354.151,316.391) (354.294,316.391) (354.436,316.391) (354.578,316.391) (354.720,316.391) (354.862,316.391) (355.004,316.391) (355.146,316.391) (355.289,316.391) (355.431,316.391) (355.573,316.391) (355.715,316.391) (355.857,322.461) (355.999,322.461) (356.141,322.461) (356.284,322.461) (356.426,322.461) (356.568,322.461) (356.710,322.461) (356.852,322.461) (356.994,322.461) (357.136,322.461) (357.279,322.461) (357.421,322.461) (357.563,322.461) (357.705,322.461) (357.847,322.461) (357.989,322.461) (358.131,322.461) (358.274,322.461) (358.416,322.461) (358.558,322.461) (358.700,322.461) (358.842,322.461) (358.984,322.461) (359.126,322.461) (359.269,322.461) (359.411,322.461) (359.553,322.461) (359.695,322.461) (359.837,322.461) (359.979,322.461) (360.121,322.461) (360.264,322.461) (360.406,322.461) (360.548,322.461) (360.690,322.461) (360.832,322.461) (360.974,322.461) (361.116,322.461) (361.259,322.461) (361.401,322.461) (361.543,322.461) (361.685,322.461) (361.827,322.461) (361.969,322.461) (362.111,322.461) (362.254,322.461) (362.396,322.461) (362.538,322.461) (362.680,322.461) (362.822,322.461) (362.964,322.461) (363.106,322.461) (363.249,322.461) (363.391,322.461) (363.533,322.461) (363.675,322.461) (363.817,322.461) (363.959,322.461) (364.101,322.461) (364.244,322.461) (364.386,322.461) (364.528,322.461) (364.670,322.461) (364.812,322.461) (364.954,322.461) (365.096,322.461) (365.239,322.461) (365.381,322.461) (365.523,322.461) (365.665,322.461) (365.807,322.461) (365.949,322.461) (366.091,322.461) (366.234,322.461) (366.376,322.461) (366.518,322.461) (366.660,322.461) (366.802,322.461) (366.944,322.461) (367.086,322.461) (367.229,322.461) (367.371,322.461) (367.513,322.461) (367.655,322.461) (367.797,322.461) (367.939,322.461) (368.081,322.461) (368.224,322.461) (368.366,322.461) (368.508,322.461) (368.650,322.461) (368.792,322.461) (368.934,322.461) (369.076,322.461) (369.219,322.461) (369.361,322.461) (369.503,322.461) (369.645,322.461) (369.787,322.461) (369.929,322.461) (370.071,322.461) (370.214,322.461) (370.356,322.461) (370.498,322.461) (370.640,322.461) (370.782,322.461) (370.924,322.461) (371.066,322.461) (371.209,322.461) (371.351,322.461) (371.493,322.461) (371.635,322.461) (371.777,322.461) (371.919,322.461) (372.061,322.461) (372.204,322.461) (372.346,322.461) (372.488,322.461) (372.630,322.461) (372.772,322.461) (372.914,322.461) (373.056,322.461) (373.199,322.461) (373.341,322.461) (373.483,322.461) (373.625,322.461) (373.767,322.461) (373.909,322.461) (374.051,322.461) (374.194,322.461) (374.336,322.461) (374.478,322.461) (374.620,322.461) (374.762,322.461) (374.904,322.461) (375.046,322.461) (375.189,322.461) (375.331,322.461) (375.473,322.461) (375.615,322.461) (375.757,322.461) (375.899,322.461) (376.041,322.461) (376.184,322.461) (376.326,322.461) (376.468,322.461) (376.610,322.461) (376.752,322.461) (376.894,322.461) (377.036,322.461) (377.179,322.461) (377.321,322.461) (377.463,322.461) (377.605,322.461) (377.747,322.461) (377.889,322.461) (378.031,322.461) (378.174,322.461) (378.316,322.461) (378.458,322.461) (378.600,322.461) (378.742,322.461) (378.884,322.461) (379.026,322.461) (379.169,322.461) (379.311,322.461) (379.453,322.461) (379.595,322.461) (379.737,322.461) (379.879,322.461) (380.021,332.126) (380.164,332.126) (380.306,332.126) (380.448,332.126) (380.590,332.126) (380.732,332.126) (380.874,332.126) (381.016,332.126) (381.159,332.126) (381.301,332.126) (381.443,332.126) (381.585,332.126) (381.727,332.126) (381.869,332.126) (382.011,332.126) (382.154,332.126) (382.296,332.126) (382.438,332.126) (382.580,332.126) (382.722,332.126) (382.864,332.126) (383.006,332.126) (383.149,332.126) (383.291,332.126) (383.433,332.126) (383.575,332.126) (383.717,332.126) (383.859,332.126) (384.001,332.126) (384.144,332.126) (384.286,332.126) (384.428,332.126) (384.570,332.126) (384.712,332.126) (384.854,332.126) (384.996,332.126) (385.139,332.126) (385.281,332.126) (385.423,332.126) (385.565,332.126) (385.707,332.126) (385.849,332.126) (385.991,332.126) (386.134,332.126) (386.276,332.126) (386.418,332.126) (386.560,332.126) (386.702,332.126) (386.844,332.126) (386.986,332.126) (387.129,332.126) (387.271,332.126) (387.413,332.126) (387.555,332.126) (387.697,332.126) (387.839,332.126) (387.981,332.126) (388.124,332.126) (388.266,332.126) (388.408,332.126) (388.550,332.126) (388.692,332.126) (388.834,332.126) (388.976,332.126) (389.119,332.126) (389.261,332.126) (389.403,332.126) (389.545,332.126) (389.687,332.126) (389.829,332.126) (389.971,332.126) (390.114,332.126) (390.256,332.126) (390.398,332.126) (390.540,332.126) (390.682,332.126) (390.824,332.126) (390.966,332.126) (391.109,332.126) (391.251,332.126) (391.393,332.126) (391.535,332.126) (391.677,332.126) (391.819,332.126) (391.961,332.126) (392.104,332.126) (392.246,332.126) (392.388,332.126) (392.530,332.126) (392.672,332.126) (392.814,332.126) (392.956,332.126) (393.099,332.126) (393.241,332.126) (393.383,332.126) (393.525,332.126) (393.667,332.126) (393.809,332.126) (393.951,332.126) (394.094,332.126) (394.236,332.126) (394.378,332.126) (394.520,332.126) (394.662,332.126) (394.804,332.126) (394.946,332.126) (395.089,332.126) (395.231,332.126) (395.373,332.126) (395.515,332.126) (395.657,332.126) (395.799,332.126) (395.941,332.126) (396.084,332.126) (396.226,332.126) (396.368,332.126) (396.510,332.126) (396.652,332.126) (396.794,332.126) (396.936,332.126) (397.079,332.126) (397.221,332.126) (397.363,332.126) (397.505,332.126) (397.647,332.126) (397.789,332.126) (397.931,332.126) (398.074,332.126) (398.216,332.126) (398.358,332.126) (398.500,332.126) (398.642,332.126) (398.784,332.126) (398.926,332.126) (399.069,332.126) (399.211,332.126) (399.353,332.126) (399.495,332.126) (399.637,332.126) (399.779,332.126) (399.921,332.126) (400.064,332.126) (400.206,332.126) (400.348,332.126) (400.490,332.126) (400.632,332.126) (400.774,332.126) (400.916,332.126) (401.059,332.126) (401.201,332.126) (401.343,332.126) (401.485,332.126) (401.627,332.126) (401.769,332.126) (401.911,332.126) (402.054,332.126) (402.196,332.126) (402.338,332.126) (402.480,332.126) (402.622,332.126) (402.764,332.126) (402.906,332.126) (403.049,332.126) (403.191,332.126) (403.333,332.126) (403.475,332.126) (403.617,332.126) (403.759,332.126) (403.901,332.126) (404.044,332.126) (404.186,332.126) (404.328,332.126) (404.470,332.126) (404.612,332.126) (404.754,332.126) (404.896,332.126) (405.039,332.126) (405.181,332.126) (405.323,332.126) (405.465,332.126) (405.607,332.126) (405.749,332.126) (405.891,332.126) (406.034,332.126) (406.176,332.126) (406.318,332.126) (406.460,332.126) (406.602,332.126) (406.744,332.126) (406.886,332.126) (407.029,332.126) (407.171,332.126) (407.313,332.126) (407.455,332.126) (407.597,332.126) (407.739,332.126) (407.881,332.126) (408.024,332.126) (408.166,332.126) (408.308,332.126) (408.450,332.126) (408.592,332.126) (408.734,332.126) (408.876,332.126) (409.019,332.126) (409.161,332.126) (409.303,332.126) (409.445,332.126) (409.587,332.126) (409.729,332.126) (409.871,332.126) (410.014,332.126) (410.156,332.126) (410.298,332.126) (410.440,332.126) (410.582,332.126) (410.724,332.126) (410.866,332.126) (411.009,332.126) (411.151,332.126) (411.293,332.126) (411.435,332.126) (411.577,332.126) (411.719,332.126) (411.861,332.126) (412.004,332.126) (412.146,332.126) (412.288,332.126) (412.430,332.126) (412.572,332.126) (412.714,332.126) (412.856,332.126) (412.999,332.126) (413.141,332.126) (413.283,332.126) (413.425,332.126) (413.567,332.126) (413.709,332.126) (413.851,332.126) (413.994,332.126) (414.136,332.126) (414.278,332.126) (414.420,332.126) (414.562,332.126) (414.704,332.126) (414.846,332.126) (414.989,332.126) (415.131,332.126) (415.273,332.126) (415.415,332.126) (415.557,332.126) (415.699,332.126) (415.841,332.126) (415.984,332.126) (416.126,332.126) (416.268,332.126) (416.410,332.126) (416.552,332.126) (416.694,332.126) (416.836,332.126) (416.979,332.126) (417.121,332.126) (417.263,332.126) (417.405,332.126) (417.547,332.126) (417.689,332.126) (417.831,332.126) (417.974,332.126) (418.116,332.126) (418.258,332.126) (418.400,332.126) (418.542,332.126) (418.684,332.126) (418.826,332.126) (418.969,332.126) (419.111,332.126) (419.253,332.126) (419.395,352.227) (419.537,352.227) (419.679,352.227) (419.821,352.227) (419.964,352.227) (420.106,352.227) (420.248,352.227) (420.390,352.227) (420.532,352.227) (420.674,352.227) (420.816,352.227) (420.959,352.227) (421.101,352.227) (421.243,352.227) (421.385,352.227) (421.527,352.227) (421.669,352.227) (421.811,352.227) (421.954,352.227) (422.096,352.227) (422.238,352.227) (422.380,352.227) (422.522,352.227) (422.664,352.227) (422.806,352.227) (422.949,352.227) (423.091,352.227) (423.233,352.227) (423.375,352.227) (423.517,352.227) (423.659,352.227) (423.801,352.227) (423.944,352.227) (424.086,352.227) (424.228,352.227) (424.370,352.227) (424.512,352.227) (424.654,352.227) (424.796,352.227) (424.939,352.227) (425.081,352.227) (425.223,352.227) (425.365,352.227) (425.507,352.227) (425.649,352.227) (425.791,352.227) (425.934,352.227) (426.076,352.227) (426.218,352.227) (426.360,352.227) (426.502,352.227) (426.644,352.227) (426.786,352.227) (426.929,367.492) (427.071,367.492) (427.213,367.492) (427.355,367.492) (427.497,367.492) (427.639,367.492) (427.781,367.492) (427.924,367.492) (428.066,367.492) (428.208,367.492) (428.350,367.492) (428.492,367.492) (428.634,367.492) (428.776,367.492) (428.919,367.492) (429.061,367.492) (429.203,367.492) (429.345,367.492) (429.487,367.492) (429.629,367.492) (429.771,367.492) (429.914,367.492) (430.056,367.492) (430.198,367.492) (430.340,367.492) (430.482,367.492) (430.624,367.492) (430.766,367.492) (430.909,367.492) (431.051,367.492) (431.193,367.492) (431.335,367.492) (431.477,367.492) (431.619,367.492) (431.761,367.492) (431.904,367.492) (432.046,367.492) (432.188,367.492) (432.330,367.492) (432.472,367.492) (432.614,367.492) (432.756,367.492) (432.899,367.492) (433.041,367.492) (433.183,367.492) (433.325,367.492) (433.467,367.492) (433.609,367.492) (433.751,367.492) (433.894,367.492) (434.036,367.492) (434.178,367.492) (434.320,367.492) (434.462,367.492) (434.604,367.492) (434.746,367.492) (434.889,367.492) (435.031,367.492) (435.173,367.492) (435.315,367.492) (435.457,367.492) (435.599,367.492) (435.741,367.492) (435.884,367.492) (436.026,367.492) (436.168,367.492) (436.310,367.492) (436.452,367.492) (436.594,367.492) (436.736,367.492) (436.879,367.492) (437.021,367.492) (437.163,367.492) (437.305,367.492) (437.447,367.492) (437.589,367.492) (437.731,367.492) (437.874,367.492) (438.016,367.492) (438.158,367.492) (438.300,367.492) (438.442,367.492) (438.584,367.492) (438.726,367.492) (438.869,367.492) (439.011,367.492) (439.153,367.492) (439.295,367.492) (439.437,367.492) (439.579,367.492) (439.721,367.492) (439.864,367.492) (440.006,367.492) (440.148,367.492) (440.290,367.492) (440.432,367.492) (440.574,367.492) (440.716,367.492) (440.859,367.492) (441.001,367.492) (441.143,367.492) (441.285,367.492) (441.427,367.492) (441.569,367.492) (441.711,367.492) (441.854,367.492) (441.996,367.492) (442.138,367.492) (442.280,367.492) (442.422,367.492) (442.564,367.492) (442.706,367.492) (442.849,367.492) (442.991,367.492) (443.133,367.492) (443.275,367.492) (443.417,367.492) (443.559,367.492) (443.701,367.492) (443.844,367.492) (443.986,367.492) (444.128,367.492) (444.270,367.492) (444.412,367.492) (444.554,367.492) (444.696,367.492) (444.839,367.492) (444.981,367.492) (445.123,367.492) (445.265,367.492) (445.407,367.492) (445.549,367.492) (445.691,367.492) (445.834,367.492) (445.976,367.492) (446.118,367.492) (446.260,367.492) (446.402,367.492) (446.544,367.492) (446.686,367.492) (446.829,367.492) (446.971,367.492) (447.113,367.492) (447.255,367.492) (447.397,367.492) (447.539,367.492) (447.681,367.492) (447.824,367.492) (447.966,367.492) (448.108,367.492) (448.250,367.492) (448.392,367.492) (448.534,367.492) (448.676,367.492) (448.819,367.492) (448.961,367.492) (449.103,367.492) (449.245,367.492) (449.387,367.492) (449.529,367.492) (449.671,367.492) (449.814,367.492) (449.956,367.492) (450.098,367.492) (450.240,367.492) (450.382,367.492) (450.524,367.492) (450.666,367.492) (450.809,367.492) (450.951,367.492) (451.093,367.492) (451.235,367.492) (451.377,367.492) (451.519,367.492) (451.661,367.492) (451.804,367.492) (451.946,367.492) (452.088,367.492) (452.230,367.492) (452.372,367.492) (452.514,367.492) (452.656,367.492) (452.799,367.492) (452.941,367.492) (453.083,367.492) (453.225,367.492) (453.367,367.492) (453.509,367.492) (453.651,367.492) (453.794,367.492) (453.936,367.492) (454.078,367.492) (454.220,367.492) (454.362,367.492) (454.504,367.492) (454.646,367.492) (454.789,367.492) (454.931,367.492) (455.073,367.492) (455.215,367.492) (455.357,367.492) (455.499,367.492) (455.641,367.492) (455.784,367.492) (455.926,367.492) (456.068,367.492) (456.210,367.492) (456.352,367.492) (456.494,367.492) (456.636,367.492) (456.779,367.492) (456.921,367.492) (457.063,367.492) (457.205,367.492) (457.347,367.492) (457.489,367.492) (457.631,367.492) (457.774,367.492) (457.916,367.492) (458.058,367.492) (458.200,367.492) (458.342,367.492) (458.484,367.492) (458.626,367.492) (458.769,367.492) (458.911,367.492) (459.053,367.492) (459.195,367.492) (459.337,367.492) (459.479,367.492) (459.621,367.492) (459.764,367.492) (459.906,367.492) (460.048,367.492) (460.190,367.492) (460.332,367.492) (460.474,367.492) (460.616,367.492) (460.759,367.492) (460.901,367.492) (461.043,367.492) (461.185,367.492) (461.327,367.492) (461.469,367.492) (461.611,367.492) (461.754,367.492) (461.896,367.492) (462.038,367.492) (462.180,367.492) (462.322,367.492) (462.464,367.492) (462.606,367.492) (462.749,367.492) (462.891,367.492) (463.033,367.492) (463.175,367.492) (463.317,367.492) (463.459,367.492) (463.601,367.492) (463.744,367.492) (463.886,367.492) (464.028,367.492) (464.170,367.492) (464.312,367.492) (464.454,367.492) (464.596,367.492) (464.739,367.492) (464.881,367.492) (465.023,367.492) (465.165,367.492) (465.307,367.492) (465.449,367.492) (465.591,367.492) (465.734,367.492) (465.876,367.492) (466.018,367.492) (466.160,367.492) (466.302,367.492) (466.444,367.492) (466.586,367.492) (466.729,367.492) (466.871,367.492) (467.013,367.492) (467.155,396.247) (467.297,396.247) (467.439,396.247) (467.581,396.247) (467.724,396.247) (467.866,396.247) (468.008,396.247) (468.150,396.247) (468.292,396.247) (468.434,396.247) (468.576,396.247) (468.719,396.247) (468.861,396.247) (469.003,396.247) (469.145,396.247) (469.287,396.247) (469.429,396.247) (469.571,396.247) (469.714,396.247) (469.856,396.247) (469.998,396.247) (470.140,396.247) (470.282,396.247) (470.424,396.247) (470.566,396.247) (470.709,396.247) (470.851,396.247) (470.993,396.247) (471.135,396.247) (471.277,396.247) (471.419,396.247) (471.561,396.247) (471.704,396.247) (471.846,396.247) (471.988,396.247) (472.130,396.247) (472.272,396.247) (472.414,396.247) (472.556,396.247) (472.699,396.247) (472.841,396.247) (472.983,396.247) (473.125,396.247) (473.267,396.247) (473.409,396.247) (473.551,396.247) (473.694,396.247) (473.836,396.247) (473.978,396.247) (474.120,396.247) (474.262,396.247) (474.404,396.247) (474.546,396.247) (474.689,396.247) (474.831,396.247) (474.973,396.247) (475.115,396.247) (475.257,396.247) (475.399,396.247) (475.541,396.247) (475.684,396.247) (475.826,396.247) (475.968,396.247) (476.110,396.247) (476.252,396.247) (476.394,396.247) (476.536,396.247) (476.679,396.247) (476.821,396.247) (476.963,396.247) (477.105,396.247) (477.247,396.247) (477.389,396.247) (477.531,396.247) (477.674,396.247) (477.816,396.247) (477.958,396.247) (478.100,396.247) (478.242,396.247) (478.384,396.247) (478.526,396.247) (478.669,396.247) (478.811,396.247) (478.953,396.247) (479.095,396.247) (479.237,396.247) (479.379,396.247) (479.521,396.247) (479.664,396.247) (479.806,396.247) (479.948,396.247) (480.090,396.247) (480.232,396.247) (480.374,396.247) (480.516,396.247) (480.659,396.247) (480.801,396.247) (480.943,396.247) (481.085,396.247) (481.227,396.247) (481.369,396.247) (481.511,396.247) (481.654,396.247) (481.796,396.247) (481.938,396.247) (482.080,396.247) (482.222,396.247) (482.364,396.247) (482.506,396.247) (482.649,396.247) (482.791,396.247) (482.933,396.247) (483.075,396.247) (483.217,396.247) (483.359,396.247) (483.501,396.247) (483.644,396.247) (483.786,396.247) (483.928,396.247) (484.070,396.247) (484.212,396.247) (484.354,396.247) (484.496,396.247) (484.639,396.247) (484.781,396.247) (484.923,396.247) (485.065,396.247) (485.207,396.247) (485.349,396.247) (485.491,396.247) (485.634,396.247) (485.776,396.247) (485.918,396.247) (486.060,396.247) (486.202,396.247) (486.344,396.247) (486.486,396.247) (486.629,396.247) (486.771,396.247) (486.913,396.247) (487.055,396.247) (487.197,396.247) (487.339,396.247) (487.481,396.247) (487.624,396.247) (487.766,396.247) (487.908,396.247) (488.050,396.247) (488.192,396.247) (488.334,396.247) (488.476,396.247) (488.619,396.247) (488.761,396.247) (488.903,396.247) (489.045,396.247) (489.187,396.247) (489.329,396.247) (489.471,396.247) (489.614,396.247) (489.756,396.247) (489.898,396.247) (490.040,396.247) (490.182,396.247) (490.324,396.247) (490.466,396.247) (490.609,396.247) (490.751,396.247) (490.893,396.247) (491.035,396.247) (491.177,396.247) (491.319,396.247) (491.461,396.247) (491.604,396.247) (491.746,396.247) (491.888,396.247) (492.030,396.247) (492.172,396.247) (492.314,396.247) (492.456,396.247) (492.599,396.247) (492.741,396.247) (492.883,396.247) (493.025,396.247) (493.167,396.247) (493.309,396.247) (493.451,396.247) (493.594,396.247) (493.736,396.247) (493.878,396.247) (494.020,396.247) (494.162,396.247) (494.304,396.247) (494.446,396.247) (494.589,396.247) (494.731,396.247) (494.873,396.247) (495.015,396.247) (495.157,396.247) (495.299,396.247) (495.441,396.247) (495.584,396.247) (495.726,396.247) (495.868,396.247) (496.010,396.247) (496.152,396.247) (496.294,396.247) (496.436,396.247) (496.579,396.247) (496.721,396.247) (496.863,396.247) (497.005,396.247) (497.147,396.247) (497.289,396.247) (497.431,396.247) (497.574,396.247) (497.716,396.247) (497.858,396.247) (498.000,396.247)};
\draw[methodA,line width=1.05pt,dashed] plot coordinates {(299.000,320.193) (299.142,320.142) (299.284,320.092) (299.426,320.042) (299.569,319.992) (299.711,319.943) (299.853,319.893) (299.995,319.844) (300.137,319.795) (300.279,319.747) (300.421,319.699) (300.564,319.651) (300.706,319.603) (300.848,319.555) (300.990,319.508) (301.132,319.461) (301.274,319.414) (301.416,319.368) (301.559,319.322) (301.701,319.275) (301.843,319.230) (301.985,319.184) (302.127,319.139) (302.269,319.094) (302.411,319.049) (302.554,319.005) (302.696,318.960) (302.838,318.916) (302.980,318.873) (303.122,318.829) (303.264,318.786) (303.406,318.743) (303.549,318.700) (303.691,318.657) (303.833,318.615) (303.975,318.573) (304.117,318.531) (304.259,318.490) (304.401,318.448) (304.544,318.407) (304.686,318.367) (304.828,318.326) (304.970,318.286) (305.112,318.246) (305.254,318.206) (305.396,318.166) (305.539,318.127) (305.681,318.088) (305.823,318.049) (305.965,318.010) (306.107,317.972) (306.249,317.934) (306.391,317.896) (306.534,317.858) (306.676,317.821) (306.818,317.784) (306.960,317.747) (307.102,317.710) (307.244,317.674) (307.386,317.638) (307.529,317.602) (307.671,317.566) (307.813,317.531) (307.955,317.495) (308.097,317.460) (308.239,317.426) (308.381,317.391) (308.524,317.357) (308.666,317.323) (308.808,317.289) (308.950,317.256) (309.092,317.222) (309.234,317.189) (309.376,317.156) (309.519,317.124) (309.661,317.092) (309.803,317.059) (309.945,317.028) (310.087,316.996) (310.229,316.965) (310.371,316.933) (310.514,316.903) (310.656,316.872) (310.798,316.841) (310.940,316.811) (311.082,316.781) (311.224,316.752) (311.366,316.722) (311.509,316.693) (311.651,316.664) (311.793,316.635) (311.935,316.607) (312.077,316.578) (312.219,316.550) (312.361,316.522) (312.504,316.495) (312.646,316.467) (312.788,316.440) (312.930,316.413) (313.072,316.387) (313.214,316.360) (313.356,316.334) (313.499,316.308) (313.641,316.282) (313.783,316.257) (313.925,316.232) (314.067,316.207) (314.209,316.182) (314.351,316.157) (314.494,316.133) (314.636,316.109) (314.778,316.085) (314.920,316.061) (315.062,316.038) (315.204,316.015) (315.346,315.992) (315.489,315.969) (315.631,315.946) (315.773,315.924) (315.915,315.902) (316.057,315.880) (316.199,315.859) (316.341,315.837) (316.484,315.816) (316.626,315.795) (316.768,315.775) (316.910,315.754) (317.052,315.734) (317.194,315.714) (317.336,315.694) (317.479,315.675) (317.621,315.656) (317.763,315.637) (317.905,315.618) (318.047,315.599) (318.189,315.581) (318.331,315.563) (318.474,315.545) (318.616,315.527) (318.758,315.510) (318.900,315.492) (319.042,315.475) (319.184,315.458) (319.326,315.442) (319.469,315.426) (319.611,315.409) (319.753,315.394) (319.895,315.378) (320.037,315.362) (320.179,315.347) (320.321,315.332) (320.464,315.317) (320.606,315.303) (320.748,315.288) (320.890,315.274) (321.032,315.260) (321.174,315.247) (321.316,315.233) (321.459,315.220) (321.601,315.207) (321.743,315.194) (321.885,315.182) (322.027,315.169) (322.169,315.157) (322.311,315.145) (322.454,315.134) (322.596,315.122) (322.738,315.111) (322.880,315.100) (323.022,315.089) (323.164,315.079) (323.306,315.068) (323.449,315.058) (323.591,315.048) (323.733,315.039) (323.875,315.029) (324.017,315.020) (324.159,315.011) (324.301,315.002) (324.444,314.993) (324.586,314.985) (324.728,314.977) (324.870,314.969) (325.012,314.961) (325.154,314.953) (325.296,314.946) (325.439,314.939) (325.581,314.932) (325.723,314.925) (325.865,314.919) (326.007,314.913) (326.149,314.907) (326.291,314.901) (326.434,314.895) (326.576,314.890) (326.718,314.885) (326.860,314.880) (327.002,314.875) (327.144,314.871) (327.286,314.866) (327.429,314.862) (327.571,314.858) (327.713,314.855) (327.855,314.851) (327.997,314.848) (328.139,314.845) (328.281,314.842) (328.424,314.840) (328.566,314.837) (328.708,314.835) (328.850,314.833) (328.992,314.831) (329.134,314.830) (329.276,314.828) (329.419,314.827) (329.561,314.826) (329.703,314.826) (329.845,314.825) (329.987,314.825) (330.129,314.825) (330.271,314.825) (330.414,314.825) (330.556,314.826) (330.698,314.826) (330.840,314.827) (330.982,314.829) (331.124,314.830) (331.266,314.832) (331.409,314.833) (331.551,314.835) (331.693,314.838) (331.835,314.840) (331.977,314.843) (332.119,314.845) (332.261,314.848) (332.404,314.852) (332.546,314.855) (332.688,314.859) (332.830,314.863) (332.972,314.867) (333.114,314.871) (333.256,314.875) (333.399,314.880) (333.541,314.885) (333.683,314.890) (333.825,314.895) (333.967,314.901) (334.109,314.907) (334.251,314.913) (334.394,314.919) (334.536,314.925) (334.678,314.931) (334.820,314.938) (334.962,314.945) (335.104,314.952) (335.246,314.960) (335.389,314.967) (335.531,314.975) (335.673,314.983) (335.815,314.991) (335.957,314.999) (336.099,315.008) (336.241,315.016) (336.384,315.025) (336.526,315.034) (336.668,315.044) (336.810,315.053) (336.952,315.063) (337.094,315.073) (337.236,315.083) (337.379,315.093) (337.521,315.104) (337.663,315.115) (337.805,315.126) (337.947,315.137) (338.089,315.148) (338.231,315.160) (338.374,315.171) (338.516,315.183) (338.658,315.195) (338.800,315.208) (338.942,315.220) (339.084,315.233) (339.226,315.246) (339.369,315.259) (339.511,315.272) (339.653,315.286) (339.795,315.299) (339.937,315.313) (340.079,315.327) (340.221,315.341) (340.364,315.356) (340.506,315.371) (340.648,315.385) (340.790,315.400) (340.932,315.416) (341.074,315.431) (341.216,315.447) (341.359,315.463) (341.501,315.479) (341.643,315.495) (341.785,315.511) (341.927,315.528) (342.069,315.545) (342.211,315.561) (342.354,315.579) (342.496,315.596) (342.638,315.614) (342.780,315.631) (342.922,315.649) (343.064,315.667) (343.206,315.686) (343.349,315.704) (343.491,315.723) (343.633,315.742) (343.775,315.761) (343.917,315.780) (344.059,315.799) (344.201,315.819) (344.344,315.839) (344.486,315.859) (344.628,315.879) (344.770,315.900) (344.912,315.920) (345.054,315.941) (345.196,315.962) (345.339,315.983) (345.481,316.004) (345.623,316.026) (345.765,316.048) (345.907,316.069) (346.049,316.092) (346.191,316.114) (346.334,316.136) (346.476,316.159) (346.618,316.182) (346.760,316.205) (346.902,316.228) (347.044,316.251) (347.186,316.275) (347.329,316.299) (347.471,316.323) (347.613,316.347) (347.755,316.371) (347.897,316.395) (348.039,316.420) (348.181,316.445) (348.324,316.470) (348.466,316.495) (348.608,316.521) (348.750,316.546) (348.892,316.572) (349.034,316.598) (349.176,316.624) (349.319,316.650) (349.461,316.677) (349.603,316.704) (349.745,316.730) (349.887,316.757) (350.029,316.785) (350.171,316.812) (350.314,316.840) (350.456,316.867) (350.598,316.895) (350.740,316.923) (350.882,316.952) (351.024,316.980) (351.166,317.009) (351.309,317.038) (351.451,317.067) (351.593,317.096) (351.735,317.125) (351.877,317.155) (352.019,317.185) (352.161,317.215) (352.304,317.245) (352.446,317.275) (352.588,317.305) (352.730,317.336) (352.872,317.367) (353.014,317.398) (353.156,317.429) (353.299,317.460) (353.441,317.492) (353.583,317.523) (353.725,317.555) (353.867,317.587) (354.009,317.619) (354.151,317.652) (354.294,317.684) (354.436,317.717) (354.578,317.750) (354.720,317.783) (354.862,317.816) (355.004,317.850) (355.146,317.883) (355.289,317.917) (355.431,317.951) (355.573,317.985) (355.715,318.019) (355.857,318.054) (355.999,318.088) (356.141,318.123) (356.284,318.158) (356.426,318.193) (356.568,318.229) (356.710,318.264) (356.852,318.300) (356.994,318.336) (357.136,318.372) (357.279,318.408) (357.421,318.444) (357.563,318.481) (357.705,318.517) (357.847,318.554) (357.989,318.591) (358.131,318.628) (358.274,318.666) (358.416,318.703) (358.558,318.741) (358.700,318.779) (358.842,318.817) (358.984,318.855) (359.126,318.893) (359.269,318.932) (359.411,318.971) (359.553,319.010) (359.695,319.049) (359.837,319.088) (359.979,319.127) (360.121,319.167) (360.264,319.206) (360.406,319.246) (360.548,319.286) (360.690,319.326) (360.832,319.367) (360.974,319.407) (361.116,319.448) (361.259,319.489) (361.401,319.530) (361.543,319.571) (361.685,319.613) (361.827,319.654) (361.969,319.696) (362.111,319.738) (362.254,319.780) (362.396,319.822) (362.538,319.864) (362.680,319.907) (362.822,319.949) (362.964,319.992) (363.106,320.035) (363.249,320.078) (363.391,320.121) (363.533,320.165) (363.675,320.209) (363.817,320.252) (363.959,320.296) (364.101,320.340) (364.244,320.385) (364.386,320.429) (364.528,320.474) (364.670,320.518) (364.812,320.563) (364.954,320.608) (365.096,320.654) (365.239,320.699) (365.381,320.745) (365.523,320.790) (365.665,320.836) (365.807,320.882) (365.949,320.928) (366.091,320.975) (366.234,321.021) (366.376,321.068) (366.518,321.115) (366.660,321.162) (366.802,321.209) (366.944,321.256) (367.086,321.303) (367.229,321.351) (367.371,321.399) (367.513,321.447) (367.655,321.495) (367.797,321.543) (367.939,321.591) (368.081,321.640) (368.224,321.689) (368.366,321.737) (368.508,321.786) (368.650,321.836) (368.792,321.885) (368.934,321.934) (369.076,321.984) (369.219,322.034) (369.361,322.084) (369.503,322.134) (369.645,322.184) (369.787,322.234) (369.929,322.285) (370.071,322.335) (370.214,322.386) (370.356,322.437) (370.498,322.488) (370.640,322.540) (370.782,322.591) (370.924,322.643) (371.066,322.694) (371.209,322.746) (371.351,322.798) (371.493,322.850) (371.635,322.903) (371.777,322.955) (371.919,323.008) (372.061,323.061) (372.204,323.114) (372.346,323.167) (372.488,323.220) (372.630,323.273) (372.772,323.327) (372.914,323.380) (373.056,323.434) (373.199,323.488) (373.341,323.542) (373.483,323.597) (373.625,323.651) (373.767,323.706) (373.909,323.760) (374.051,323.815) (374.194,323.870) (374.336,323.925) (374.478,323.981) (374.620,324.036) (374.762,324.092) (374.904,324.147) (375.046,324.203) (375.189,324.259) (375.331,324.315) (375.473,324.372) (375.615,324.428) (375.757,324.485) (375.899,324.541) (376.041,324.598) (376.184,324.655) (376.326,324.712) (376.468,324.770) (376.610,324.827) (376.752,324.885) (376.894,324.942) (377.036,325.000) (377.179,325.058) (377.321,325.116) (377.463,325.175) (377.605,325.233) (377.747,325.292) (377.889,325.351) (378.031,325.409) (378.174,325.468) (378.316,325.528) (378.458,325.587) (378.600,325.646) (378.742,325.706) (378.884,325.766) (379.026,325.825) (379.169,325.885) (379.311,325.945) (379.453,326.006) (379.595,326.066) (379.737,326.127) (379.879,326.187) (380.021,326.248) (380.164,326.309) (380.306,326.370) (380.448,326.431) (380.590,326.493) (380.732,326.554) (380.874,326.616) (381.016,326.678) (381.159,326.740) (381.301,326.802) (381.443,326.864) (381.585,326.926) (381.727,326.988) (381.869,327.051) (382.011,327.114) (382.154,327.177) (382.296,327.240) (382.438,327.303) (382.580,327.366) (382.722,327.429) (382.864,327.493) (383.006,327.557) (383.149,327.620) (383.291,327.684) (383.433,327.748) (383.575,327.813) (383.717,327.877) (383.859,327.941) (384.001,328.006) (384.144,328.071) (384.286,328.135) (384.428,328.200) (384.570,328.266) (384.712,328.331) (384.854,328.396) (384.996,328.462) (385.139,328.527) (385.281,328.593) (385.423,328.659) (385.565,328.725) (385.707,328.791) (385.849,328.857) (385.991,328.924) (386.134,328.990) (386.276,329.057) (386.418,329.124) (386.560,329.191) (386.702,329.258) (386.844,329.325) (386.986,329.392) (387.129,329.460) (387.271,329.527) (387.413,329.595) (387.555,329.663) (387.697,329.731) (387.839,329.799) (387.981,329.867) (388.124,329.935) (388.266,330.004) (388.408,330.072) (388.550,330.141) (388.692,330.210) (388.834,330.279) (388.976,330.348) (389.119,330.417) (389.261,330.486) (389.403,330.556) (389.545,330.625) (389.687,330.695) (389.829,330.765) (389.971,330.835) (390.114,330.905) (390.256,330.975) (390.398,331.046) (390.540,331.116) (390.682,331.187) (390.824,331.257) (390.966,331.328) (391.109,331.399) (391.251,331.470) (391.393,331.541) (391.535,331.612) (391.677,331.684) (391.819,331.755) (391.961,331.827) (392.104,331.899) (392.246,331.971) (392.388,332.043) (392.530,332.115) (392.672,332.187) (392.814,332.259) (392.956,332.332) (393.099,332.405) (393.241,332.477) (393.383,332.550) (393.525,332.623) (393.667,332.696) (393.809,332.769) (393.951,332.843) (394.094,332.916) (394.236,332.990) (394.378,333.063) (394.520,333.137) (394.662,333.211) (394.804,333.285) (394.946,333.359) (395.089,333.434) (395.231,333.508) (395.373,333.582) (395.515,333.657) (395.657,333.732) (395.799,333.807) (395.941,333.881) (396.084,333.957) (396.226,334.032) (396.368,334.107) (396.510,334.182) (396.652,334.258) (396.794,334.334) (396.936,334.409) (397.079,334.485) (397.221,334.561) (397.363,334.637) (397.505,334.713) (397.647,334.790) (397.789,334.866) (397.931,334.943) (398.074,335.019) (398.216,335.096) (398.358,335.173) (398.500,335.250) (398.642,335.327) (398.784,335.404) (398.926,335.482) (399.069,335.559) (399.211,335.637) (399.353,335.714) (399.495,335.792) (399.637,335.870) (399.779,335.948) (399.921,336.026) (400.064,336.104) (400.206,336.182) (400.348,336.261) (400.490,336.339) (400.632,336.418) (400.774,336.497) (400.916,336.576) (401.059,336.654) (401.201,336.734) (401.343,336.813) (401.485,336.892) (401.627,336.971) (401.769,337.051) (401.911,337.130) (402.054,337.210) (402.196,337.290) (402.338,337.370) (402.480,337.450) (402.622,337.530) (402.764,337.610) (402.906,337.690) (403.049,337.771) (403.191,337.851) (403.333,337.932) (403.475,338.013) (403.617,338.094) (403.759,338.175) (403.901,338.256) (404.044,338.337) (404.186,338.418) (404.328,338.500) (404.470,338.581) (404.612,338.663) (404.754,338.744) (404.896,338.826) (405.039,338.908) (405.181,338.990) (405.323,339.072) (405.465,339.154) (405.607,339.236) (405.749,339.319) (405.891,339.401) (406.034,339.484) (406.176,339.567) (406.318,339.649) (406.460,339.732) (406.602,339.815) (406.744,339.898) (406.886,339.982) (407.029,340.065) (407.171,340.148) (407.313,340.232) (407.455,340.315) (407.597,340.399) (407.739,340.483) (407.881,340.567) (408.024,340.651) (408.166,340.735) (408.308,340.819) (408.450,340.903) (408.592,340.987) (408.734,341.072) (408.876,341.156) (409.019,341.241) (409.161,341.326) (409.303,341.411) (409.445,341.496) (409.587,341.581) (409.729,341.666) (409.871,341.751) (410.014,341.836) (410.156,341.922) (410.298,342.007) (410.440,342.093) (410.582,342.178) (410.724,342.264) (410.866,342.350) (411.009,342.436) (411.151,342.522) (411.293,342.608) (411.435,342.695) (411.577,342.781) (411.719,342.867) (411.861,342.954) (412.004,343.041) (412.146,343.127) (412.288,343.214) (412.430,343.301) (412.572,343.388) (412.714,343.475) (412.856,343.562) (412.999,343.649) (413.141,343.737) (413.283,343.824) (413.425,343.912) (413.567,343.999) (413.709,344.087) (413.851,344.175) (413.994,344.263) (414.136,344.351) (414.278,344.439) (414.420,344.527) (414.562,344.615) (414.704,344.703) (414.846,344.792) (414.989,344.880) (415.131,344.969) (415.273,345.058) (415.415,345.146) (415.557,345.235) (415.699,345.324) (415.841,345.413) (415.984,345.502) (416.126,345.592) (416.268,345.681) (416.410,345.770) (416.552,345.860) (416.694,345.949) (416.836,346.039) (416.979,346.129) (417.121,346.218) (417.263,346.308) (417.405,346.398) (417.547,346.488) (417.689,346.578) (417.831,346.669) (417.974,346.759) (418.116,346.849) (418.258,346.940) (418.400,347.030) (418.542,347.121) (418.684,347.212) (418.826,347.302) (418.969,347.393) (419.111,347.484) (419.253,347.575) (419.395,347.666) (419.537,347.758) (419.679,347.849) (419.821,347.940) (419.964,348.032) (420.106,348.123) (420.248,348.215) (420.390,348.306) (420.532,348.398) (420.674,348.490) (420.816,348.582) (420.959,348.674) (421.101,348.766) (421.243,348.858) (421.385,348.950) (421.527,349.043) (421.669,349.135) (421.811,349.228) (421.954,349.320) (422.096,349.413) (422.238,349.505) (422.380,349.598) (422.522,349.691) (422.664,349.784) (422.806,349.877) (422.949,349.970) (423.091,350.063) (423.233,350.156) (423.375,350.250) (423.517,350.343) (423.659,350.437) (423.801,350.530) (423.944,350.624) (424.086,350.717) (424.228,350.811) (424.370,350.905) (424.512,350.999) (424.654,351.093) (424.796,351.187) (424.939,351.281) (425.081,351.375) (425.223,351.470) (425.365,351.564) (425.507,351.658) (425.649,351.753) (425.791,351.847) (425.934,351.942) (426.076,352.037) (426.218,352.132) (426.360,352.226) (426.502,352.321) (426.644,352.416) (426.786,352.511) (426.929,352.606) (427.071,352.702) (427.213,352.797) (427.355,352.892) (427.497,352.988) (427.639,353.083) (427.781,353.179) (427.924,353.274) (428.066,353.370) (428.208,353.466) (428.350,353.562) (428.492,353.658) (428.634,353.753) (428.776,353.849) (428.919,353.946) (429.061,354.042) (429.203,354.138) (429.345,354.234) (429.487,354.331) (429.629,354.427) (429.771,354.524) (429.914,354.620) (430.056,354.717) (430.198,354.813) (430.340,354.910) (430.482,355.007) (430.624,355.104) (430.766,355.201) (430.909,355.298) (431.051,355.395) (431.193,355.492) (431.335,355.589) (431.477,355.687) (431.619,355.784) (431.761,355.881) (431.904,355.979) (432.046,356.076) (432.188,356.174) (432.330,356.272) (432.472,356.369) (432.614,356.467) (432.756,356.565) (432.899,356.663) (433.041,356.761) (433.183,356.859) (433.325,356.957) (433.467,357.055) (433.609,357.153) (433.751,357.252) (433.894,357.350) (434.036,357.448) (434.178,357.547) (434.320,357.645) (434.462,357.744) (434.604,357.843) (434.746,357.941) (434.889,358.040) (435.031,358.139) (435.173,358.238) (435.315,358.337) (435.457,358.436) (435.599,358.535) (435.741,358.634) (435.884,358.733) (436.026,358.832) (436.168,358.931) (436.310,359.031) (436.452,359.130) (436.594,359.230) (436.736,359.329) (436.879,359.429) (437.021,359.528) (437.163,359.628) (437.305,359.728) (437.447,359.827) (437.589,359.927) (437.731,360.027) (437.874,360.127) (438.016,360.227) (438.158,360.327) (438.300,360.427) (438.442,360.527) (438.584,360.628) (438.726,360.728) (438.869,360.828) (439.011,360.929) (439.153,361.029) (439.295,361.129) (439.437,361.230) (439.579,361.331) (439.721,361.431) (439.864,361.532) (440.006,361.633) (440.148,361.733) (440.290,361.834) (440.432,361.935) (440.574,362.036) (440.716,362.137) (440.859,362.238) (441.001,362.339) (441.143,362.440) (441.285,362.542) (441.427,362.643) (441.569,362.744) (441.711,362.845) (441.854,362.947) (441.996,363.048) (442.138,363.150) (442.280,363.251) (442.422,363.353) (442.564,363.455) (442.706,363.556) (442.849,363.658) (442.991,363.760) (443.133,363.862) (443.275,363.963) (443.417,364.065) (443.559,364.167) (443.701,364.269) (443.844,364.371) (443.986,364.473) (444.128,364.576) (444.270,364.678) (444.412,364.780) (444.554,364.882) (444.696,364.985) (444.839,365.087) (444.981,365.189) (445.123,365.292) (445.265,365.394) (445.407,365.497) (445.549,365.600) (445.691,365.702) (445.834,365.805) (445.976,365.908) (446.118,366.010) (446.260,366.113) (446.402,366.216) (446.544,366.319) (446.686,366.422) (446.829,366.525) (446.971,366.628) (447.113,366.731) (447.255,366.834) (447.397,366.937) (447.539,367.040) (447.681,367.144) (447.824,367.247) (447.966,367.350) (448.108,367.454) (448.250,367.557) (448.392,367.660) (448.534,367.764) (448.676,367.867) (448.819,367.971) (448.961,368.075) (449.103,368.178) (449.245,368.282) (449.387,368.386) (449.529,368.489) (449.671,368.593) (449.814,368.697) (449.956,368.801) (450.098,368.905) (450.240,369.008) (450.382,369.112) (450.524,369.216) (450.666,369.320) (450.809,369.425) (450.951,369.529) (451.093,369.633) (451.235,369.737) (451.377,369.841) (451.519,369.945) (451.661,370.050) (451.804,370.154) (451.946,370.258) (452.088,370.363) (452.230,370.467) (452.372,370.572) (452.514,370.676) (452.656,370.781) (452.799,370.885) (452.941,370.990) (453.083,371.094) (453.225,371.199) (453.367,371.304) (453.509,371.409) (453.651,371.513) (453.794,371.618) (453.936,371.723) (454.078,371.828) (454.220,371.933) (454.362,372.037) (454.504,372.142) (454.646,372.247) (454.789,372.352) (454.931,372.457) (455.073,372.562) (455.215,372.668) (455.357,372.773) (455.499,372.878) (455.641,372.983) (455.784,373.088) (455.926,373.193) (456.068,373.299) (456.210,373.404) (456.352,373.509) (456.494,373.615) (456.636,373.720) (456.779,373.825) (456.921,373.931) (457.063,374.036) (457.205,374.142) (457.347,374.247) (457.489,374.353) (457.631,374.459) (457.774,374.564) (457.916,374.670) (458.058,374.775) (458.200,374.881) (458.342,374.987) (458.484,375.092) (458.626,375.198) (458.769,375.304) (458.911,375.410) (459.053,375.516) (459.195,375.621) (459.337,375.727) (459.479,375.833) (459.621,375.939) (459.764,376.045) (459.906,376.151) (460.048,376.257) (460.190,376.363) (460.332,376.469) (460.474,376.575) (460.616,376.681) (460.759,376.787) (460.901,376.893) (461.043,377.000) (461.185,377.106) (461.327,377.212) (461.469,377.318) (461.611,377.424) (461.754,377.531) (461.896,377.637) (462.038,377.743) (462.180,377.849) (462.322,377.956) (462.464,378.062) (462.606,378.168) (462.749,378.275) (462.891,378.381) (463.033,378.488) (463.175,378.594) (463.317,378.700) (463.459,378.807) (463.601,378.913) (463.744,379.020) (463.886,379.126) (464.028,379.233) (464.170,379.340) (464.312,379.446) (464.454,379.553) (464.596,379.659) (464.739,379.766) (464.881,379.873) (465.023,379.979) (465.165,380.086) (465.307,380.193) (465.449,380.299) (465.591,380.406) (465.734,380.513) (465.876,380.620) (466.018,380.726) (466.160,380.833) (466.302,380.940) (466.444,381.047) (466.586,381.154) (466.729,381.260) (466.871,381.367) (467.013,381.474) (467.155,381.581) (467.297,381.688) (467.439,381.795) (467.581,381.902) (467.724,382.009) (467.866,382.115) (468.008,382.222) (468.150,382.329) (468.292,382.436) (468.434,382.543) (468.576,382.650) (468.719,382.757) (468.861,382.864) (469.003,382.971) (469.145,383.078) (469.287,383.185) (469.429,383.292) (469.571,383.399) (469.714,383.507) (469.856,383.614) (469.998,383.721) (470.140,383.828) (470.282,383.935) (470.424,384.042) (470.566,384.149) (470.709,384.256) (470.851,384.363) (470.993,384.470) (471.135,384.578) (471.277,384.685) (471.419,384.792) (471.561,384.899) (471.704,385.006) (471.846,385.113) (471.988,385.221) (472.130,385.328) (472.272,385.435) (472.414,385.542) (472.556,385.649) (472.699,385.757) (472.841,385.864) (472.983,385.971) (473.125,386.078) (473.267,386.185) (473.409,386.293) (473.551,386.400) (473.694,386.507) (473.836,386.614) (473.978,386.722) (474.120,386.829) (474.262,386.936) (474.404,387.043) (474.546,387.151) (474.689,387.258) (474.831,387.365) (474.973,387.472) (475.115,387.580) (475.257,387.687) (475.399,387.794) (475.541,387.901) (475.684,388.009) (475.826,388.116) (475.968,388.223) (476.110,388.330) (476.252,388.438) (476.394,388.545) (476.536,388.652) (476.679,388.759) (476.821,388.867) (476.963,388.974) (477.105,389.081) (477.247,389.188) (477.389,389.296) (477.531,389.403) (477.674,389.510) (477.816,389.617) (477.958,389.725) (478.100,389.832) (478.242,389.939) (478.384,390.046) (478.526,390.154) (478.669,390.261) (478.811,390.368) (478.953,390.475) (479.095,390.583) (479.237,390.690) (479.379,390.797) (479.521,390.904) (479.664,391.011) (479.806,391.118) (479.948,391.226) (480.090,391.333) (480.232,391.440) (480.374,391.547) (480.516,391.654) (480.659,391.761) (480.801,391.869) (480.943,391.976) (481.085,392.083) (481.227,392.190) (481.369,392.297) (481.511,392.404) (481.654,392.511) (481.796,392.618) (481.938,392.725) (482.080,392.833) (482.222,392.940) (482.364,393.047) (482.506,393.154) (482.649,393.261) (482.791,393.368) (482.933,393.475) (483.075,393.582) (483.217,393.689) (483.359,393.796) (483.501,393.903) (483.644,394.010) (483.786,394.117) (483.928,394.224) (484.070,394.330) (484.212,394.437) (484.354,394.544) (484.496,394.651) (484.639,394.758) (484.781,394.865) (484.923,394.972) (485.065,395.079) (485.207,395.185) (485.349,395.292) (485.491,395.399) (485.634,395.506) (485.776,395.613) (485.918,395.719) (486.060,395.826) (486.202,395.933) (486.344,396.040) (486.486,396.146) (486.629,396.253) (486.771,396.360) (486.913,396.466) (487.055,396.573) (487.197,396.680) (487.339,396.786) (487.481,396.893) (487.624,396.999) (487.766,397.106) (487.908,397.212) (488.050,397.319) (488.192,397.425) (488.334,397.532) (488.476,397.638) (488.619,397.745) (488.761,397.851) (488.903,397.958) (489.045,398.064) (489.187,398.170) (489.329,398.277) (489.471,398.383) (489.614,398.489) (489.756,398.596) (489.898,398.702) (490.040,398.808) (490.182,398.914) (490.324,399.021) (490.466,399.127) (490.609,399.233) (490.751,399.339) (490.893,399.445) (491.035,399.551) (491.177,399.658) (491.319,399.764) (491.461,399.870) (491.604,399.976) (491.746,400.082) (491.888,400.188) (492.030,400.294) (492.172,400.400) (492.314,400.505) (492.456,400.611) (492.599,400.717) (492.741,400.823) (492.883,400.929) (493.025,401.035) (493.167,401.140) (493.309,401.246) (493.451,401.352) (493.594,401.457) (493.736,401.563) (493.878,401.669) (494.020,401.774) (494.162,401.880) (494.304,401.986) (494.446,402.091) (494.589,402.197) (494.731,402.302) (494.873,402.407) (495.015,402.513) (495.157,402.618) (495.299,402.724) (495.441,402.829) (495.584,402.934) (495.726,403.040) (495.868,403.145) (496.010,403.250) (496.152,403.355) (496.294,403.460) (496.436,403.566) (496.579,403.671) (496.721,403.776) (496.863,403.881) (497.005,403.986) (497.147,404.091) (497.289,404.196) (497.431,404.301) (497.574,404.406) (497.716,404.511) (497.858,404.615) (498.000,404.720)};
\fill[methodM] (323.481,316.391) circle (2pt);
\node[anchor=south east,font=\scriptsize,text=methodM,fill=white,inner sep=1pt] at (323.481,321.391) {$v_q$};
\fill[methodS] (351.756,316.391) circle (2pt);
\node[anchor=south west,font=\scriptsize,text=methodS,fill=white,inner sep=1pt] at (351.756,321.391) {$\bar v$};
\fill[methodM] (467.187,396.247) circle (2pt);
\node[anchor=south,font=\scriptsize,text=methodM,fill=white,inner sep=1pt] at (467.187,401.247) {$v_f$};
\end{scope}
\node[anchor=west,font=\small\bfseries,text=navy] at (264,226) {(d) NO$_x$: cumulative contribution};
\draw[rulegray,line width=.25pt] (299,29.000) -- (498,29.000);
\node[anchor=east,font=\scriptsize] at (295,29.000) {0};
\draw[rulegray,line width=.25pt] (299,69.227) -- (498,69.227);
\node[anchor=east,font=\scriptsize] at (295,69.227) {500};
\draw[rulegray,line width=.25pt] (299,109.455) -- (498,109.455);
\node[anchor=east,font=\scriptsize] at (295,109.455) {1,000};
\draw[rulegray,line width=.25pt] (299,149.682) -- (498,149.682);
\node[anchor=east,font=\scriptsize] at (295,149.682) {1,500};
\draw[rulegray,line width=.25pt] (299,189.909) -- (498,189.909);
\node[anchor=east,font=\scriptsize] at (295,189.909) {2,000};
\draw[black,line width=.35pt] (299.000,29) -- (299.000,26);
\node[anchor=north,font=\scriptsize] at (299.000,25) {0};
\draw[black,line width=.35pt] (348.750,29) -- (348.750,26);
\node[anchor=north,font=\scriptsize] at (348.750,25) {25};
\draw[black,line width=.35pt] (398.500,29) -- (398.500,26);
\node[anchor=north,font=\scriptsize] at (398.500,25) {50};
\draw[black,line width=.35pt] (448.250,29) -- (448.250,26);
\node[anchor=north,font=\scriptsize] at (448.250,25) {75};
\draw[black,line width=.35pt] (498.000,29) -- (498.000,26);
\node[anchor=north,font=\scriptsize] at (498.000,25) {100};
\draw[black,line width=.45pt] (299,206) -- (299,29) -- (498,29);
\node[anchor=north,font=\scriptsize] at (398.5,10) {Cumulative counted-vehicle share (\%)};
\node[rotate=90,font=\scriptsize] at (266,117.5) {Cumulative emission (g/lane)};
\begin{scope}
\clip (299,29) rectangle (498,206);
\draw[methodO,line width=1.05pt,] plot coordinates {(299.000,29.000) (304.216,33.250) (309.212,36.487) (314.041,39.670) (318.607,42.679) (323.049,44.433) (327.392,46.201) (331.641,48.011) (335.916,49.853) (340.045,51.667) (344.172,53.531) (348.099,55.441) (351.904,57.538) (355.631,59.055) (359.375,60.693) (363.109,62.377) (366.659,64.046) (369.974,65.252) (373.063,66.582) (376.019,68.010) (378.943,69.619) (381.942,71.333) (384.913,72.994) (387.698,74.582) (390.366,76.127) (393.071,77.855) (395.980,79.829) (398.871,81.949) (401.622,84.024) (404.131,85.859) (406.571,87.620) (408.993,89.554) (411.539,91.705) (414.077,93.857) (416.665,95.879) (419.250,97.783) (421.934,99.538) (424.817,101.241) (427.815,102.866) (430.924,104.397) (434.062,105.751) (437.312,106.977) (440.680,108.167) (444.165,109.796) (447.696,111.291) (451.306,112.739) (454.944,114.694) (458.659,116.518) (462.401,118.223) (466.207,119.908) (470.191,121.537) (474.249,124.249) (478.341,127.645) (482.365,130.840) (486.456,133.831) (490.466,136.641) (494.359,139.256) (498.000,141.702)};
\draw[methodR,line width=1.05pt,] plot coordinates {(299.000,29.000) (299.005,29.005) (299.017,29.016) (299.035,29.033) (299.060,29.057) (299.091,29.087) (299.130,29.123) (299.174,29.166) (299.226,29.215) (299.284,29.270) (299.349,29.332) (299.420,29.400) (299.498,29.474) (299.582,29.555) (299.673,29.641) (299.771,29.734) (299.876,29.834) (299.986,29.939) (300.104,30.051) (300.228,30.170) (300.359,30.294) (300.496,30.425) (300.640,30.561) (300.790,30.705) (300.947,30.854) (301.110,31.009) (301.280,31.171) (301.456,31.339) (301.639,31.513) (301.828,31.693) (302.024,31.880) (302.226,32.072) (302.434,32.209) (302.649,32.350) (302.871,32.495) (303.099,32.645) (303.333,32.798) (303.573,32.956) (303.820,33.118) (304.074,33.285) (304.333,33.455) (304.599,33.630) (304.871,33.808) (305.150,33.991) (305.435,34.179) (305.726,34.370) (306.023,34.566) (306.326,34.765) (306.636,34.969) (306.952,35.178) (307.274,35.390) (307.602,35.607) (307.937,35.828) (308.277,36.053) (308.624,36.283) (308.976,36.517) (309.335,36.755) (309.700,36.998) (310.071,37.245) (310.448,37.496) (310.830,37.753) (311.219,38.013) (311.614,38.279) (312.015,38.549) (312.422,38.824) (312.834,39.103) (313.253,39.388) (313.677,39.678) (314.108,39.972) (314.544,40.272) (314.986,40.577) (315.434,40.888) (315.887,41.204) (316.347,41.526) (316.812,41.854) (317.283,42.187) (317.759,42.527) (318.241,42.873) (318.729,43.226) (319.223,43.585) (319.722,43.951) (320.227,44.324) (320.738,44.705) (321.254,45.092) (321.776,45.488) (322.303,45.891) (322.835,46.303) (323.374,46.724) (323.917,47.153) (324.466,47.591) (325.021,48.038) (325.581,48.496) (326.146,48.963) (326.716,49.441) (327.292,49.812) (327.873,50.191) (328.460,50.579) (329.051,50.976) (329.648,51.213) (330.250,51.455) (330.857,51.704) (331.470,51.958) (332.087,52.218) (332.709,52.485) (333.337,52.758) (333.969,53.038) (334.606,53.325) (335.248,53.619) (335.895,53.921) (336.547,54.230) (337.203,54.546) (337.864,54.871) (338.530,55.204) (339.200,55.545) (339.875,55.895) (340.555,56.254) (341.239,56.622) (341.927,56.999) (342.620,57.253) (343.317,57.513) (344.019,57.780) (344.724,58.053) (345.434,58.333) (346.148,58.620) (346.866,58.914) (347.588,59.215) (348.313,59.524) (349.043,59.840) (349.776,60.163) (350.514,60.494) (351.254,60.833) (351.999,61.180) (352.747,61.534) (353.498,61.897) (354.253,62.146) (355.011,62.400) (355.772,62.660) (356.537,62.926) (357.304,63.197) (358.075,63.474) (358.848,63.756) (359.624,64.045) (360.404,64.339) (361.185,64.638) (361.970,64.944) (362.757,65.255) (363.546,65.572) (364.338,65.895) (365.132,66.224) (365.928,66.558) (366.726,66.898) (367.527,67.244) (368.329,67.595) (369.133,67.952) (369.939,68.314) (370.747,68.682) (371.556,69.055) (372.367,69.433) (373.179,69.816) (373.993,70.205) (374.808,70.599) (375.624,70.997) (376.441,71.401) (377.259,71.808) (378.078,72.221) (378.898,72.638) (379.719,73.059) (380.540,73.485) (381.362,73.914) (382.184,74.348) (383.007,74.785) (383.831,75.225) (384.654,75.669) (385.478,76.116) (386.301,76.566) (387.125,77.020) (387.949,77.475) (388.773,77.934) (389.596,78.394) (390.419,78.857) (391.242,79.321) (392.064,79.788) (392.886,80.256) (393.708,80.725) (394.528,81.195) (395.348,81.667) (396.168,82.139) (396.986,82.611) (397.804,83.084) (398.621,83.557) (399.436,84.030) (400.251,84.503) (401.065,84.975) (401.877,85.446) (402.688,85.917) (403.498,86.387) (404.307,86.855) (405.114,87.322) (405.920,87.787) (406.724,88.251) (407.527,88.712) (408.328,89.171) (409.128,89.628) (409.926,90.082) (410.722,90.534) (411.517,90.983) (412.310,91.428) (413.101,91.871) (413.890,92.310) (414.678,92.746) (415.463,93.178) (416.247,93.606) (417.028,94.030) (417.808,94.451) (418.586,94.867) (419.362,95.279) (420.135,95.687) (420.907,96.090) (421.676,96.488) (422.444,96.882) (423.209,97.271) (423.972,97.656) (424.733,98.035) (425.491,98.409) (426.248,98.779) (427.002,99.143) (427.753,99.502) (428.503,99.856) (429.250,100.205) (429.995,100.548) (430.738,100.886) (431.478,101.219) (432.216,101.546) (432.951,101.868) (433.684,102.184) (434.414,102.496) (435.143,102.801) (435.868,103.101) (436.591,103.396) (437.312,103.686) (438.030,103.970) (438.746,104.249) (439.459,104.522) (440.169,104.790) (440.877,105.053) (441.582,105.311) (442.284,105.563) (442.984,105.810) (443.681,106.052) (444.376,106.290) (445.068,106.522) (445.757,106.749) (446.443,107.080) (447.126,107.404) (447.807,107.721) (448.484,108.031) (449.159,108.334) (449.831,108.631) (450.500,108.920) (451.166,109.203) (451.829,109.480) (452.489,109.750) (453.146,110.014) (453.800,110.272) (454.451,110.524) (455.098,110.771) (455.742,111.011) (456.383,111.246) (457.021,111.595) (457.656,111.936) (458.287,112.270) (458.915,112.595) (459.539,112.913) (460.160,113.224) (460.778,113.527) (461.392,113.823) (462.002,114.112) (462.609,114.395) (463.212,114.671) (463.811,114.941) (464.406,115.205) (464.998,115.463) (465.586,115.715) (466.170,115.961) (466.750,116.202) (467.326,116.437) (467.899,116.668) (468.467,116.894) (469.031,117.272) (469.590,117.642) (470.146,118.005) (470.697,118.360) (471.244,118.818) (471.787,119.267) (472.325,119.706) (472.859,120.137) (473.388,120.560) (473.913,120.974) (474.433,121.380) (474.949,121.779) (475.460,122.170) (475.966,122.554) (476.467,122.930) (476.964,123.300) (477.456,123.664) (477.942,124.021) (478.424,124.371) (478.901,124.716) (479.373,125.054) (479.840,125.387) (480.301,125.714) (480.757,126.035) (481.209,126.352) (481.655,126.662) (482.095,126.968) (482.530,127.269) (482.960,127.564) (483.385,127.855) (483.804,128.141) (484.217,128.422) (484.625,128.699) (485.028,128.971) (485.425,129.238) (485.816,129.501) (486.201,129.759) (486.581,130.013) (486.955,130.263) (487.323,130.508) (487.685,130.750) (488.042,130.986) (488.393,131.219) (488.737,131.447) (489.076,131.672) (489.409,131.892) (489.736,132.107) (490.057,132.319) (490.372,132.527) (490.680,132.730) (490.983,132.929) (491.279,133.124) (491.570,133.315) (491.854,133.502) (492.132,133.685) (492.404,133.863) (492.669,134.037) (492.928,134.208) (493.181,134.374) (493.428,134.536) (493.668,134.693) (493.902,134.847) (494.130,134.996) (494.351,135.141) (494.566,135.282) (494.775,135.419) (494.977,135.611) (495.172,135.798) (495.361,135.978) (495.544,136.152) (495.720,136.320) (495.890,136.481) (496.054,136.637) (496.210,136.786) (496.361,136.929) (496.504,137.066) (496.641,137.197) (496.772,137.321) (496.896,137.439) (497.014,137.551) (497.124,137.657) (497.229,137.756) (497.327,137.849) (497.418,137.936) (497.502,138.017) (497.580,138.091) (497.651,138.159) (497.716,138.220) (497.774,138.276) (497.826,138.325) (497.870,138.367) (497.909,138.404) (497.940,138.434) (497.965,138.458) (497.983,138.475) (497.995,138.486) (498.000,138.491)};
\draw[methodS,line width=1.05pt,dashed] plot coordinates {(299.000,29.000) (498.000,98.334)};
\draw[methodM,line width=1.05pt,] plot coordinates {(299.000,29.000) (299.005,29.005) (299.017,29.016) (299.035,29.033) (299.060,29.057) (299.091,29.087) (299.130,29.123) (299.174,29.166) (299.226,29.215) (299.284,29.270) (299.349,29.332) (299.420,29.400) (299.498,29.474) (299.582,29.555) (299.673,29.641) (299.771,29.734) (299.876,29.834) (299.986,29.939) (300.104,30.051) (300.228,30.170) (300.359,30.294) (300.496,30.425) (300.640,30.561) (300.790,30.705) (300.947,30.854) (301.110,31.009) (301.280,31.171) (301.456,31.339) (301.639,31.513) (301.828,31.693) (302.024,31.880) (302.226,32.072) (302.434,32.271) (302.649,32.476) (302.871,32.686) (303.099,32.903) (303.333,33.126) (303.573,33.356) (303.820,33.591) (304.074,33.832) (304.333,34.079) (304.599,34.332) (304.871,34.591) (305.150,34.857) (305.435,35.128) (305.726,35.405) (306.023,35.688) (306.326,35.977) (306.636,36.271) (306.952,36.572) (307.274,36.879) (307.602,37.191) (307.937,37.509) (308.277,37.833) (308.624,38.163) (308.976,38.498) (309.335,38.839) (309.700,39.186) (310.071,39.539) (310.448,39.897) (310.830,40.261) (311.219,40.631) (311.614,41.006) (312.015,41.386) (312.422,41.773) (312.834,42.164) (313.253,42.561) (313.677,42.964) (314.108,43.372) (314.544,43.785) (314.986,44.204) (315.434,44.628) (315.887,45.057) (316.347,45.492) (316.812,45.932) (317.283,46.376) (317.759,46.826) (318.241,47.281) (318.729,47.741) (319.223,48.207) (319.722,48.676) (320.227,49.151) (320.738,49.631) (321.254,50.115) (321.776,50.604) (322.303,51.098) (322.835,51.596) (323.374,52.099) (323.917,52.607) (324.466,53.118) (325.021,53.634) (325.581,54.154) (326.146,54.679) (326.716,55.207) (327.292,55.740) (327.873,56.276) (328.460,56.817) (329.051,57.361) (329.648,57.909) (330.250,58.460) (330.857,59.015) (331.470,59.573) (332.087,60.135) (332.709,60.700) (333.337,61.267) (333.969,61.838) (334.606,62.412) (335.248,62.989) (335.895,63.568) (336.547,64.150) (337.203,64.734) (337.864,65.321) (338.530,65.910) (339.200,66.501) (339.875,67.094) (340.555,67.689) (341.239,68.285) (341.927,68.883) (342.620,69.483) (343.317,70.084) (344.019,70.686) (344.724,71.290) (345.434,71.894) (346.148,72.500) (346.866,73.106) (347.588,73.713) (348.313,74.320) (349.043,74.927) (349.776,75.535) (350.514,76.143) (351.254,76.751) (351.999,77.359) (352.747,77.967) (353.498,78.574) (354.253,79.181) (355.011,79.787) (355.772,80.392) (356.537,80.997) (357.304,81.601) (358.075,82.204) (358.848,82.805) (359.624,83.406) (360.404,84.005) (361.185,84.602) (361.970,85.198) (362.757,85.793) (363.546,86.386) (364.338,86.977) (365.132,87.566) (365.928,88.153) (366.726,88.739) (367.527,89.322) (368.329,89.904) (369.133,90.483) (369.939,91.060) (370.747,91.635) (371.556,92.207) (372.367,92.777) (373.179,93.345) (373.993,93.911) (374.808,94.474) (375.624,95.035) (376.441,95.594) (377.259,96.150) (378.078,96.703) (378.898,97.255) (379.719,97.804) (380.540,98.350) (381.362,98.895) (382.184,99.437) (383.007,99.976) (383.831,100.514) (384.654,101.049) (385.478,101.582) (386.301,102.113) (387.125,102.642) (387.949,103.169) (388.773,103.693) (389.596,104.216) (390.419,104.738) (391.242,105.257) (392.064,105.775) (392.886,106.291) (393.708,106.805) (394.528,107.318) (395.348,107.830) (396.168,108.340) (396.986,108.849) (397.804,109.357) (398.621,109.863) (399.436,110.369) (400.251,110.874) (401.065,111.378) (401.877,111.882) (402.688,112.384) (403.498,112.886) (404.307,113.388) (405.114,113.890) (405.920,114.391) (406.724,114.892) (407.527,115.392) (408.328,115.893) (409.128,116.394) (409.926,116.895) (410.722,117.396) (411.517,117.898) (412.310,118.400) (413.101,118.902) (413.890,119.405) (414.678,119.909) (415.463,120.413) (416.247,120.918) (417.028,121.424) (417.808,121.931) (418.586,122.439) (419.362,122.948) (420.135,123.457) (420.907,123.968) (421.676,124.480) (422.444,124.994) (423.209,125.508) (423.972,126.024) (424.733,126.541) (425.491,127.060) (426.248,127.579) (427.002,128.101) (427.753,128.623) (428.503,129.147) (429.250,129.673) (429.995,130.200) (430.738,130.728) (431.478,131.258) (432.216,131.790) (432.951,132.323) (433.684,132.857) (434.414,133.393) (435.143,133.930) (435.868,134.468) (436.591,135.008) (437.312,135.549) (438.030,136.092) (438.746,136.636) (439.459,137.181) (440.169,137.727) (440.877,138.274) (441.582,138.823) (442.284,139.372) (442.984,139.923) (443.681,140.474) (444.376,141.026) (445.068,141.580) (445.757,142.133) (446.443,142.688) (447.126,143.243) (447.807,143.799) (448.484,144.355) (449.159,144.912) (449.831,145.469) (450.500,146.026) (451.166,146.583) (451.829,147.140) (452.489,147.697) (453.146,148.254) (453.800,148.811) (454.451,149.368) (455.098,149.924) (455.742,150.479) (456.383,151.034) (457.021,151.589) (457.656,152.142) (458.287,152.695) (458.915,153.246) (459.539,153.796) (460.160,154.346) (460.778,154.894) (461.392,155.440) (462.002,155.985) (462.609,156.528) (463.212,157.070) (463.811,157.610) (464.406,158.147) (464.998,158.683) (465.586,159.217) (466.170,159.748) (466.750,160.277) (467.326,160.804) (467.899,161.328) (468.467,161.849) (469.031,162.368) (469.590,162.883) (470.146,163.396) (470.697,163.906) (471.244,164.413) (471.787,164.916) (472.325,165.416) (472.859,165.913) (473.388,166.406) (473.913,166.896) (474.433,167.382) (474.949,167.864) (475.460,168.343) (475.966,168.818) (476.467,169.288) (476.964,169.755) (477.456,170.217) (477.942,170.675) (478.424,171.129) (478.901,171.579) (479.373,172.024) (479.840,172.464) (480.301,172.900) (480.757,173.332) (481.209,173.759) (481.655,174.181) (482.095,174.598) (482.530,175.010) (482.960,175.418) (483.385,175.820) (483.804,176.218) (484.217,176.610) (484.625,176.997) (485.028,177.379) (485.425,177.756) (485.816,178.128) (486.201,178.494) (486.581,178.855) (486.955,179.210) (487.323,179.560) (487.685,179.905) (488.042,180.244) (488.393,180.578) (488.737,180.906) (489.076,181.228) (489.409,181.545) (489.736,181.856) (490.057,182.161) (490.372,182.461) (490.680,182.755) (490.983,183.043) (491.279,183.325) (491.570,183.602) (491.854,183.872) (492.132,184.137) (492.404,184.395) (492.669,184.648) (492.928,184.895) (493.181,185.136) (493.428,185.371) (493.668,185.600) (493.902,185.823) (494.130,186.039) (494.351,186.250) (494.566,186.455) (494.775,186.653) (494.977,186.846) (495.172,187.032) (495.361,187.212) (495.544,187.386) (495.720,187.554) (495.890,187.716) (496.054,187.871) (496.210,188.021) (496.361,188.164) (496.504,188.301) (496.641,188.431) (496.772,188.556) (496.896,188.674) (497.014,188.786) (497.124,188.891) (497.229,188.991) (497.327,189.084) (497.418,189.170) (497.502,189.251) (497.580,189.325) (497.651,189.393) (497.716,189.455) (497.774,189.510) (497.826,189.559) (497.870,189.602) (497.909,189.638) (497.940,189.668) (497.965,189.692) (497.983,189.709) (497.995,189.720) (498.000,189.725)};
\draw[methodA,line width=1.05pt,dashed] plot coordinates {(299.000,29.000) (299.005,29.004) (299.017,29.013) (299.035,29.028) (299.060,29.048) (299.091,29.073) (299.130,29.104) (299.174,29.140) (299.226,29.181) (299.284,29.227) (299.349,29.279) (299.420,29.336) (299.498,29.399) (299.582,29.467) (299.673,29.540) (299.771,29.618) (299.876,29.701) (299.986,29.790) (300.104,29.884) (300.228,29.984) (300.359,30.089) (300.496,30.198) (300.640,30.314) (300.790,30.434) (300.947,30.560) (301.110,30.690) (301.280,30.826) (301.456,30.968) (301.639,31.114) (301.828,31.266) (302.024,31.423) (302.226,31.585) (302.434,31.752) (302.649,31.924) (302.871,32.101) (303.099,32.284) (303.333,32.471) (303.573,32.664) (303.820,32.862) (304.074,33.065) (304.333,33.273) (304.599,33.486) (304.871,33.704) (305.150,33.927) (305.435,34.155) (305.726,34.388) (306.023,34.626) (306.326,34.869) (306.636,35.117) (306.952,35.370) (307.274,35.628) (307.602,35.891) (307.937,36.159) (308.277,36.431) (308.624,36.709) (308.976,36.991) (309.335,37.278) (309.700,37.570) (310.071,37.867) (310.448,38.168) (310.830,38.474) (311.219,38.785) (311.614,39.101) (312.015,39.421) (312.422,39.746) (312.834,40.076) (313.253,40.410) (313.677,40.749) (314.108,41.092) (314.544,41.440) (314.986,41.793) (315.434,42.150) (315.887,42.511) (316.347,42.877) (316.812,43.247) (317.283,43.622) (317.759,44.000) (318.241,44.384) (318.729,44.771) (319.223,45.163) (319.722,45.559) (320.227,45.959) (320.738,46.363) (321.254,46.771) (321.776,47.183) (322.303,47.599) (322.835,48.019) (323.374,48.443) (323.917,48.871) (324.466,49.302) (325.021,49.737) (325.581,50.176) (326.146,50.618) (326.716,51.064) (327.292,51.514) (327.873,51.967) (328.460,52.423) (329.051,52.882) (329.648,53.345) (330.250,53.811) (330.857,54.279) (331.470,54.751) (332.087,55.226) (332.709,55.704) (333.337,56.184) (333.969,56.667) (334.606,57.153) (335.248,57.641) (335.895,58.131) (336.547,58.624) (337.203,59.119) (337.864,59.617) (338.530,60.116) (339.200,60.618) (339.875,61.121) (340.555,61.626) (341.239,62.133) (341.927,62.641) (342.620,63.151) (343.317,63.662) (344.019,64.174) (344.724,64.688) (345.434,65.203) (346.148,65.719) (346.866,66.236) (347.588,66.753) (348.313,67.271) (349.043,67.790) (349.776,68.310) (350.514,68.829) (351.254,69.349) (351.999,69.870) (352.747,70.390) (353.498,70.910) (354.253,71.431) (355.011,71.951) (355.772,72.471) (356.537,72.990) (357.304,73.509) (358.075,74.028) (358.848,74.546) (359.624,75.063) (360.404,75.580) (361.185,76.096) (361.970,76.610) (362.757,77.124) (363.546,77.637) (364.338,78.149) (365.132,78.659) (365.928,79.168) (366.726,79.676) (367.527,80.183) (368.329,80.688) (369.133,81.191) (369.939,81.694) (370.747,82.194) (371.556,82.693) (372.367,83.191) (373.179,83.687) (373.993,84.181) (374.808,84.674) (375.624,85.165) (376.441,85.654) (377.259,86.141) (378.078,86.627) (378.898,87.111) (379.719,87.593) (380.540,88.074) (381.362,88.553) (382.184,89.030) (383.007,89.505) (383.831,89.979) (384.654,90.451) (385.478,90.922) (386.301,91.391) (387.125,91.858) (387.949,92.324) (388.773,92.789) (389.596,93.252) (390.419,93.713) (391.242,94.173) (392.064,94.632) (392.886,95.090) (393.708,95.546) (394.528,96.001) (395.348,96.456) (396.168,96.909) (396.986,97.361) (397.804,97.812) (398.621,98.262) (399.436,98.711) (400.251,99.160) (401.065,99.608) (401.877,100.055) (402.688,100.502) (403.498,100.948) (404.307,101.394) (405.114,101.839) (405.920,102.284) (406.724,102.729) (407.527,103.173) (408.328,103.618) (409.128,104.062) (409.926,104.506) (410.722,104.951) (411.517,105.395) (412.310,105.840) (413.101,106.284) (413.890,106.729) (414.678,107.175) (415.463,107.620) (416.247,108.067) (417.028,108.513) (417.808,108.960) (418.586,109.408) (419.362,109.856) (420.135,110.305) (420.907,110.755) (421.676,111.205) (422.444,111.656) (423.209,112.108) (423.972,112.560) (424.733,113.013) (425.491,113.468) (426.248,113.923) (427.002,114.378) (427.753,114.835) (428.503,115.293) (429.250,115.751) (429.995,116.211) (430.738,116.671) (431.478,117.132) (432.216,117.594) (432.951,118.057) (433.684,118.521) (434.414,118.986) (435.143,119.451) (435.868,119.918) (436.591,120.385) (437.312,120.853) (438.030,121.322) (438.746,121.792) (439.459,122.262) (440.169,122.733) (440.877,123.205) (441.582,123.677) (442.284,124.150) (442.984,124.623) (443.681,125.097) (444.376,125.571) (445.068,126.046) (445.757,126.521) (446.443,126.996) (447.126,127.472) (447.807,127.947) (448.484,128.423) (449.159,128.899) (449.831,129.375) (450.500,129.850) (451.166,130.326) (451.829,130.801) (452.489,131.276) (453.146,131.751) (453.800,132.225) (454.451,132.699) (455.098,133.172) (455.742,133.644) (456.383,134.116) (457.021,134.587) (457.656,135.057) (458.287,135.526) (458.915,135.994) (459.539,136.461) (460.160,136.927) (460.778,137.392) (461.392,137.855) (462.002,138.316) (462.609,138.776) (463.212,139.235) (463.811,139.692) (464.406,140.147) (464.998,140.600) (465.586,141.051) (466.170,141.500) (466.750,141.947) (467.326,142.392) (467.899,142.835) (468.467,143.275) (469.031,143.713) (469.590,144.149) (470.146,144.582) (470.697,145.012) (471.244,145.439) (471.787,145.864) (472.325,146.286) (472.859,146.705) (473.388,147.121) (473.913,147.534) (474.433,147.943) (474.949,148.350) (475.460,148.753) (475.966,149.153) (476.467,149.550) (476.964,149.943) (477.456,150.332) (477.942,150.718) (478.424,151.100) (478.901,151.479) (479.373,151.854) (479.840,152.225) (480.301,152.592) (480.757,152.955) (481.209,153.315) (481.655,153.670) (482.095,154.021) (482.530,154.368) (482.960,154.711) (483.385,155.050) (483.804,155.384) (484.217,155.715) (484.625,156.040) (485.028,156.362) (485.425,156.679) (485.816,156.992) (486.201,157.300) (486.581,157.604) (486.955,157.903) (487.323,158.197) (487.685,158.487) (488.042,158.773) (488.393,159.053) (488.737,159.329) (489.076,159.601) (489.409,159.867) (489.736,160.129) (490.057,160.386) (490.372,160.638) (490.680,160.885) (490.983,161.127) (491.279,161.365) (491.570,161.597) (491.854,161.825) (492.132,162.048) (492.404,162.265) (492.669,162.478) (492.928,162.686) (493.181,162.888) (493.428,163.086) (493.668,163.278) (493.902,163.466) (494.130,163.648) (494.351,163.826) (494.566,163.998) (494.775,164.165) (494.977,164.327) (495.172,164.483) (495.361,164.635) (495.544,164.781) (495.720,164.923) (495.890,165.059) (496.054,165.189) (496.210,165.315) (496.361,165.435) (496.504,165.550) (496.641,165.660) (496.772,165.765) (496.896,165.864) (497.014,165.959) (497.124,166.047) (497.229,166.131) (497.327,166.209) (497.418,166.282) (497.502,166.350) (497.580,166.413) (497.651,166.470) (497.716,166.521) (497.774,166.568) (497.826,166.609) (497.870,166.645) (497.909,166.676) (497.940,166.701) (497.965,166.721) (497.983,166.736) (497.995,166.745) (498.000,166.749)};
\end{scope}
\end{tikzpicture}
}\end{center}
\textbf{Rate curves (a,b).} Black is the zero-acceleration lookup; orange dashed is its cubic approximation.
Markers identify the queue speed (13.611 mph), observed mean speed (23.557 mph) and free-flow speed (64.161 mph).
The queue and observed mean speeds select the same lookup rate for both pollutants. The free-flow rate is much higher.

\textbf{Cumulative curves (c,d).} The horizontal coordinate is the cumulative share of counted vehicles,
not clock time. O accumulates each bin's complete-link emission in chronological order;
R accumulates reconstructed group emissions in entry order using its vehicle weights.
M and A use the same group delays $w(u)$ to allocate each group's free/queue time; this illustrative
extension preserves their analytical episode totals exactly. S assigns the same emission to every vehicle,
so its cumulative curve is straight.
For O, straight segments within a five-minute group distribute that group's total uniformly over its vehicle share.
Every endpoint equals its episode total on page 3.

\notebox{\textbf{What this example shows.} R is closest to O for CO$_2$ and NO$_x$ in this case.
S does not win every episode. M's higher free-flow rate can offset its lower total vehicle-hours relative to O.
These plots explain the calculations; none supplies independent vehicle trajectories or measured emissions.}
